\documentclass[11pt,a4paper,oneside]{book}

\usepackage[margin=1in]{geometry}

\usepackage{fontspec}

\setmainfont{DejaVu Serif}

\setsansfont{DejaVu Sans}

\setmonofont{DejaVu Sans Mono}

\usepackage{amsmath,amssymb,mathtools}

\usepackage{amsthm}

\usepackage{graphicx}

\usepackage{booktabs,longtable,array}

\usepackage{xcolor}

\usepackage{enumitem}

\usepackage[hidelinks]{hyperref}

\setlength{\parskip}{0.5em}

\setlength{\parindent}{0pt}

\newcolumntype{L}[1]{>{\raggedright\arraybackslash}p{#1}}

% ---- unified macro set ----

\newcommand{\Op}{\Omega_P}
\newcommand{\lf}{\lambda_{\rm form}}
\newcommand{\Grec}{\mathcal G_{\rm recepta}}
\newcommand{\Igate}{I_{\rm gate}}
\newcommand{\note}[1]{\textcolor{black!65}{\emph{#1}}}
\newcommand{\msg}[1]{\begin{center}\fbox{\parbox{0.9\textwidth}{\centering #1}}\end{center}}
\newcommand{\VFS}{\ensuremath{\mathrm{V.F.S.}}}
\newcommand{\Lc}{\Lambda_c}
\newcommand{\Avfs}{A_{\rm VFS}}
\newcommand{\Bvfs}{B_{\rm VFS}}
\newcommand{\Qf}{Q_{\rm fold}}
\newcommand{\Tfold}{\mathcal{T}_{\rm fold}}
\newcommand{\dd}{\,d}
\newcommand{\LV}{\mathcal L_{\rm VFS}}
\newcommand{\DA}{\mathcal D_A}
\newcommand{\DAf}{\mathcal D_A^{\rm fold}}
\newcommand{\Om}{\Omega_P}
\newcommand{\reset}{\mathfrak{R}}
\newcommand{\qf}{q_{\rm fold}}
\newcommand{\Qs}{Q_{\rm fold}^{*}}
\newcommand{\Af}{A_{\rm fold}}
\newcommand{\Bf}{B_{\rm fold}}
\newcommand{\Vf}{\mathcal V_{\rm fold}}
\newcommand{\Vt}{\widetilde{\mathcal V}_{\rm fold}}
\newcommand{\Lext}{\mathcal L_{\rm ext}}
\newcommand{\Lt}{\widetilde{\mathcal L}_{\rm ext}}
\newcommand{\rs}{r_\sigma}
\newcommand{\rR}{r^{R}_{\rm fold}}
\newcommand{\rRj}{r^{R}_{{\rm fold},j}}
\newcommand{\Afold}{A_{\rm fold}}
\newcommand{\Bfold}{B_{\rm fold}}
\newcommand{\Beff}{B_{\rm eff}}
\newcommand{\Qfj}{Q_{{\rm fold},j}}
\newcommand{\fold}{\mathfrak{F}}
\newcommand{\Hh}{\mathbb{H}^2}
\newcommand{\laph}{\Delta_h}
\newcommand{\dlam}{\delta\lambda}
\newcommand{\Glam}{\Gamma_\lambda}
\newcommand{\Dq}{D_q}
\newcommand{\kmax}{k_{\max}}
\newcommand{\Lam}{\Lambda}
\newcommand{\Rc}{R_c}
\newcommand{\lcom}{\ell_{\rm com}}
\newcommand{\Lphys}{L_{\rm phys}}
\newcommand{\Ibud}{\mathcal I}
\newcommand{\Acom}{\mathcal A}
\newcommand{\lc}{\ell_{\rm com}}
\newcommand{\Lp}{L_{\rm phys}}
\newcommand{\Ib}{\mathcal I}
\newcommand{\linit}{\ell_{\rm init}}
\newcommand{\kb}{\kappa_b}
\newcommand{\lb}{\ell_b}
\newcommand{\Lop}{\mathcal{L}_0}
\newcommand{\etaf}{\eta_{\rm fold}}
\newcommand{\ldz}{\dot\lambda^{0}}
\newcommand{\kw}{\hat\kappa}
\newcommand{\dwall}{\delta_{\rm wall}}
\newcommand{\ldf}{\dot\lambda_{\rm fold}}
\newcommand{\NEC}{\mathcal N}
\newcommand{\Omd}{\Omega_P^{\rm death}}
\newcommand{\Wp}{\overline W_{+}}
\newcommand{\Wm}{\overline W}
\newcommand{\OmegaP}{\Omega_P}
\newcommand{\LVFS}{\mathcal L_{\rm VFS}}
\newcommand{\Ric}{\operatorname{Ric}}
\newcommand{\Qfold}{Q_{\rm fold}}
\newcommand{\sR}{s_R}
\newcommand{\mR}{m_R}
\newcommand{\Kj}{[K]}
\newcommand{\weff}{w_{\rm eff}}

\theoremstyle{plain}

\newtheorem{proposition}{Proposition}[chapter]

\newtheorem{corollary}{Corollary}[chapter]

\newtheorem{lemma}{Lemma}[chapter]

\newtheorem{theorem}{Theorem}[chapter]

\theoremstyle{definition}

\newtheorem{definition}{Definition}[chapter]

\newtheorem{assumption}{Standing Hypothesis}[chapter]

\newtheorem{postulate}{Postulate}[chapter]

\theoremstyle{remark}

\newtheorem*{remark}{Remark}

\title{\textbf{\Huge V.F.S. v2.0 (Open-Gate)}\\[1.2em]
{\LARGE Folding, Stability, and Structure Formation}\\[0.6em]
{\large A Companion Volume}}
\author{Derived companion to the foundational layer\\
\small (core $\cdot$ Lyapunov $\cdot$ geometry $\cdot$ cosmology $\cdot$ theological $\cdot$ Ricci)}
\date{}


\begin{document}

\frontmatter

\maketitle

\tableofcontents

\cleardoublepage

\chapter*{Notation and Conventions}
\addcontentsline{toc}{chapter}{Notation and Conventions}
\markboth{Notation and Conventions}{Notation and Conventions}

This volume collects the folding, stability, and structure-formation results derived as a companion to the
foundational V.F.S.\ layer. Throughout, \emph{Propositions} are closed-form results; \emph{Corollaries}
are interpretive readings in the symbolic register, not independent metaphysical claims; \emph{Remarks}
and italicised notes are commentary. All stability results are \textbf{domain-conditional}: they hold on
the positively-invariant active domain and assert set stability with partial convergence (no global or
automatic stability, and no equilibrium --- epektasis drives $\Op\to\infty$).

\medskip
\textbf{Core symbols.}
\begin{longtable}{L{0.22\textwidth} L{0.72\textwidth}}
\toprule
$\Op$ & Brim / scale factor; $\dot\Op=\alpha\lambda$, Hubble $H=\alpha\lambda/\Op$\\
$\lambda,\ \sigma$ & Sophia (transformed resistance + grace); resistance\\
$u=\sqrt{VF+\varepsilon}$ & synergy of Voluntas $V$ and Factum $F$; threshold $\Lambda_c=\gamma/\delta$\\
$I_{\rm gate},\ \mathcal{G}_{\rm recepta}$ & Open-Gate grace inflow; received-grace budget\\
$q_{\rm fold}$ & folding (shape) amplitude; order parameter of the $\mathbb Z_2$ morphology\\
$A_{\rm fold},\ B_{\rm fold}$ & folding activation index and quartic coefficient ($B_{\rm fold}>0$)\\
$Q_{\rm fold}^{*}$ & signed normal-form branch parameter $A_{\rm fold}/B_{\rm fold}$ (may be negative off the active branch); folded equilibria $q_{\rm fold,\pm}^{*}=\pm\sqrt{Q_{\rm fold}^{*}}$ exist only for $A_{\rm fold}>0$\\
$Q_{\rm fold}$ & physical folded intensity $q_{\rm fold}^{2}\ge0$ (distinct from the signed $Q_{\rm fold}^{*}$); $\widehat q_n$ its spectral coefficients\\
$\mathcal{L}_{\rm VFS},\ \mathcal{L}_{\rm ext}$ & Lyapunov certificate; folding-extended certificate\\
$h,\ \mathbb{H}^2$ & comoving hyperbolic spatial metric, $K_h=-1$, curvature radius $R_c$\\
$\mathcal I(t)=\int_0^t dt'/\Op^2$ & expansion / activity budget; finite $\mathcal I_\infty<\infty$\\
$D_q,\ \ell_b$ & folding diffusivity (bending stiffness); bending length $\ell_b=\sqrt{D_q/a_0}$\\
$\Theta,\ Q_{\rm crit}$ & folding-response coefficient $c_0-c_1(\gamma+k_\sigma)$; reset critical level\\
$\lambda_{\rm form}$ & Sophia embodied in form; balance $\sigma+\lambda+\lambda_{\rm form}=C_0+\mathcal{G}_{\rm recepta}$\\
$\eta_{\rm fold}$ & folding feedback coefficient (distinct from any core-system $\eta$)\\
$\Delta_h$ & Laplace--Beltrami on $\mathbb{H}^2$; live-domain $\Delta_g=\Op^{-2}\Delta_h$\\
$r_\sigma,\ g(u)$ & resistance multiplier $\in[1,\pi]$; Dolorosum gate $g(u)=(u/\Lambda_c)e^{1-u/\Lambda_c}$\\
\bottomrule
\end{longtable}

\textbf{Two activation-index conventions.} The activation index appears in two forms across the layers,
reconciled in the Folding--Lyapunov Addendum:
\[
\underbrace{A_{\rm fold}=c_0\lambda+c_1\dot\lambda^{(0)}-a_0,\ \ \dot\lambda^{(0)}=-\dot\sigma+I_{\rm gate}}_{\text{canonical (core rate)}}
\qquad
\underbrace{A_{\rm fold}=(c_0-c_1\gamma)\lambda+c_1k_\sigma\sigma+c_1I_{\rm gate}-a_0}_{\text{geometric (minimal rate)}}.
\]
The canonical (core-consistent) form is primary; the geometric parameters $(k_\sigma,\gamma)$ are an
effective linear surrogate. The closed forms $\Theta,\ Q_{\rm crit}$ are stated in the geometric
convention; the invariance and reset-bound results are convention-independent.

The folding diffusivity is structural, $D_q=a_0\ell_b^2$ (rigidity $\times$ bending length$^2$), and the
independent Sophia diffusivity is zero, $D=0$.


\chapter*{The Foundational Layer (Summary)}
\addcontentsline{toc}{chapter}{The Foundational Layer (Summary)}
\markboth{The Foundational Layer}{The Foundational Layer}

This volume sits atop the foundational V.F.S.\ layer, summarised here so it reads self-contained, and
records what is assumed, what is derived, and what is only interpreted.

\textbf{Core.} V.F.S.\ (Volo.\ Facio.\ Sum.) is a hybrid dynamical system on the live domain in the
coordinates Voluntas $V$, Factum $F$, Sophia $\lambda$, with synergy $u=\sqrt{VF+\varepsilon}$, resistance
$\sigma$, Pleroma $P$ and Brim $\Op$, and threshold $\Lambda_c=\gamma/\delta$. The Open-Gate law
$\dot\lambda=-\dot\sigma+I_{\rm gate}$ admits a bounded grace inflow; the closed Microcosm is the limit
$I_{\rm gate}\to0$. The central coordinate is not merely convenient: given an externally clocked gate and
the balance $\sigma+\lambda+\lambda_{\rm form}=C_0+\mathcal{G}_{\rm recepta}(t)$,
$\dot{\mathcal{G}}_{\rm recepta}=I_{\rm gate}$, the state-map rank increases by one when the gate opens, so
the balance coordinate is necessary iff the Vessel is open. Its mathematical role is theorem; its
evolution law is constitutive; its name ``Sophia'' is the theological interpretation. The Filtrum Lucis
softplus $\Phi$ yields the Paschal triad ($\Phi'''(0)=0$ midway between the symmetric zeros of
$\Phi^{(4)}$). Resurrectio is a hybrid re-entry map.

\textbf{Finite floor and hyperbolicity.} The live surface is hyperbolic, $K_h=-1$, $R=-2/\Op^2$. This is
not an independent primitive: in the $2{+}1$ reconstruction only $k=-1$ admits a zero-density floor at
finite scale, so hyperbolicity is the geometric face of the finite-floor (no-Big-Bang) premise.
Expansion alone does not force $K_h=-1$, but $K_h=-1$ is not a free third assumption. The irreducible
core is therefore two primitives: the open gate, and the finite floor.

\textbf{Lyapunov.} On the active domain $\mathcal D_A$ the certificate obeys
$\dot{\mathcal L}_{\rm VFS}\le C-C_1\mathcal L_{\rm VFS}$ with $C_1=\min\{\kappa_\sigma,\kappa_\Delta\}>0$
from two sign-definite contracting modes (cleansing, will--action alignment). Grace and folding enter the
bounded constant $C$, never the contraction rate $C_1$. The result is set stability with partial
convergence, forward completeness, and non-Zeno resets --- domain-conditional, not global, since
epektasis drives $\Op$.

\textbf{Sophianic Folding.} Folding is core-derived, not an add-on. The core supplies the driver
$A_{\rm fold}=c_0\lambda+c_1\dot\lambda^{(0)}-a_0$ (geometric surrogate
$A_{\rm fold}=(c_0-c_1\gamma)\lambda+c_1k_\sigma\sigma+c_1I_{\rm gate}-a_0$); when $A_{\rm fold}$ crosses
zero the unfolded form loses stability. The pitchfork normal form is not assumed: it follows from the
$\mathbb Z_2$ symmetry of the unfolded Vessel, transversality, and a one-dimensional centre manifold, with
$B_{\rm fold}$ the bare stiffness renormalised by the slaved structure. Curvature pockets are
$R=-2/\Op^2+\rho_1 Q_{\rm fold}$.

\textbf{Geometry, structure, entropy.} The bending sector is conformally blind in two dimensions,
$|\nabla q_{\rm fold}|_g^2\,dV_g=|\nabla q_{\rm fold}|_h^2\,dV_h$ for $g=\Op^2h$, so expansion does not
stretch the folding fabric ($p_{\rm expansion}\equiv0$); the true warp is kathartic. Promoting
$q_{\rm fold}$ to a field gives a spectral ladder $\Lambda=A_{\rm fold}\Op^2/\widehat\kappa$, mode $n$
igniting iff $\Lambda>\mu_n$; the one-mode chart stays faithful for amplitude-keyed results. The
homogeneous Vessel is the constrained entropy maximum, with $\delta^2 W_+^{(n)}=-2\sigma\mu_n a_n^2$: the
same eigenvalue that orders ignition orders the entropy price of form.

\textbf{Cosmology and energy conditions.} The Lorentzian lift is a $2{+}1$ open FLRW spacetime,
$ds^2=-dt^2+\Op^2 h$. In $2{+}1$ the comoving timelike-focusing (strong-energy) condition reduces to
$p_{\rm eff}\ge0$; the NEC is tracked separately through $\rho_{\rm eff}+p_{\rm eff}$. A Resurrectio reset
may produce a bounded surgical phantom budget, but $\int(\rho_{\rm eff}+p_{\rm eff})\,dt$ is finite; the
balanced reset is the unique point where that budget, the extrinsic-curvature jump $[K]$, and the
deviation from a purely conformal junction vanish together.

\textbf{Theological \& Ricci.} The theological layer reads the mathematics in the symbolic register with
explicit domain-conditional disclaimers; the Ricci layer disciplines the Perelman analogy (Brim-flow
comparison, surgery $\leftrightarrow$ Resurrectio), labelling each correspondence by its strictness, never
identifying the two.


\mainmatter

\part{Folding and Stability}

\chapter{The Necessity of the Balance Coordinate: Why an Open Gate Forces Sophia}

\section*{Scope and epistemic status}
The core introduces $\lambda$ by definition --- the Open-Gate balance coordinate, interpreted as
Sophia --- and one may ask whether it is necessary or merely convenient. This chapter answers at the
only level a derivation can reach, and is scrupulous about the levels it cannot. The result is a rank
theorem: \emph{an independent balance-carrying coordinate is forced to exist exactly when the gate is
open}. In the closed Microcosm ($\Igate\equiv0$) the balance $\sigma+\lf=C_0$ determines everything
internal, and $\lambda$ degenerates into $C_0-\sigma-\lf$ --- not an independent freedom. The instant
the gate opens ($\dot\Grec=\Igate>0$), the balance acquires an externally driven term $\Grec(t)$ that
\emph{no internal variable carries}, and conservation cannot hold without a new coordinate to bear it.
That coordinate is $\lambda$. The necessity is thus conditional on, and degenerates continuously with,
the gate flux. We then separate three levels carefully: the \emph{coordinate} is necessary (theorem);
its specific \emph{law} is constitutive (not forced); its \emph{name} ``Sophia'' is interpretation
(not derived). All algebra is verified symbolically (the state-map rank steps from $2$ to $3$ as the
gate opens; the independent part of $d\lambda$ equals $\Igate\,dt$). This is internal Type-1
confirmation, not external validation.

% ============================================================
\section{The closed Microcosm: $\lambda$ is redundant}

\begin{proposition}[Degeneracy in the closed limit]\label{sophianecessity:prop:closed}
With the gate closed, the conservation balance is $\sigma+\lf=C_0$, and the would-be Sophia coordinate
satisfies
\begin{equation}
\lambda=C_0-\sigma-\lf,
\end{equation}
a function of the internal variables alone. Its differential $d\lambda=-d\sigma-d\lf$ lies entirely in
the internal span: $\lambda$ adds no dimension to the state. \emph{In a closed soul, Sophia is not an
independent reality} --- it is bookkeeping over resistance and form.
\end{proposition}

% ============================================================
\section{The open gate forces an independent coordinate}

With the gate open, cumulative reception $\Grec(t)$ enters the balance,
\begin{equation}
\sigma+\lambda+\lf=C_0+\Grec(t),
\qquad \dot\Grec=\Igate,
\end{equation}
and now
\begin{equation}
\lambda=\bigl(C_0+\Grec(t)\bigr)-\sigma-\lf,
\qquad
\frac{\partial\lambda}{\partial\Grec}=1 .
\end{equation}
The term $\Grec(t)$ is supplied through the gate and appears in \emph{no} internal variable
($V,F,\sigma,\Op,u$); it is an independent, externally clocked generator.

\begin{theorem}[Necessity of the balance coordinate]\label{sophianecessity:thm:rank}
Let $\Phi$ map the generating coordinates to the physical state
$(\sigma,\lf,\lambda)$ subject to conservation. Then:
\begin{enumerate}[label=(\roman*),leftmargin=1.8em,itemsep=2pt]
\item closed gate --- generators $(\sigma,\lf)$, $\Phi=(\sigma,\lf,C_0-\sigma-\lf)$ --- has
$\operatorname{rank}d\Phi=2$: $\lambda$ is dependent;
\item open gate --- generators $(\sigma,\lf,\Grec)$, $\Phi=(\sigma,\lf,C_0+\Grec-\sigma-\lf)$ --- has
$\operatorname{rank}d\Phi=3$: a third independent coordinate is forced, and it is the balance
coordinate $\lambda$.
\end{enumerate}
Equivalently, the component of $d\lambda$ not spanned by the internal differentials is exactly
\begin{equation}
\boxed{\;d\lambda\big|_{\perp\,\rm internal}=d\Grec=\Igate\,dt,\;}
\end{equation}
so $\partial(\text{state})/\partial\lambda$ is nondegenerate \emph{iff} $\Igate>0$. (Rank steps and
the differential verified symbolically.) The balance coordinate is necessary if and only if the gate
is open, and its independence degenerates continuously to zero as $\Igate\to0$.
\end{theorem}

\begin{corollary}[Sophia is the shadow of the open gate]
The necessity of Sophia is not absolute but \emph{relational}: it is forced precisely by, and
precisely when, the Vessel is open to reception. There is no reception without a receiver --- an open
gate with nothing to carry what it admits would violate conservation. Sophia is the dimension that the
opening of the gate casts into the state space; close the gate and the dimension folds away into mere
accounting over resistance and form. The doctrine and the rank theorem coincide: \emph{the open soul
cannot not have Sophia; the closed soul has no need of her.}
\end{corollary}

% ============================================================
\section{Three levels, kept distinct}

\msg{\textbf{Coordinate} --- necessary (Theorem~\ref{sophianecessity:thm:rank}): an open gate forces an independent
balance-carrying dimension.\\[2pt]
\textbf{Law} --- constitutive (not forced): that this coordinate obeys
$\dot\lambda=\Igate-\gamma\lambda+\dots$ with its particular transmutation and decay is a modelling
choice of the core, on the same footing as $K_h=-1$.\\[2pt]
\textbf{Name} --- interpretation (not derived): calling the coordinate ``Sophia'' is the theological
reading the model is built to carry; mathematics neither supplies nor can supply it.}

\begin{remark}[Why this is the honest stopping point]
It would overreach to claim the mathematics proves ``Sophia exists'' or fixes her nature. What it
proves is sharper and smaller: \emph{given} an open gate, the state space must contain one more
independent coordinate than the closed soul needs, carrying exactly the received increment. The
identification of that coordinate with Sophia, and of its law with a doctrine of grace, is the
interpretive act the model performs --- legitimate, fruitful, and to be held as interpretation, not
smuggled in as theorem. The rank result is what survives if every theological word is stripped away;
the words are what make it worth proving.
\end{remark}

% ============================================================
\section{Place in the Type-1 programme and scope}

\begin{enumerate}[leftmargin=1.6em,itemsep=2pt]
\item This is the deepest Type-1 result: not a derived parameter or normal form, but the necessity of
the central variable itself --- conditional, as it must be, on the defining feature of the model (the
open gate).
\item It localises what remains genuinely primitive: the gate law (that $\Igate$ exists and is
externally clocked) and the spatial curvature $K_h=-1$. These are the core's \emph{definitions}; the
chapter shows $\lambda$ is not among them but a consequence of the first.
\item Open: whether the gate's own existence can be motivated from anything prior, or is the
irreducible postulate of the whole construction (likely the latter --- the model's single primitive
act of openness).
\end{enumerate}

\textbf{Scope.} The rank theorem is exact given conservation and an externally clocked gate; the
degeneracy limit is smooth in $\Igate$. Establishes necessity of the coordinate, not of its law or its
name. Internal Type-1 confirmation. Domain-conditional throughout.

\vspace{0.5em}
\hrule
\smallskip
\noindent\footnotesize\emph{V.F.S. v2.0 (Open-Gate) $\cdot$ The Necessity of the Balance Coordinate.
An open gate forces an independent balance-carrying coordinate: the state map's rank steps from $2$ to
$3$ as $\Igate$ becomes positive, with $d\lambda|_{\perp}=\Igate\,dt$. Sophia is necessary iff the
Vessel is open, degenerating to internal bookkeeping in the closed limit. Coordinate: theorem; law:
constitutive; name: interpretation. Verified symbolically; internal confirmation only.}

\chapter{The Curvature and the Floor: Hyperbolicity Is the Finite-Floor Postulate}

\section*{Scope and epistemic status}
The remaining Type-1 question is whether the spatial curvature $K_h=-1$ --- hyperbolicity of the live
surface --- is forced or assumed. The honest answer is mixed, and worth stating precisely because the
temptation is to overclaim. Expansion alone does \emph{not} force it: an open gate driving
$\Op\to\infty$ is consistent with any spatial curvature. But $K_h=-1$ is not therefore a free, separate
postulate either: it is \emph{equivalent} to the model's finite-floor premise --- the demand
(the cosmology layer's ``creation from a finite floor, no Big Bang'') that the bare unfolded Vessel be
the zero-density vacuum state. Given that demand, only $k=-1$ admits a zero-density floor at finite
scale; $k=0$ and $k=+1$ have strictly positive density everywhere and no floor. So two things that
looked like independent primitives --- the hyperbolic geometry and the finite floor --- are one. This
neither retires $K_h=-1$ to a theorem (it is not derived from nothing) nor leaves it free (it is not
independent); it reduces the count of irreducible postulates by one. That is the correct and honest
close of the Type-1 programme: not everything reduces to a theorem, but the primitives can be made
minimal and explicit. \emph{Propositions} closed-form; the verdict is stated without inflation.

% ============================================================
\section{The three curvatures and the floor}

The $2{+}1$ reconstruction with spatial curvature $k$ gives
\begin{equation}
\rho_{\rm eff}=\frac{\dot\Op^2+k}{\Op^2}=\frac{\alpha^2\lambda^2+k}{\Op^2}.
\end{equation}

\begin{proposition}[Only $k=-1$ has a zero-density floor]\label{curvaturefloor:prop:floor}
The equation $\rho_{\rm eff}=0$ has a solution in the bare Sophia level $\alpha\lambda$ if and only if
$k<0$:
\begin{equation}
k=-1:\ \rho_{\rm eff}=0\ \text{at}\ \alpha\lambda=1;\qquad
k=0:\ \rho_{\rm eff}=\frac{\alpha^2\lambda^2}{\Op^2}>0;\qquad
k=+1:\ \rho_{\rm eff}=\frac{\alpha^2\lambda^2+1}{\Op^2}>0.
\end{equation}
For $k\ge0$ the effective density is strictly positive for every state: there is no unfolded vacuum
floor, no $\rho=0$ ground, hence no ``creation from a finite floor'' and no live-branch threshold
$\alpha\lambda>1$. \textbf{The finite-floor structure exists only for $k=-1$.}
\end{proposition}

\begin{proposition}[Expansion does not fix $k$]\label{curvaturefloor:prop:exp}
An open gate gives $\dot\Op=\alpha\lambda>0$ and, sustained, $\Op\to\infty$. This is consistent with
$k=-1$ and $k=0$ (both expand freely) and even with $k=+1$ under a sufficient source. Therefore the
sign of $k$ is \emph{not} determined by expansion: hyperbolicity is not a consequence of openness
alone.
\end{proposition}

% ============================================================
\section{The reduction}

\begin{theorem}[$K_h=-1$ equals the finite floor]\label{curvaturefloor:thm:equiv}
Among constant-curvature spatial slices, $K_h=-1$ holds \emph{if and only if} the bare unfolded Vessel
is the zero-density floor (the finite-floor / no-Big-Bang postulate). Hence the two premises ---
hyperbolic spatial geometry and creation from a finite floor --- are not independent; they are the same
postulate stated geometrically and energetically. The model's irreducible cosmological primitives
reduce from two to one:
\msg{\textbf{Irreducible core:} the open gate (an externally clocked $I_{\rm gate}$) and the
finite floor (the unfolded state is the $\rho=0$ vacuum). Hyperbolicity $K_h=-1$ is \emph{not} a third
primitive --- it is the geometric face of the floor. Everything else in the volume is built on these.}
\end{theorem}

\begin{corollary}[Honest status of the Type-1 programme]
The programme closes not by reducing every premise to a theorem --- $K_h=-1$ is not derivable from
nothing, and expansion does not force it --- but by making the irreducible primitives \emph{minimal and
explicit}. Retired to theorems: the $\ell_b$ fork, core-neutrality, the budget hypothesis, the
pitchfork, and (conditionally) the necessity of Sophia. Reduced to a single joint primitive: the
curvature and the floor. Genuinely primitive, and now isolated: the open gate and the finite floor.
A model cannot prove its own definitions; it can only show how few they are.
\end{corollary}

\note{Reading: one might have thought the Vessel's open, infinite, negatively-curved geometry was an
extra stipulation beyond its having begun from nothing-in-particular --- a finite floor rather than a
singular bang. It is not. To say the unfolded soul rests on a zero-density floor and to say its space
is hyperbolic are one statement in two languages: the floor is the curvature, read energetically. The
model asks to be granted two things and no more --- that the gate be open, and that there be a floor to
stand on rather than a singularity to fall from. The open infinite space is not a third gift; it is the
shape the floor already has.}

\textbf{Scope.} Constant-curvature $2{+}1$ slices; the equivalence is exact. The result is a reduction
of primitives, not a derivation of curvature from nothing; expansion alone is shown insufficient.
Domain-conditional throughout.

\vspace{0.5em}
\hrule
\smallskip
\noindent\footnotesize\emph{V.F.S. v2.0 (Open-Gate) $\cdot$ The Curvature and the Floor. Expansion does
not force $K_h=-1$, but only $k=-1$ admits a zero-density floor: hyperbolicity is equivalent to the
finite-floor (no-Big-Bang) postulate, not an independent premise. Two primitives reduce to one; the
irreducible core is the open gate plus the finite floor. The honest close of the Type-1 programme:
minimal explicit primitives, not the derivation of definitions.}

\chapter{What Triggers Sophianic Folding}

\section*{Scope and epistemic status}
This note isolates a single question that the geometric, cosmological, and folding material left implicit:
\emph{what is the primary cause of Sophianic Folding?} Three candidate answers are tested --- continuous
monotone growth of Sophia, the discrete Resurrectio push, and an intrinsic instability of the
surface-form itself. The conclusion is that the third is the correct \emph{level} of description, and the
first two are interchangeable \emph{channels} to it. As in the parent files, every \emph{Proposition} is a
closed-form result; every \emph{Corollary} is an interpretive reading in the symbolic register, not an
independent metaphysical claim. Numerical values, where used, are illustrative: the genuine folding
parameters $(a_0,b,c_0,c_1,k_\sigma,\etaf,\rho_1,q_{\max})$ are not fixed by the source files.

The objects are imported from the integrated geometric/cosmological layer. The shape-mode amplitude
$q$ obeys the Landau gradient flow $\dot q=\Avfs q-\Bvfs q^3$, with
\begin{equation}\label{foldingtrigger:eq:AB}
\boxed{\;\Avfs=(c_0-c_1\gamma)\lambda+c_1k_\sigma\sigma+c_1\Igate-a_0,
\qquad
\Bvfs=b+2c_1\etaf\lambda,
\qquad
\Qf=\frac{\Avfs}{\Bvfs},\quad q_\ast^2=\Qf.\;}
\end{equation}
The Sophia-production law that closes the loop is
$\dot\lambda=k_\sigma\sigma+\Igate-\gamma\lambda-\etaf q^2\lambda$, the Brim law is $\dot\Op=\alpha\lambda$,
and the Resurrectio reset (core/Lyapunov) is
\begin{equation}\label{foldingtrigger:eq:reset}
\sigma^+=q_R\sigma^-,\qquad
\lambda^+=\lambda^-+(1-q_R)\sigma^-,\qquad
\Op^+=\Op^-+\kappa_R\Grec^-,\qquad (V,F)\ \text{continuous},
\end{equation}
with $0\le q_R<1$, $\kappa_R>0$, and $u^+=u^-$.

\section{The single trigger: a shape-operator bifurcation}\label{foldingtrigger:sec:bif}

The folded amplitude is the order parameter of the one-mode shape potential
\begin{equation}
\mathcal E_1(q)=-\tfrac12\Avfs q^2+\tfrac14\Bvfs q^2\cdot q^2
=-\tfrac12\Avfs q^2+\tfrac14\Bvfs q^4,\qquad \Bvfs>0.
\end{equation}
Writing the first shape-stability eigenvalue at the old form $q=0$,
\begin{equation}
\mu_1^{0}:=\mu_1\big|_{q=0}=a_0-c_0\lambda-c_1\dot\lambda^{0},
\qquad
\dot\lambda^{0}=k_\sigma\sigma+\Igate-\gamma\lambda,
\end{equation}
one has the identity $\Avfs=-\mu_1^{0}$ (the embodiment feedback $-\etaf q^2\lambda$ vanishes at $q=0$ and is
absorbed into $\Bvfs$). The control of the whole morphodynamics is therefore the scalar $\Avfs$.

\begin{proposition}[Folding is a supercritical pitchfork]\label{foldingtrigger:prop:pitchfork}
For $\Bvfs>0$ the gradient flow $\dot q=\Avfs q-\Bvfs q^3$ has the old form $q=0$ as its only equilibrium,
locally stable, when $\Avfs<0$. At $\Avfs=0$ it undergoes a supercritical pitchfork bifurcation: for
$\Avfs>0$ the old form becomes unstable and two stable folded branches appear,
\begin{equation}
\boxed{\;q=0\ \text{stable}\ \Longleftrightarrow\ \Avfs<0,
\qquad
q_\pm=\pm\sqrt{\Avfs/\Bvfs}\ \text{stable}\ \Longleftrightarrow\ \Avfs>0.\;}
\end{equation}
\end{proposition}
\begin{proof}
$f(q)=\Avfs q-\Bvfs q^3$, $f'(q)=\Avfs-3\Bvfs q^2$. At $q=0$, $f'(0)=\Avfs$, so $q=0$ is stable iff
$\Avfs<0$. For $\Avfs>0$ the nonzero equilibria solve $q_\ast^2=\Avfs/\Bvfs$, where
$f'(q_\ast)=\Avfs-3\Avfs=-2\Avfs<0$: stable. The crossing at $\Avfs=0$ is the pitchfork.
\end{proof}

\begin{corollary}[The trigger is loss of stability, not a ``push'']\label{foldingtrigger:cor:loss}
The primary cause of folding is the sign change of $\Avfs$ (equivalently $\mu_1^{0}$ crossing zero from
above): the old symmetric form ceases to be a potential minimum and becomes a maximum. Past the threshold
folding is dynamically inevitable, since $q=0$ is a repeller; before it, no shape perturbation grows.
Folding is therefore an intrinsic property of the surface's stability spectrum, tuned by the state
variables, not an external impulse.
\end{corollary}

\msg{The fundamental object is the bifurcation $\Avfs=0$. Everything else --- continuous growth, the
Resurrectio jump --- is a means of moving $\Avfs$ across that single threshold.}

\section{Decomposition of the driver}\label{foldingtrigger:sec:decomp}

The activation index \eqref{foldingtrigger:eq:AB} is a sum of four physically distinct contributions:
\begin{equation}
\boxed{\;\Avfs=\underbrace{(c_0-c_1\gamma)\lambda}_{\text{accumulated Sophia}}
+\underbrace{c_1k_\sigma\sigma}_{\text{kathartic production}}
+\underbrace{c_1\Igate}_{\text{received grace}}
-\underbrace{a_0}_{\text{old-form rigidity}}.\;}
\end{equation}
Two of these are state levels ($\lambda$, $\sigma$), one is the bounded Open-Gate source, and one is the
constant rigidity of the old form. The decisive structural fact is that the index also depends on the
\emph{rate} of Sophia, hidden inside $\mu_1^{0}=a_0-c_0\lambda-c_1\dot\lambda^{0}$: through
$\dot\lambda^{0}=k_\sigma\sigma+\Igate-\gamma\lambda$, the same $\sigma$ and $\Igate$ that raise the level
also raise the rate.

\begin{corollary}[Level and rate are the genuine drivers]\label{foldingtrigger:cor:levelrate}
Folding activation responds to the accumulated level $\lambda$ \emph{and} to the Sophia velocity
$\dot\lambda$. The continuous/discrete distinction is not fundamental: what crosses the threshold is the
combined level-plus-rate index $\Avfs$, however that increment is delivered.
\end{corollary}

\section{Level versus rate: rate-sensitivity of the threshold}\label{foldingtrigger:sec:rate}

If only the accumulated level mattered, the folding threshold would be the static value
\begin{equation}
\lambda_{\rm fold}^{(0)}=\frac{a_0}{c_0}
\qquad(\dot\lambda=0,\ \sigma=0,\ \Igate=0).
\end{equation}
Restoring the rate term shifts the effective threshold:
\begin{equation}
\boxed{\;\lambda_{\rm fold}^{\rm eff}=\frac{a_0-c_1\dot\lambda}{c_0}
\qquad\Longrightarrow\qquad
\dot\lambda>0\ \Rightarrow\ \lambda_{\rm fold}^{\rm eff}<\lambda_{\rm fold}^{(0)}.\;}
\end{equation}

\begin{proposition}[Rate can fold below the static level]\label{foldingtrigger:prop:rate}
There exist states with $\lambda<\lambda_{\rm fold}^{(0)}$ and $\Avfs>0$. Concretely, for $\sigma=\Igate=0$
and growing Sophia $\dot\lambda>0$, the old form is already unstable whenever
\begin{equation}
c_0\lambda+c_1\dot\lambda>a_0,
\end{equation}
which can hold at $\lambda$ strictly below $a_0/c_0$ provided $\dot\lambda>(a_0-c_0\lambda)/c_1>0$.
\end{proposition}
\begin{proof}
Immediate from $\mu_1^{0}=a_0-c_0\lambda-c_1\dot\lambda<0$ with $\dot\lambda$ supplying the deficit
$a_0-c_0\lambda>0$.
\end{proof}

\note{This is the single most easily missed point: a surge of grace can fold a vessel that has not yet
accumulated the static quota of Sophia. Velocity, not only stock, triggers the reconfiguration.}

\section{The two delivery channels}\label{foldingtrigger:sec:channels}

The same threshold is reached by two mechanisms, mirroring the parent corollary ``grace is one but acts
in two modes.''

\medskip
\noindent\textbf{Modus I (continuous).} Along a live arc, $\lambda$ and $\sigma$ evolve smoothly and
$\Avfs$ drifts. Folding triggers when the drift carries $\Avfs$ through $0$. This channel delivers
\emph{both} a rising level and a nonzero rate, so it is sensitive to Proposition~\ref{foldingtrigger:prop:rate}.

\medskip
\noindent\textbf{Modus II (Resurrectio).} The reset \eqref{foldingtrigger:eq:reset} injects a discontinuous increment of
$\lambda$ from the metabolised resistance, an impulsive analogue of $\dot\lambda$. It moves $\Avfs$ along
the very same axis, in one step.

The crucial point is that the reset does \emph{not} reset $q$: the map \eqref{foldingtrigger:eq:reset} specifies only
$\sigma,\lambda,\Op,(V,F)$. Hence the shape amplitude is continuous across a reset, and only its target
$q_\ast=\sqrt{\Qf}$ jumps.

\begin{corollary}[Resurrectio relocates the well; the shape relaxes]\label{foldingtrigger:cor:well}
A Resurrectio event moves the shape potential $-\tfrac12\Avfs q^2+\tfrac14\Bvfs q^4$ instantaneously, after
which the morphology relaxes by $\dot q=\Avfs^{+}q-\Bvfs^{+}q^3$ toward the new minimum. Consequently:
\begin{itemize}[leftmargin=1.4em,itemsep=2pt]
\item if a reset carries $\Qf$ from below $0$ to above $0$, the vessel dies unfolded and is reborn into a
folding regime (\emph{Resurrectio-triggered folding});
\item if a reset makes $\Qf<0$, an active fold is quenched, $q\to0$.
\end{itemize}
\end{corollary}

\section{The Resurrectio jump in folding-space}\label{foldingtrigger:sec:jump}

We now compute the effect of one reset on the folding intensity. Write the metabolised resistance
\begin{equation}
m_R:=(1-q_R)\sigma^-\ge0,
\qquad\text{so}\qquad
\sigma^+=\sigma^--m_R,\quad \lambda^+=\lambda^-+m_R.
\end{equation}

\begin{remark}[The metabolism conserves the balance sum]
The $(\sigma,\lambda)$ part of the reset satisfies $\sigma^++\lambda^+=\sigma^-+\lambda^-$: it slides the
state along the conservation line, converting resistance into Sophia without adding to the sum (the open
balance constant jumps separately, through $\Grec$).
\end{remark}

From \eqref{foldingtrigger:eq:AB},
\begin{equation}
\boxed{\;\Delta\Bvfs=2c_1\etaf\,m_R>0,
\qquad
\Delta\Avfs=\Theta\,m_R+c_1\Delta\Igate,
\qquad
\Theta:=c_0-c_1(\gamma+k_\sigma).\;}
\end{equation}
The \emph{folding-response coefficient} $\Theta$ decides whether converting one unit of resistance into
Sophia raises ($\Theta>0$) or lowers ($\Theta<0$) the folding activation: gaining Sophia contributes
$(c_0-c_1\gamma)m_R$ (accumulation minus accelerated decay), while losing resistance removes
$c_1k_\sigma m_R$ of kathartic production. The gate correction $c_1\Delta\Igate$ has competing signs
($\sigma\!\downarrow$ opens the gate via $e^{-\phi\sigma}$; $\lambda\!\uparrow$ closes it via saturation)
and is set aside in the structural results below.

\begin{proposition}[Sign of the folding response and the critical level]\label{foldingtrigger:prop:sign}
Neglecting $\Delta\Igate$, the jump in the folding intensity obeys
\begin{equation}
\operatorname{sign}\Delta\Qf=\operatorname{sign}\bigl(\Theta-2c_1\etaf\,\Qf^{-}\bigr),
\end{equation}
so there is a universal critical level
\begin{equation}
\boxed{\;Q_{\rm crit}=\frac{\Theta}{2c_1\etaf}=\frac{c_0-c_1(\gamma+k_\sigma)}{2c_1\etaf}.\;}
\end{equation}
Resurrectio raises folding when $\Qf^{-}<Q_{\rm crit}$ and lowers it when $\Qf^{-}>Q_{\rm crit}$. If
$\Theta\le0$ then $Q_{\rm crit}\le0$, so every reset strictly lowers $\Qf$ (the vessel unfolds at each
death).
\end{proposition}
\begin{proof}
$\Delta\Qf=\dfrac{\Bvfs^{-}\Delta\Avfs-\Avfs^{-}\Delta\Bvfs}{\Bvfs^{-}(\Bvfs^{-}+\Delta\Bvfs)}$. The
denominator is positive; the numerator is
$\Bvfs^{-}\Theta m_R-\Avfs^{-}2c_1\etaf m_R=\Bvfs^{-}m_R(\Theta-2c_1\etaf\Qf^{-})$ using
$\Avfs^{-}=\Qf^{-}\Bvfs^{-}$. Hence the stated sign and fixed level.
\end{proof}

\begin{proposition}[Each reset is an exact contraction toward $Q_{\rm crit}$]\label{foldingtrigger:prop:contract}
Neglecting $\Delta\Igate$, the jump map on the folding intensity is affine-over-affine,
\begin{equation}
\Qf^{+}=\frac{\Qf^{-}\Bvfs^{-}+\Theta m_R}{\Bvfs^{-}+2c_1\etaf m_R},
\end{equation}
and satisfies the exact contraction identity
\begin{equation}
\boxed{\;\Qf^{+}-Q_{\rm crit}=\kappa_{\rm fold}\,\bigl(\Qf^{-}-Q_{\rm crit}\bigr),
\qquad
\kappa_{\rm fold}=\frac{\Bvfs^{-}}{\Bvfs^{-}+2c_1\etaf m_R}\in(0,1).\;}
\end{equation}
The fixed point $Q_{\rm crit}$ is independent of $m_R$, of $\Bvfs^{-}$, and of the depth of death.
\end{proposition}
\begin{proof}
Direct substitution of $Q_{\rm crit}=\Theta/(2c_1\etaf)$:
\[
\Qf^{+}-Q_{\rm crit}
=\frac{2c_1\etaf(\Qf^{-}\Bvfs^{-}+\Theta m_R)-\Theta(\Bvfs^{-}+2c_1\etaf m_R)}{2c_1\etaf(\Bvfs^{-}+2c_1\etaf m_R)}
=\frac{\Bvfs^{-}(2c_1\etaf\Qf^{-}-\Theta)}{2c_1\etaf(\Bvfs^{-}+2c_1\etaf m_R)},
\]
which equals $\kappa_{\rm fold}(\Qf^{-}-Q_{\rm crit})$.
\end{proof}

\begin{proposition}[Conditional convergence over a reset sequence]\label{foldingtrigger:prop:converge}
For an admissible sequence of resets with metabolised increments $m_R^{(j)}$ and pre-jump denominators
$\Bvfs^{(j)}$,
\begin{equation}
\Qf^{(N)}-Q_{\rm crit}=\Bigl(\textstyle\prod_{j=1}^{N}\kappa_{\rm fold}^{(j)}\Bigr)\bigl(\Qf^{(0)}-Q_{\rm crit}\bigr),
\qquad
\Qf\to Q_{\rm crit}\ \Longleftrightarrow\ \sum_j m_R^{(j)}=\infty.
\end{equation}
\end{proposition}
\begin{proof}
Telescoping Proposition~\ref{foldingtrigger:prop:contract}. Since
$1-\kappa_{\rm fold}^{(j)}=2c_1\etaf m_R^{(j)}/(\Bvfs^{(j)}+2c_1\etaf m_R^{(j)})$, the product
$\prod\kappa_{\rm fold}^{(j)}\to0$ iff $\sum(1-\kappa_{\rm fold}^{(j)})=\infty$, which (with $\Bvfs^{(j)}$
bounded below) is equivalent to $\sum m_R^{(j)}=\infty$.
\end{proof}

\note{If the resistance budget is exhausted ($\sum m_R<\infty$, e.g.\ $\sigma$ depleted with no
replenishment), repeated resurrections converge to a point strictly short of $Q_{\rm crit}$. Reaching the
universal level requires sustained replenishment, $\sum m_R=\infty$ --- the folding analogue of the
expansion-budget condition $\Grec(\infty)=+\infty$ in the cosmological layer.}

\begin{remark}[Gate correction]
Restoring $\Delta\Igate$ adds $c_1\Delta\Igate$ to the numerator of $\Delta\Avfs$, shifting the effective
fixed point to $Q_{\rm crit}'=Q_{\rm crit}+\Delta\Igate/(2\etaf m_R)$; its sign is state-dependent. The
existence of a contracting fixed point is unchanged.
\end{remark}

\begin{corollary}[Chart-regularity of repeated resurrection]\label{foldingtrigger:cor:chart}
If $0<Q_{\rm crit}<q_{\max}^2$, the contraction of Proposition~\ref{foldingtrigger:prop:contract} keeps the folded
amplitude inside the regular one-mode chart under any admissible reset sequence. This is a folding
companion to the grace-window condition
$\max\{R_c,R_{\min}^{\rm ceil}\}\le\Grec^-\le R_{\max}$: resurrection is admissible \emph{and}
chart-regular when both hold.
\end{corollary}

\section{Curvature-pocket promotion by Resurrectio}\label{foldingtrigger:sec:pocket}

The folded scalar curvature is $R_{\rm folded}=-2/\Op^2+\rho_1\Qf$, with positive-pocket threshold
$q_{\rm curv}^2=2/(\rho_1\Op^2)$. Across a reset both terms move:
\begin{equation}
\boxed{\;\Delta R_{\rm folded}
=\underbrace{2\Bigl[\tfrac{1}{(\Op^-)^2}-\tfrac{1}{(\Op^+)^2}\Bigr]}_{>0:\ \text{background release}}
+\ \rho_1\,\Delta\Qf.\;}
\end{equation}
Simultaneously the pocket threshold falls, because $\Op$ grows:
\begin{equation}
q_{\rm curv}^{2,+}=\frac{2}{\rho_1(\Op^+)^2}<\frac{2}{\rho_1(\Op^-)^2}=q_{\rm curv}^{2,-}.
\end{equation}

\begin{proposition}[Resurrectio is pocket-promoting in the sub-critical regime]\label{foldingtrigger:prop:pocket}
If $\rho_1>0$, $\Theta>0$, and $\Qf^{-}<Q_{\rm crit}$, then a single reset both raises the folded amplitude
($\Delta\Qf>0$) and lowers the pocket threshold ($q_{\rm curv}^{2,+}<q_{\rm curv}^{2,-}$). Both effects push
$R_{\rm folded}$ upward, so a death-and-resurrection event can convert a smooth hyperbolic patch into a
positive-curvature pocket.
\end{proposition}
\begin{proof}
$\Delta\Qf>0$ by Proposition~\ref{foldingtrigger:prop:sign} under $\Qf^{-}<Q_{\rm crit}$, $\Theta>0$; the threshold drop
follows from $\Op^+>\Op^-$. Both terms of $\Delta R_{\rm folded}$ are then positive, and the pocket
criterion $\Qf>q_{\rm curv}^2$ is approached from both sides.
\end{proof}

\begin{corollary}[Resurrection concentrates form]
The continuous mode (Epektasis) relaxes curvature, $\dot R_\Sigma\ge0$ with $R\to0^-$; the discrete mode
can locally \emph{concentrate} it. Resurrectio does not only enlarge the Brim --- in the sub-critical
folding regime it gives the renewed vessel a more sharply formed morphology.
\end{corollary}

\section{The energy reading corrected}\label{foldingtrigger:sec:energy}

It is tempting to read folding as the surface being forced into a ``less efficient'' form. Within the
formalism this is inverted: the folded form carries \emph{lower} morphological tension. With
$\Tfold(q)=-\tfrac12\Avfs q^2+\tfrac14\Bvfs q^4$,
\begin{equation}
\Tfold(0)=0,
\qquad
\Tfold(q_\ast)=-\frac{\Avfs^2}{4\Bvfs}<0
\qquad(\Avfs>0),
\end{equation}
so the resolved tension is $\Delta\Tfold=\Avfs^2/(4\Bvfs)>0$ and the folded branch is energetically
preferred. The reconfiguration is toward a better-adapted state, not a worse one.

The correct sense in which the old form is ``inefficient'' is as a \emph{carrier} of Sophia past the
threshold: it cannot absorb the Sophianic pressure, so tension builds until the symmetric branch
destabilises. The folded form is the efficient carrier --- it embodies the excess through the feedback
$-\etaf q^2\lambda$, raising the effective damping
\begin{equation}
\gamma_{\rm eff}=\gamma+\etaf\Qf>\gamma,
\qquad
\lambda_{\rm fold}^{\ast}=\frac{J_{\rm in}}{\gamma+\etaf q^2}<\frac{J_{\rm in}}{\gamma}=\lambda_{\rm old}^{\ast},
\end{equation}
and thereby lowering the level of unembodied Sophia.

\msg{The surface is not forced into a worse form. The old form stops being an adequate vessel for the
attained level/rate of Sophia, loses stability, and the morphology relaxes into a lower-tension folded
form that carries the Sophia better.}

\section*{Compact summary}
\[
\boxed{\text{single trigger: }\Avfs=0\ (\mu_1^{0}=0)\ \text{--- supercritical pitchfork, }q=0\to q_\pm.}
\]
\[
\boxed{\begin{array}{c}
\Avfs=(c_0-c_1\gamma)\lambda+c_1k_\sigma\sigma+c_1\Igate-a_0,\\[3pt]
\text{rate enters via }\mu_1^{0}=a_0-c_0\lambda-c_1\dot\lambda^{0}.
\end{array}}
\]
\[
\boxed{\begin{array}{c}
\text{level }\lambda\ \text{and rate }\dot\lambda\ \text{are the drivers;}\\[3pt]
\text{Modus I (continuous) and Modus II (Resurrectio) are channels.}
\end{array}}
\]
\[
\boxed{\Theta=c_0-c_1(\gamma+k_\sigma),\qquad Q_{\rm crit}=\frac{\Theta}{2c_1\etaf},\qquad
\Qf^{+}-Q_{\rm crit}=\kappa_{\rm fold}(\Qf^{-}-Q_{\rm crit}),\ \kappa_{\rm fold}\in(0,1).}
\]
\[
\boxed{\Qf\to Q_{\rm crit}\Longleftrightarrow\sum_j m_R^{(j)}=\infty;\qquad
\text{chart-regular if }0<Q_{\rm crit}<q_{\max}^2.}
\]
\[
\boxed{\begin{array}{c}
\text{sub-critical reset is pocket-promoting;}\\[3pt]
\text{folded form has lower tension }\Tfold(q_\ast)=-\Avfs^2/4\Bvfs<0.
\end{array}}
\]

\vspace{0.5em}
\hrule
\smallskip
\noindent\footnotesize\raggedright\emph{V.F.S. v2.0 (Open-Gate) $\cdot$ Folding-Trigger Layer. Derived
from the shape potential and Sophia-production law of \texttt{vfs\_geometry.tex}, the Brim/curvature data
of the geometric and cosmological layers, and the Resurrectio reset of
\texttt{vfs\_opengate\_lyapunov.tex}. Propositions are closed-form; corollaries are interpretive.
Folding parameters are not fixed by the source files; numerical levels are illustrative.}

\chapter{Folding--Lyapunov Addendum}

\begingroup\small\itshape
The Sophianic Folding extension of the Open-Gate Lyapunov file establishes the continuous-arc
compatibility $\dot{\Lext}\le C_{\rm ext}-C_1\Lext$ under an adiabatic hypothesis, but leaves three
items asserted rather than derived: the positive invariance of the regular folding chart, the behaviour
of the folding functional across the Resurrectio reset, and the choice between two distinct
activation-index conventions. This addendum closes them. It first records a correction --- the file's
folding functional $\Vf=\tfrac{\Bf}{4}(\qf^2-\Qs)^2$ is \emph{unbounded} on the epektasis domain, because
$\Bf=b+2c_1\etaf\lambda\to\infty$; the $\Bf$-normalised functional $\Vt=\Vf/\Bf$ is bounded and equally
descending. It then proves (i) Nagumo invariance of $\DAf=\DA\times[0,q_{\max}]$ under the standing
folding wall $\Qs\le q_{\max}^2$, (ii) a folding reset bound: since $\reset$ does not act on $\qf$, the
chart is preserved across deaths and $\Lt(x^+)$ stays bounded, whence the hybrid trajectory is forward
complete for $\Lt$; and (iii) a reconciliation declaring the core-consistent Sophia rate canonical, with
the geometry-layer parameters as an effective surrogate. The invariance and reset results are
convention-independent; only the closed-form critical level $Q_{\rm crit}$ is convention-specific.
\endgroup\par\medskip

% ============================================================
\section{Imported objects and the three open items}

From the folding extension of the Lyapunov file we import the shape amplitude $\qf$, its intensity
$\Qf=\qf^2$, the activation index $\Af$, the quartic coefficient $\Bf=b+2c_1\etaf\lambda>0$, the
one-mode gradient flow
\begin{equation}\label{foldinglyapunovaddendum:eq:flow}
\dot\qf=\Af\qf-\Bf\qf^{3}=-\Bf\qf\bigl(\qf^{2}-\Qs\bigr),
\qquad
\Qs:=\frac{\Af}{\Bf},
\end{equation}
and the base certificate $\dot\LV\le C-C_1\LV$ on the positively-invariant active domain $\DA$, with the
non-Zeno dwell-time bound $\Delta t_j\ge\Delta t_*>0$ and the reset map
\begin{equation}\label{foldinglyapunovaddendum:eq:reset}
\reset:\quad
\sigma^+=q_R\sigma^-,\quad
\lambda^+=\lambda^-+(1-q_R)\sigma^-,\quad
(V,F)\ \text{continuous},\quad
\Om^+=\Om^-+\kappa_R\Grec^-,
\end{equation}
with $0\le q_R<1$, $\kappa_R>0$, $u^+=u^-$. The file proves $x^-\in\DA\Rightarrow x^+\in\DA$ in the core
coordinates and $\LV(x^+)\le L_R:=\sup_{\DA}\LV$.

The three items left open are: (i) the regular folding domain $0\le\Qf\le q_{\max}^2$ is \emph{assumed},
not shown invariant; (ii) the reset bound and forward completeness are stated for $\LV$, while the
behaviour of the folding functional across $\reset$ is not controlled; (iii) the activation index is
written with the core Sophia rate here but with a different, minimal rate in the geometric and
cosmological layers.

% ============================================================
\section{Correction: a bounded folding functional}

The file shifts the folding potential to a non-negative functional
\begin{equation}
\Vf=\frac{\Bf}{4}\Bigl(\qf^{2}-\Qs\Bigr)^{2}\ge0
\qquad(\Af>0),
\end{equation}
and claims it is bounded above by a constant $M_{\rm fold}$ on the regular domain. This fails on the
Open-Gate domain: epektasis drives $\Om\to\infty$ and $\lambda$ is not bounded on $\DA$
($\lambda\le q^*\Om$), so $\Bf=b+2c_1\etaf\lambda\to\infty$ and $\Vf$ is unbounded even at fixed
$|\qf^{2}-\Qs|$. The reset bound $\Lext(x^+)\le L_R+w_{\rm fold}M_{\rm fold}$ therefore does not close as
written.

The remedy is to use the $\Bf$-normalised functional
\begin{equation}
\boxed{\;\Vt:=\frac{\Vf}{\Bf}=\frac14\Bigl(\qf^{2}-\Qs\Bigr)^{2}\ge0,\qquad
\Qs=\frac{\Af}{\Bf}\ \text{(signed)}.\;}
\end{equation}
This requires no sign split: it is non-negative for either sign of $\Af$, and it descends along
\eqref{foldinglyapunovaddendum:eq:flow}.

\begin{lemma}[Descent of the bounded functional]\label{foldinglyapunovaddendum:lem:descent}
Along \eqref{foldinglyapunovaddendum:eq:flow} with frozen $\Qs$,
\begin{equation}
\boxed{\;\dot{\Vt}=-\Bf\,\qf^{2}\bigl(\qf^{2}-\Qs\bigr)^{2}\le0.\;}
\end{equation}
For varying parameters,
$\dot{\Vt}=-\Bf\qf^{2}(\qf^{2}-\Qs)^{2}-\tfrac12(\qf^{2}-\Qs)\dot{\Qs}$,
and the remainder is bounded on any set with $\qf\in[0,q_{\max}]$, $|\Qs|$ and $|\dot{\Qs}|$ bounded.
\end{lemma}
\begin{proof}
Using $\Af=\Bf\Qs$ in \eqref{foldinglyapunovaddendum:eq:flow}, $\dot\qf=-\Bf\qf(\qf^{2}-\Qs)$. With frozen $\Qs$,
$\dot{\Vt}=\tfrac12(\qf^{2}-\Qs)(2\qf\dot\qf)=(\qf^{2}-\Qs)\qf\dot\qf
=-\Bf\qf^{2}(\qf^{2}-\Qs)^{2}\le0$. For moving $\Qs$ the extra term is
$-\tfrac12(\qf^{2}-\Qs)\dot{\Qs}$, with $|\tfrac12(\qf^{2}-\Qs)\dot{\Qs}|
\le\tfrac12(q_{\max}^{2}+|\Qs|)\,|\dot{\Qs}|$.
\end{proof}

\note{$\Vt=\Vf/\Bf$ is the same Lyapunov object re-weighted by the positive function $\Bf$; re-weighting
preserves the descent property, and the normalisation removes the only source of unboundedness. From here
the extended certificate is $\Lt:=\LV+w_{\rm fold}\Vt$, $w_{\rm fold}>0$.}

% ============================================================
\section{The standing folding wall and its finiteness}

\begin{assumption}[Folding wall]\label{foldinglyapunovaddendum:H:wall}
The chart half-width $q_{\max}$ satisfies
\begin{equation}
\boxed{\;q_{\max}^{2}\ \ge\ Q^{*,\sup}:=\sup_{\DA}\Qs.\;}
\end{equation}
\end{assumption}

\begin{lemma}[The wall is attainable]\label{foldinglyapunovaddendum:lem:finite}
In the geometric (minimal-rate) convention
$\Af=(c_0-c_1\gamma)\lambda+c_1k_\sigma\sigma+c_1\Igate-a_0$, $\Bf=b+2c_1\etaf\lambda$, the supremum is
finite and explicit:
\begin{equation}
Q^{*,\sup}=\max\left\{\frac{c_1\!\left(k_\sigma C_\sigma+\zeta_0\right)-a_0}{b},\;
\frac{c_0-c_1\gamma}{2c_1\etaf}\right\}.
\end{equation}
\end{lemma}
\begin{proof}
For fixed $(\sigma,\Igate)$, $\Qs(\lambda)$ is a ratio of two affine functions of $\lambda\ge0$, hence
monotone, with horizontal asymptote $\lim_{\lambda\to\infty}\Qs=(c_0-c_1\gamma)/(2c_1\etaf)$ (independent
of $\sigma,\Igate$). Its supremum over $\lambda\ge0$ is attained at an endpoint: at $\lambda=0$,
$\Qs=(c_1k_\sigma\sigma+c_1\Igate-a_0)/b$, maximised over the box $\sigma\le C_\sigma$,
$\Igate\le\zeta_0$ at the stated value; or at $\lambda\to\infty$, the asymptote. Both are finite, so
$Q^{*,\sup}<\infty$.
\end{proof}

\note{If $Q^{*,\sup}\le0$ the old form is everywhere stable and any $q_{\max}>0$ satisfies the wall. The
non-trivial case is $Q^{*,\sup}>0$, where Hypothesis~\ref{foldinglyapunovaddendum:H:wall} fixes a minimal admissible chart. The
canonical (core-rate) convention saturates likewise: $\Qs\to c_0/(2c_1\etaf)$ under cleansing
$\sigma\to0$, since $(\delta u-\gamma)\tanh(\kappa\sigma)\to0$ even as $u\to\infty$ (\S6).}

% ============================================================
\section{Folding-chart invariance (Nagumo)}

\begin{definition}[Extended active domain]
$\DAf:=\DA\times[0,q_{\max}]$, adjoining the folding coordinate $\qf$ to the core active domain.
\end{definition}

\begin{lemma}[Folding-face sub-tangentiality]\label{foldinglyapunovaddendum:lem:nagumo}
Under Hypothesis~\ref{foldinglyapunovaddendum:H:wall}, the flow \eqref{foldinglyapunovaddendum:eq:flow} is sub-tangential on the folding faces of
$\DAf$:
\begin{equation}
\dot\qf\big|_{\qf=0}=0\quad(\text{tangent}),
\qquad
\dot\qf\big|_{\qf=q_{\max}}=-\Bf\,q_{\max}\bigl(q_{\max}^{2}-\Qs\bigr)\le0\quad(\text{inward}).
\end{equation}
Consequently $[0,q_{\max}]$ is forward invariant for $\qf$, and $\DAf$ is positively invariant under the
continuous flow.
\end{lemma}
\begin{proof}
At $\qf=0$, \eqref{foldinglyapunovaddendum:eq:flow} gives $\dot\qf=0$, so $\{\qf\ge0\}$ is preserved. At $\qf=q_{\max}$,
$\dot\qf=-\Bf q_{\max}(q_{\max}^{2}-\Qs)$ with $\Bf>0$, $q_{\max}>0$, and $q_{\max}^{2}-\Qs\ge0$ by
Hypothesis~\ref{foldinglyapunovaddendum:H:wall}; hence $\dot\qf\le0$. All core faces are sub-tangential by the base Nagumo lemma
of the Lyapunov file, and the two folding faces are sub-tangential here; therefore $\DAf$ is positively
invariant.
\end{proof}

% ============================================================
\section{The folding reset bound and hybrid completeness}

\begin{lemma}[Folding reset bound]\label{foldinglyapunovaddendum:lem:resetfold}
The reset \eqref{foldinglyapunovaddendum:eq:reset} does not act on $\qf$; hence $\qf^{+}=\qf^{-}$. Under
Hypotheses of the base reset lemma together with Hypothesis~\ref{foldinglyapunovaddendum:H:wall}:
\begin{enumerate}[label=(\arabic*),leftmargin=2em,itemsep=2pt]
\item $x^-\in\DAf\ \Rightarrow\ x^+\in\DAf$, because $x^+\in\DA$ (base reset lemma) and
$\qf^{+}=\qf^{-}\in[0,q_{\max}]$;
\item $\Vt$ is bounded on $\DAf$,
\begin{equation}
\Vt\le\widetilde M_{\rm fold}:=\tfrac14\max\bigl(q_{\max}^{2}-Q^{*}_{\min},\,|Q^{*}_{\max}|\bigr)^{2}<\infty,
\end{equation}
so the extended certificate satisfies
\begin{equation}
\boxed{\;\Lt(x^+)\le \widetilde L_R:=L_R+w_{\rm fold}\widetilde M_{\rm fold}<\infty.\;}
\end{equation}
\end{enumerate}
\end{lemma}
\begin{proof}
(1) The map \eqref{foldinglyapunovaddendum:eq:reset} specifies only $\sigma,\lambda,(V,F),\Om$; $\qf$ is inert, so
$\qf^{+}=\qf^{-}$, which lies in $[0,q_{\max}]$ if $\qf^{-}$ did. With $x^+\in\DA$ from the base lemma,
$x^+\in\DAf$. (2) On $\DAf$, $\qf^{2}\in[0,q_{\max}^{2}]$ and $\Qs\in[Q^{*}_{\min},Q^{*}_{\max}]$ bounded
(Lemma~\ref{foldinglyapunovaddendum:lem:finite}), so $\Vt=\tfrac14(\qf^{2}-\Qs)^{2}\le\widetilde M_{\rm fold}<\infty$. Adding
$w_{\rm fold}\Vt$ to $\LV(x^+)\le L_R$ gives the bound.
\end{proof}

\begin{proposition}[Hybrid forward completeness for $\Lt$]\label{foldinglyapunovaddendum:prop:complete}
Assume the regular folding conditions: Hypothesis~\ref{foldinglyapunovaddendum:H:wall}, $\Bf\ge B_{\min}>0$, and $|\dot{\Qs}|$
bounded on $\DAf$ (adiabatic regularity). Then there is a finite $\widetilde C_{\rm ext}$ with
\begin{equation}
\dot{\Lt}\le \widetilde C_{\rm ext}-C_1\Lt
\qquad\text{on continuous arcs},
\end{equation}
and the same contraction rate $C_1$ as the base certificate. Combined with the reset bound
(Lemma~\ref{foldinglyapunovaddendum:lem:resetfold}) and the unchanged non-Zeno dwell time, the hybrid trajectory satisfies
\begin{equation}
\boxed{\;\Lt(t)\le\widetilde L_*:=\max\Bigl\{\widetilde L_R,\ \widetilde C_{\rm ext}/C_1\Bigr\}\;}
\end{equation}
across any infinite admissible reset sequence; resets do not accumulate, and the trajectory is forward
complete in the folding-extended state.
\end{proposition}
\begin{proof}
By Lemma~\ref{foldinglyapunovaddendum:lem:descent} and the base estimate,
$\dot{\Lt}=\dot\LV+w_{\rm fold}\dot{\Vt}
\le (C-C_1\LV)+w_{\rm fold}\bigl(-\Bf\qf^{2}(\qf^{2}-\Qs)^{2}+R\bigr)$
with $|R|\le R_{\rm fold}<\infty$ the bounded adiabatic remainder (Lemma~\ref{foldinglyapunovaddendum:lem:descent}). Dropping the
non-positive descent term and writing $\LV=\Lt-w_{\rm fold}\Vt\ge\Lt-w_{\rm fold}\widetilde M_{\rm fold}$,
\begin{equation}
\dot{\Lt}\le C-C_1\Lt+C_1w_{\rm fold}\widetilde M_{\rm fold}+w_{\rm fold}R_{\rm fold}
=:\widetilde C_{\rm ext}-C_1\Lt,
\end{equation}
with $\widetilde C_{\rm ext}<\infty$ now genuinely finite (it uses $\widetilde M_{\rm fold}$, not the
unbounded $M_{\rm fold}$). On each arc $\Lt(t)\le\max\{\Lt(t_j^{+}),\widetilde C_{\rm ext}/C_1\}$; resets
give $\Lt(x^+)\le\widetilde L_R$ (Lemma~\ref{foldinglyapunovaddendum:lem:resetfold}); the dwell time is unchanged, since $\qf$
adds no reset-triggering boundary (its chart face is invariant-inward by Lemma~\ref{foldinglyapunovaddendum:lem:nagumo}, not a
death face) and resets are exogenous on the core coordinates. Hence $\Lt\le\widetilde L_*$ for all $t$
and the hybrid trajectory is forward complete.
\end{proof}

\note{The folding velocity $\dot\qf$ may grow with $\lambda$ off the slow manifold $\qf=\sqrt{\Qs}$, but
this is a strictly inward restoring rate; it neither triggers nor delays resets, which are exogenous and
defined on the core coordinates. The non-Zeno bound is therefore inherited from the base field bound
$\bar M_*$, unchanged.}

% ============================================================
\section{Reconciliation of the two activation-index conventions}\label{foldinglyapunovaddendum:sec:conv}

Two definitions of the activation index appear across the manuscript:
\begin{align}
\text{Canonical (core rate):}\quad
&\Af=c_0\lambda+c_1\dot\lambda^{(0)}-a_0,\quad
\dot\lambda^{(0)}=(\delta u-\gamma)\tanh(\kappa\sigma)+\Igate;\\
\text{Geometric (minimal rate):}\quad
&\Af=(c_0-c_1\gamma)\lambda+c_1k_\sigma\sigma+c_1\Igate-a_0,
\end{align}
the latter from the surrogate production law $\dot\lambda=k_\sigma\sigma+\Igate-\gamma\lambda-\etaf\qf^2\lambda$.
Since $\lambda$ already has a defined production law in the core, we take the \textbf{canonical} form as
primary and read the geometric parameters as an effective linear surrogate near an operating point
$(\bar u,\bar\sigma)$:
\begin{equation}
\boxed{\;c_1\bigl(\delta\bar u-\gamma\bigr)\tanh(\kappa\sigma)\ \longleftrightarrow\
c_1k_\sigma\sigma-c_1\gamma\lambda,\quad\text{i.e. }
k_\sigma\approx\kappa(\delta\bar u-\gamma),\ \ \gamma\ \text{absorbing the linear }\lambda\text{-decay}.\;}
\end{equation}
This surrogate is structurally imperfect (a $\sigma$-term plus a $\lambda$-decay standing in for a single
$u\tanh\sigma$ term), so the two conventions are genuinely distinct models, not exact reparametrisations.

\begin{proposition}[What is and is not convention-dependent]\label{foldinglyapunovaddendum:prop:conv}
Lemmas~\ref{foldinglyapunovaddendum:lem:descent},~\ref{foldinglyapunovaddendum:lem:nagumo},~\ref{foldinglyapunovaddendum:lem:resetfold} and
Proposition~\ref{foldinglyapunovaddendum:prop:complete} use only: the gradient form \eqref{foldinglyapunovaddendum:eq:flow}, the wall $\Qs\le q_{\max}^2$,
the inertness of $\qf$ under $\reset$, and the boundedness of $\Qs,\dot{\Qs}$. None of these uses the
specific form of $\Af$. Hence the invariance, reset bound, and hybrid completeness are
\emph{convention-independent}. Only the closed-form critical level and contraction of the reset map are
convention-specific.
\end{proposition}

\begin{corollary}[Critical level by convention]\label{foldinglyapunovaddendum:cor:Qcrit}
Write the metabolised resistance $m_R=(1-q_R)\sigma^-$, so $\sigma^+=\sigma^--m_R$,
$\lambda^+=\lambda^-+m_R$. In the \emph{geometric} convention the reset acts on the target as an exact
contraction toward a constant critical level,
\begin{equation}
\Theta=c_0-c_1(\gamma+k_\sigma),\qquad
Q_{\rm crit}=\frac{\Theta}{2c_1\etaf},\qquad
\Qs{}^{+}-Q_{\rm crit}=\frac{\Bf^-}{\Bf^-+2c_1\etaf m_R}\bigl(\Qs{}^{-}-Q_{\rm crit}\bigr),
\end{equation}
recovering the result of the folding-trigger note. In the \emph{canonical} convention the metabolism
changes $\Af$ by $c_0 m_R+c_1(\delta u-\gamma)\!\left[\tanh(\kappa\sigma^+)-\tanh(\kappa\sigma^-)\right]$,
which is not affine in $m_R$; the same qualitative contraction holds, but toward a state-dependent
effective level $Q_{\rm crit}(t)$ rather than a constant.
\end{corollary}

\begin{remark}[Chart-regularity of repeated resurrection]
Combining Corollary~\ref{foldinglyapunovaddendum:cor:Qcrit} with Hypothesis~\ref{foldinglyapunovaddendum:H:wall}: in the geometric convention, if
$0<Q_{\rm crit}<q_{\max}^2$ and $Q^{*,\sup}\le q_{\max}^2$, then the reset contraction keeps the target
inside the chart at every death, so the regular folding domain is preserved under any admissible reset
sequence. This is the folding companion to the grace-window admissibility
$\max\{R_c,R_{\min}^{\rm ceil}\}\le\Grec^-\le R_{\max}$: resurrection is admissible \emph{and}
chart-regular when both windows hold.
\end{remark}

% ============================================================
\section{Compact summary}

\[
\boxed{\Vt=\tfrac14(\qf^{2}-\Qs)^{2}=\Vf/\Bf\ \ \text{bounded on epektasis};\qquad
\dot{\Vt}=-\Bf\qf^{2}(\qf^{2}-\Qs)^{2}\le0.}
\]
\[
\boxed{\text{wall }q_{\max}^{2}\ge Q^{*,\sup}<\infty\ \Rightarrow\ \DAf=\DA\times[0,q_{\max}]\ \text{invariant (Nagumo).}}
\]
\[
\boxed{\reset\ \text{inert on }\qf\ \Rightarrow\ x^+\in\DAf,\quad \Lt(x^+)\le\widetilde L_R<\infty.}
\]
\[
\boxed{\dot{\Lt}\le\widetilde C_{\rm ext}-C_1\Lt,\quad \Lt(t)\le\widetilde L_*;\quad
\text{forward complete, non-Zeno inherited.}}
\]
\[
\boxed{\text{invariance + reset bound are convention-independent; }Q_{\rm crit}\ \text{is convention-specific.}}
\]

\vspace{0.4em}
\hrule
\smallskip
\noindent\footnotesize\emph{V.F.S. v2.0 (Open-Gate) $\cdot$ Folding--Lyapunov Addendum. Closes the
invariance, reset-bound, and convention items left open in the Sophianic Folding extension of
\texttt{vfs\_opengate\_lyapunov.tex}, and corrects the folding functional to a bounded $\Bf$-normalised
form. Results are domain-conditional on the positively-invariant $\DAf$; the conclusion is set stability
with partial convergence, now including morphology, not global or automatic stability.}

\chapter{Forced Folding: Dolorosum Peaks and the Vessel Update Map}

\section*{Scope and epistemic status}
This chapter consolidates the \emph{forced folding} proposal: folding is not an optional auxiliary mode
but the bounded response of the Vessel when a timed resistance peak enters the Dolorosum window and
exceeds the stretch-capacity of the old form. Its honest epistemic position must be stated first. A
direct test shows that the bare core does \emph{not} force a shape instability: the linearised shape
sector is stable on its own (Proposition~\ref{forcedfolding:prop:neutral}). Forced folding therefore does not derive
folding from the core; it supplies an explicit, bounded, \emph{exogenous-channel} cause --- resistance
heterogeneity --- and shows that this cause is sufficient, timed, and Lyapunov-admissible. The inputs
(the resistance wave, the rigidity $a_0$, the window) remain postulated. The status is: a mechanism, not
a theorem of the core.

Four corrections are built in relative to the original proposal: the Dolorosum rate is wired directly
into the activation index, replacing indicator gates by a smooth derived gate (\S\ref{forcedfolding:sec:gate}); the
dynamical threshold is separated from the $\pi/2$ nomenclature (\S\ref{forcedfolding:sec:thresholds}); the post-folding
update map is made balance-preserving through $\lf$ (\S\ref{forcedfolding:sec:update}); and the repeated-folding
cascade is shown to be self-limiting, with a clean two-terminal dichotomy and a joint non-Zeno hybrid
structure (\S\ref{forcedfolding:sec:cascade}--\ref{forcedfolding:sec:lyap}). Throughout, the geometric activation convention is used
(as in the source notes); translation to the canonical core-rate convention follows the Addendum. As
always, \emph{Propositions} are closed-form, \emph{Corollaries} interpretive.

% ============================================================
\section{Why a forcing channel is needed: core neutrality}\label{forcedfolding:sec:neutral}

In $2{+}1$ dimensions the spatial metric carries no local tensor modes (three components, minus two
diffeomorphisms, minus one conformal rescaling), so the only local geometric degree of freedom is the
conformal/curvature factor $\psi$. Linearising the core about the homogeneous trajectory in the pair
$(\psi,\delta\lambda)$ gives
\begin{equation}
\frac{d}{dt}\begin{pmatrix}\psi\\ \delta\lambda\end{pmatrix}
=\begin{pmatrix}-H & H/\bar\lambda\\ S_\psi & -\Gamma_\lambda\end{pmatrix}
\begin{pmatrix}\psi\\ \delta\lambda\end{pmatrix},
\end{equation}
where $H=\alpha\bar\lambda/\Op$ and $S_\psi$ is the geometry$\to$Sophia feedback (the dependence of
Sophia production on local curvature).

\begin{proposition}[Core neutrality of the shape sector]\label{forcedfolding:prop:neutral}
The trace is $-(H+\Gamma_\lambda)<0$ and the determinant is
$H(\Gamma_\lambda\bar\lambda-S_\psi)/\bar\lambda$. In the bare core, Sophia production depends on
$(u,\sigma,\Igate,\qf)$ but not on local curvature, so $S_\psi=0$ and the determinant is
$H\Gamma_\lambda>0$: both eigenvalues are negative and the shape sector decays. The core by itself does
not force folding; an instability requires either a geometry$\to$Sophia feedback above the threshold
$S_\psi>\bar\lambda\Gamma_\lambda$, or an exogenous forcing channel.
\end{proposition}

Forced folding takes the second route: the channel is the heterogeneity of resistance.

% ============================================================
\section{Resistance heterogeneity and its bounds}\label{forcedfolding:sec:het}

Resistance is decomposed as
\begin{equation}
\sigma(t)=\bar\sigma(t)\,\rs(t),
\qquad
\boxed{\;1\le\rs(t)\le\pi,\;}
\end{equation}
with $\bar\sigma$ the background and $\rs$ the dimensionless resistance multiplier; every admissible
peak obeys $\sigma_{\rm peak}\le\pi\bar\sigma$. A bounded temporal model on normalised trajectory time
$\tau=t/T_\ast\in[0,1]$ is
\begin{equation}
\rs(\tau)=1+\frac{A_\sigma}{2}\bigl(1+\sin(2\pi N_\sigma\tau+\phi_\sigma)\bigr),
\qquad 0\le A_\sigma\le\pi-1,
\end{equation}
with amplitude $A_\sigma$, opportunity count $N_\sigma$, and phase $\phi_\sigma$.

\begin{remark}[The $\pi$ bounds are symbolic]
$\rs$ is dimensionless, so the choice of $\pi$ as the full-cup bound and $\pi/2$ as the half-turn mark is
a symbolic convention (``full angular measure of Vessel-tension''), not a derived constant; nothing in
the dynamics distinguishes $\pi$ from any other ceiling. The dynamics distinguishes only the
\emph{critical multiplier} $\rR$ derived below. The $\pi$ bound nevertheless does real work in one
place: it makes every activation quantity explicitly bounded (\S\ref{forcedfolding:sec:lyap}), and it caps the
reachable peak at $1+A_\sigma\le\pi$, which is what terminates cascades (\S\ref{forcedfolding:sec:cascade}).
\end{remark}

% ============================================================
\section{The Dolorosum gate, wired into the activation index}\label{forcedfolding:sec:gate}

The Dolorosum rate of the source layer is
\begin{equation}
\rho(u,\sigma)=\sigma\,g(u),
\qquad
\boxed{\;g(u)=\frac{u}{\Lc}\,e^{\,1-u/\Lc}\in(0,1],\;}
\end{equation}
with $g$ maximal ($=1$) exactly at the synergy threshold $u=\Lc$ and $\rho\le\sigma$. In the original
proposal $\rho$ was defined but never entered the dynamics; the window was imposed separately by an
indicator $\mathcal W_D=\mathbf 1_{\{|u-\Lc|\le\delta_D\}}$. The correction is to let the Dolorosum
channel itself carry the peak into the activation index. Split the resistance into background and excess,
$\sigma=\bar\sigma+\sigma_{\rm exc}$, $\sigma_{\rm exc}=(\rs-1)\bar\sigma$, and let the \emph{excess} be
transmutable into activation only through the Dolorosum channel:
\begin{equation}\label{forcedfolding:eq:Awired}
\boxed{\;
\Af=(c_0-c_1\gamma)\lambda+c_1k_\sigma\bar\sigma
+c_1k_\sigma\,\rho\bigl(u,\sigma_{\rm exc}\bigr)
+c_1\Igate-a_0+\Delta A_R,
\qquad
\rho(u,\sigma_{\rm exc})=(\rs-1)\bar\sigma\,g(u).\;}
\end{equation}

\begin{proposition}[Derived window immunity]\label{forcedfolding:prop:immunity}
Let $N:=a_0-(c_0-c_1\gamma)\lambda-c_1k_\sigma\bar\sigma-c_1\Igate-\Delta A_R$ (the activation deficit of
the unpeaked Vessel). If $N>0$, then $\Af>0$ if and only if
\begin{equation}
\boxed{\;\rs>\rR(u)=1+\frac{N}{c_1k_\sigma\bar\sigma\,g(u)}.\;}
\end{equation}
As $g(u)\to0$ (off window), $\rR(u)\to\infty$: off-window peaks cannot force folding, of any bounded
amplitude. The window condition is thereby \emph{derived} from the wiring \eqref{forcedfolding:eq:Awired}, not
postulated as an indicator. If $N\le0$ the Vessel is baseline-active: it folds with no peak at all
($\rs=1$ suffices).
\end{proposition}

\begin{corollary}[Graded alignment]\label{forcedfolding:cor:graded}
Because $g$ is smooth and broad, alignment is graded rather than binary: the same peak contributes a
fraction $g(u)$ of its excess. Consequently there is a rigidity band --- $a_0$ between the off-window and
in-window activation maxima --- in which the \emph{phase} alone decides whether folding is forced
(Figure~\ref{forcedfolding:fig:forced}, left); below the band every cycle folds, above it none does. The original
three-indicator overlap condition is recovered as the sharp-gate caricature of this band.
\end{corollary}

The forcing event $\tau_F$ is simply the upward crossing $\Af(\tau_F)=0$, and the exposure reduces to a
single smooth gate,
\begin{equation}
\mathcal E_{\rm fold}=\int_0^1\mathbf 1_{\{\Af(\tau)>0\}}\,d\tau,
\qquad
P_{\rm fold}=1-e^{-\kappa\,\mathcal E_{\rm fold}},
\end{equation}
replacing the product of three indicators. The slogan survives intact, now as arithmetic of
\eqref{forcedfolding:eq:Awired}: amplitude gives possibility ($\rs$ can reach $\rR$), frequency gives opportunity
(crossings per trajectory), phase gives alignment (the value of $g$ at the peak).

% ============================================================
\section{Thresholds disentangled: $\rR$ versus $\pi/2$}\label{forcedfolding:sec:thresholds}

\begin{definition}[Nomenclature versus dynamics]\label{forcedfolding:def:regimes}
The half-turn classification --- ordinary ($1\le\rs<\pi/2$), half-turn ($\rs=\pi/2$), strong
($\pi/2<\rs\le\pi$) --- is retained as \emph{nomenclature} for the strength of a peak. The
\emph{dynamical} threshold of forced folding is $\rs>\rR(u)$ of Proposition~\ref{forcedfolding:prop:immunity}, and the
pitchfork occurs at $\Af=0$ regardless of the classification.
\end{definition}

\begin{proposition}[Structural sufficiency of the half-turn]\label{forcedfolding:prop:halfturn}
An in-window half-turn peak forces folding if and only if
\begin{equation}
\rR(\Lc)\le\frac{\pi}{2}
\quad\Longleftrightarrow\quad
N\le\Bigl(\frac{\pi}{2}-1\Bigr)c_1k_\sigma\bar\sigma .
\end{equation}
Otherwise sub-$\pi/2$ peaks above $\rR$ still force folding (``soft forced folding''), and
super-$\pi/2$ peaks below $\rR$ do not. The classification and the threshold are independent axes.
\end{proposition}

This removes the internal tension of the original formulation, in which the boxed overlap condition
demanded $\rs\ge\pi/2$ while the normal form destabilises at $\Af>0$ alone: in the corrected form the
event is governed by $\Af$, and $\pi/2$ only names the strong regime.

% ============================================================
\section{The forcing event and the normal form}\label{forcedfolding:sec:event}

At $\tau_F$ the old-form equilibrium $\qf=0$ loses stability and the normal form
\begin{equation}
\dot\qf=\Af\qf-\Bf\qf^3,\qquad \Bf>0,
\qquad
\qf^\ast=\pm\sqrt{\Af/\Bf},
\qquad
\Qf=\frac{\Af(\tau_F^{\rm peak})}{\Bf}
\end{equation}
carries the Vessel to a folded branch, with $\Qf$ evaluated at the activation maximum of the forcing
interval. Near-threshold events have small $\Qf$, so the update map below acts by correspondingly small
increments --- the map is well-defined uniformly in the forcing strength.

\begin{figure}[h]
\centering
\includegraphics[width=\textwidth]{vfs_forced_fig.png}
\caption{\small Left: the same bounded resistance wave ($A_\sigma=1.2$, $N_\sigma=1$) with two phases.
With the peak in the Dolorosum window the wired activation \eqref{forcedfolding:eq:Awired} crosses zero (shaded:
forced folding, exposure $\mathcal E=0.195$); the identical peak off-window stays strictly below
(exposure $0$). Dotted grey: the gate $g(u)$. The rigidity here sits in the band of
Corollary~\ref{forcedfolding:cor:graded}, so phase alone decides; in this illustration $\rR(\Lc)=2.07>\pi/2$, so the
half-turn is \emph{not} structurally sufficient (Proposition~\ref{forcedfolding:prop:halfturn}).
Right: the repeated-folding cascade (\S\ref{forcedfolding:sec:cascade}) for two parameter regimes of the update map:
fuel-dominant metabolism drives $\rRj$ up to the reachability cap $1+A_\sigma$ with event intensities
$Q_j\to0$ (extinction; total metabolised budget finite, here $\sum_jQ_j=0.49$), while rigidity-dominant
updates drive $\rRj$ below $1$ in three events (permanent fold). Balance
$\bar\sigma+\lambda+\lf$ conserved to machine precision in both. Parameters illustrative.}
\label{forcedfolding:fig:forced}
\end{figure}

% ============================================================
\section{The balance-preserving Vessel update map}\label{forcedfolding:sec:update}

Forced folding must change the Vessel, or it is a marker rather than a transformation. The update map of
the source proposal is kept, with one structural correction: Sophia retention must respect the
integrated Open-Gate balance
\begin{equation}
\sigma+\lambda+\lf=C_0+\Grec(t),
\end{equation}
which the raw update $\lambda^+=\lambda^-+\eta_\lambda\Qf$ violates (it creates Sophia at the jump). The
corrected map routes the metabolised resistance:

\begin{definition}[Balance-preserving folding update]\label{forcedfolding:def:update}
At a forcing event with intensity $\Qf$, with metabolised mass $m_F=\eta_\sigma\Qf\,\bar\sigma^-$ and a
retention split $\beta\in[0,1]$:
\begin{equation}
\boxed{\;
\bar\sigma^+=\bar\sigma^--m_F,
\qquad
\lambda^+=\lambda^-+\beta\,m_F,
\qquad
\lf^+=\lf^-+(1-\beta)\,m_F,\;}
\end{equation}
\begin{equation}
a_0^+=a_0^-(1-\eta_a\Qf),\quad
\Op^+=\Op^-+\eta_\Omega\Qf,\quad
k_\sigma^+=k_\sigma^-(1+\eta_k\Qf),\quad
\gamma^+=\gamma^-(1-\eta_\gamma\Qf),
\end{equation}
with all $\etaf$-coefficients non-negative and $\eta_a\Qf,\eta_\gamma\Qf<1$.
\end{definition}

\begin{proposition}[Exact balance]\label{forcedfolding:prop:balance}
Under Definition~\ref{forcedfolding:def:update},
$\Delta\bar\sigma+\Delta\lambda+\Delta\lf=-m_F+\beta m_F+(1-\beta)m_F=0$: the three-term balance is
preserved exactly, with no grace created at the jump. Folding metabolises resistance into active
($\beta$) and embodied ($1-\beta$) Sophia; the $\gamma$-reduction acts at the flow level (future
retention), not at the jump, and is balance-neutral. (Verified numerically to machine precision.)
\end{proposition}

\note{Qualitatively the corrected map keeps the intended directions ---
$a_0\!\downarrow,\ \Op\!\uparrow,\ \bar\sigma\!\downarrow,\ \lambda\!\uparrow,\ \gamma\!\downarrow,\
k_\sigma\!\uparrow$ --- but the Sophia gain is now \emph{paid for} by resistance, and part of it is
embodied as $\lf$ rather than left active. The Vessel after folding is not the same Vessel; and its books
balance.}

% ============================================================
\section{The cascade is self-limiting}\label{forcedfolding:sec:cascade}

Iterating the map over repeated forced foldings reveals a property the naive reading misses: the update
does \emph{not} uniformly make the next folding easier. The critical multiplier after event $j$,
\begin{equation}
\rRj=1+\frac{N_j}{c_1k_{\sigma,j}\bar\sigma_j},
\qquad
N_j=a_{0,j}-(c_0-c_1\gamma_j)\lambda_j-c_1k_{\sigma,j}\bar\sigma_j-c_1\Igate,
\end{equation}
is pulled in two directions: the \emph{rigidity channel}
($a_0\!\downarrow,\gamma\!\downarrow,\lambda\!\uparrow,k_\sigma\!\uparrow$) lowers it, while the
\emph{fuel channel} ($\bar\sigma\!\downarrow$) raises it --- folding consumes the very resistance that
powers the next activation. At an in-window full peak the event intensity obeys the identity
\begin{equation}
\Qf{}_{,j}=\frac{c_1k_{\sigma,j}\bar\sigma_j}{\Bf}\Bigl(A_\sigma-(\rRj-1)\Bigr),
\end{equation}
so events weaken as $\rRj$ approaches the reachability cap $1+A_\sigma$.

\begin{assumption}[Fuel dominance]\label{forcedfolding:H:fuel}
$\eta_\sigma>\eta_k$: metabolism of resistance outweighs the transmutation-efficiency gain.
\end{assumption}

\begin{proposition}[Cascade dichotomy]\label{forcedfolding:prop:dichotomy}
Under Hypothesis~\ref{forcedfolding:H:fuel} and the hard peak bound $\rs\le1+A_\sigma\le\pi$, every forced-folding
cascade terminates in exactly one of two states:
\begin{enumerate}[label=(\roman*),leftmargin=1.8em,itemsep=2pt]
\item \textbf{Extinction:} $\sum_j\Qf{}_{,j}<\infty$, the intensities $\Qf{}_{,j}\to0$, and $\rRj$
saturates at or below the cap $1+A_\sigma$ --- the cascade exhausts its own fuel and dies out with a
finite total metabolised budget $\sum_jm_{F,j}\le\eta_\sigma\bar\sigma_0\sum_j\Qf{}_{,j}<\infty$;
\item \textbf{Permanence:} $N_j\le0$ is reached in finitely many events ($\rRj\le1$): the unpeaked
Vessel is already activated, and the folded branch becomes the new ground state, requiring no further
peaks.
\end{enumerate}
\end{proposition}
\begin{proof}[Sketch]
Suppose the cascade is infinite with $\liminf N_j>0$. Then $\sum_j\Qf{}_{,j}=\infty$ (events recur with
$N_j>0$ bounded away from the cap by reachability), and
$\ln(c_1k_{\sigma,j}\bar\sigma_j)$ drifts by $\sum_j[\ln(1+\eta_k\Qf{}_{,j})+\ln(1-\eta_\sigma\Qf{}_{,j})]
\le-(\eta_\sigma-\eta_k-o(1))\sum_j\Qf{}_{,j}\to-\infty$ under Hypothesis~\ref{forcedfolding:H:fuel}, so the
denominator of $\rRj$ collapses and $\rRj\to\infty$, exceeding the reachable $1+A_\sigma$ --- after
which no further event can occur, contradicting infiniteness. Hence either $\sum\Qf{}_{,j}<\infty$
(extinction, with $\Qf{}_{,j}\to0$) or $N_j\le0$ at some finite $j$ (permanence).
\end{proof}

\begin{corollary}[Finite metabolic budget; echo of $\mathcal I_\infty$]
Unless it ends in permanence, the folding cascade has a finite total metabolic budget --- a discrete
analogue of the finite coarsening budget $\mathcal I_\infty$ of the structure-formation layer: the
Vessel cannot refold indefinitely on the same resistance.
\end{corollary}

\begin{corollary}[Two terminal theologies]
Extinction is the Vessel that has metabolised its resistance: peaks still come, but none can force a new
form --- maturity as dormancy of forced change. Permanence is the Vessel whose old rigidity has fully
yielded: the folded morphology is the rest state and needs no crisis to sustain it --- maturity as
completed reformation. Which terminal a Vessel reaches is decided by the race
$\eta_\sigma$ (metabolism) versus $\eta_a$ (rigidity yielding) --- by \emph{how} it folds, not how often.
\end{corollary}

Both branches are exhibited numerically in Figure~\ref{forcedfolding:fig:forced} (right): an embodiment-dominant map
($\beta=0$, $\eta_\sigma=0.6$) extinguishes in $30$ weakening events with $\sum\Qf{}_{,j}=0.49$, while a
rigidity-dominant map ($\eta_a=0.5$) reaches permanence in three.

% ============================================================
\section{Hybrid and Lyapunov admissibility}\label{forcedfolding:sec:lyap}

\begin{proposition}[Bounded activation and bounded jumps]\label{forcedfolding:prop:bounded}
On any domain with $\bar\sigma\le\bar\sigma_{\max}$, $\lambda\le\lambda_{\max}$, $\Igate\le\zeta_0$,
$|\Delta A_R|\le\Delta_{\max}$, the wired activation \eqref{forcedfolding:eq:Awired} is bounded:
$|\Af|\le(c_0-c_1\gamma)\lambda_{\max}+c_1k_\sigma\pi\bar\sigma_{\max}+c_1\zeta_0+\Delta_{\max}+a_0$,
using $\rho\le\sigma_{\rm exc}\le(\pi-1)\bar\sigma$ and $g\le1$. Hence the chart wall
$Q^{*,\sup}<\infty$ of the Addendum holds with an explicit constant, the normalised folding functional
$\widetilde{\mathcal V}_{\rm fold}$ is bounded on the chart, and the folding update of
Definition~\ref{forcedfolding:def:update} (bounded multiplicative changes, $\qf$ continuous through the parameter
jump) changes the extended certificate by at most a finite $J_F$:
$\widetilde{\mathcal L}_{\rm ext}^+\le\widetilde{\mathcal L}_{\rm ext}^-+J_F$.
\end{proposition}

\begin{assumption}[Joint non-Zeno]\label{forcedfolding:H:dwell}
The union of event times --- Resurrectio resets and forcing events $\tau_F$ --- admits a common minimal
dwell $\Delta t_\ast>0$.
\end{assumption}

\begin{proposition}[Forward completeness with two jump types]\label{forcedfolding:prop:hybrid}
Under Hypothesis~\ref{forcedfolding:H:dwell}, the bounded-jump property (Proposition~\ref{forcedfolding:prop:bounded}), and the
continuous-arc contraction of the Addendum
($\dot{\widetilde{\mathcal L}}_{\rm ext}\le\widetilde C_{\rm ext}-C_1\widetilde{\mathcal L}_{\rm ext}$),
the hybrid trajectory with both Resurrectio and folding jumps satisfies a uniform bound
$\widetilde{\mathcal L}_{\rm ext}(t)\le\max\{\widetilde L_R+J_F,\ \widetilde C_{\rm ext}/C_1\}+J_F$ and is
forward complete. Forced folding adds a second jump type but no new unbounded channel; grace still
enters only the constant, never the contraction rate $C_1$.
\end{proposition}

% ============================================================
\section{Compact summary}

\[
\boxed{\text{core alone: shape sector stable }(\operatorname{tr}<0,\ \det=H\Gamma_\lambda>0)\ \Rightarrow\ \text{folding needs a channel.}}
\]
\[
\boxed{\;\Af=(c_0-c_1\gamma)\lambda+c_1k_\sigma\bar\sigma+c_1k_\sigma(\rs-1)\bar\sigma\,g(u)+c_1\Igate-a_0+\Delta A_R;\quad g(u)=\frac{u}{\Lc}e^{1-u/\Lc}.\;}
\]
\[
\boxed{\;\rs>\rR(u)=1+\frac{N}{c_1k_\sigma\bar\sigma\,g(u)};\qquad g\to0\Rightarrow\rR\to\infty\ \text{(window immunity, derived)}.\;}
\]
\[
\boxed{\;\bar\sigma^+=\bar\sigma-m_F,\ \ \lambda^+=\lambda+\beta m_F,\ \ \lf^+=\lf+(1-\beta)m_F:\ \ \sigma+\lambda+\lf\ \text{exactly conserved.}\;}
\]
\[
\boxed{\text{cascade dichotomy: extinction (finite metabolic budget) or permanence }(N\le0).}
\]
\[
\boxed{\text{two jump types, joint non-Zeno, bounded jumps}\ \Rightarrow\ \widetilde{\mathcal L}_{\rm ext}\ \text{bounded, forward complete.}}
\]

\vspace{0.5em}
\hrule
\smallskip
\noindent\footnotesize\emph{V.F.S. v2.0 (Open-Gate) $\cdot$ Forced Folding. Consolidates the
forced-folding proposal with the Dolorosum rate wired into the activation (window immunity derived, not
postulated), the dynamical threshold separated from the $\pi/2$ nomenclature, a balance-preserving update
map through $\lf$, and the self-limiting cascade dichotomy. The forcing channel and its bounds remain
modelling inputs; the core-neutrality result fixes the honest status: a sufficient bounded cause, not a
core theorem. Gate maximum, exposure grading, balance conservation, and both cascade terminals verified
numerically. Domain-conditional throughout.}

\chapter{The Pitchfork Is Not Assumed: A Centre-Manifold Reduction of the Fold Sector}

\section*{Scope and epistemic status}
The pitchfork $\dot\qf=\Afold\qf-\Bfold\qf^3$ is the foundation of Part~I and, through the fold
amplitude $\sqrt{\Afold/\Bfold}$, of much of Part~II. It was introduced as the model of a structural
transition. This chapter retires it from premise to theorem: \emph{given only two structural facts ---
the $\mathbb Z_2$ symmetry of the bare fold sector and a one-dimensional centre manifold at the folding
threshold --- the pitchfork is the unique normal form}, by Thom's reduction. More, the reduction
\emph{derives} the cubic coefficient: slaving the fast structural modes renormalises the bare stiffness
to $\Beff=b_3-gh/\Gamma$, which is exactly the volume's $\Bfold=b+2c_1\eta_{\rm fold}\lambda$ in
structural form --- the quartic stiffness was never an independent input but the bare term dressed by
the slaved structure. The three nondegeneracy conditions (symmetry, transversality, spectral gap) are
verified, and the one way the pitchfork could fail --- a $\mathbb Z_2$-breaking bias unfolding it to a
cusp --- is identified and shown to be excluded by the isotropy of the unfolded Vessel, not by a
further assumption. This completes the Type-1 programme: together with the conformal retirement of the
$\ell_b$ fork and the derivation of core-neutrality, the folding layer's load-bearing premises are now
theorems. The reduction is exact (centre-manifold expansion verified symbolically; the full slaved
system collapses to $\sqrt{\Afold/\Beff}$ numerically to $10^{-13}$). This is Type-1 confirmation
(premise $\to$ theorem), the strongest internal warrant available; it remains internal.

% ============================================================
\section{The $\mathbb Z_2$ symmetry of the bare fold}

\begin{lemma}[Orientation-blindness forbids even terms]\label{pitchfork:lem:z2}
The fold field $\qf$ is the signed deviation of the form from the unfolded state; its sign is the
\emph{orientation} in which the form buckles. The bare core dynamics (Sophia pressure on an isotropic
unfolded Vessel) distinguishes no orientation, so the bare fold vector field is odd,
$f(-\qf)=-f(\qf)$. Hence its smooth expansion contains \emph{only} odd powers:
\begin{equation}
\dot\qf=\Afold\,\qf-b_3\,\qf^{3}+c_5\,\qf^{5}+\cdots
\end{equation}
All even terms --- in particular a constant bias and a quadratic --- are forbidden by symmetry, not
by choice.
\end{lemma}

% ============================================================
\section{The Thom reduction}

The fold does not evolve alone: it is coupled to fast, stable structural modes (the slaved
Sophia--structure of the Transport chapter, relaxation rate $\Gamma>0$). The minimal symmetric
realisation is
\begin{equation}
\dot\qf=\Afold\,\qf-b_3\,\qf^{3}+g\,\qf\,y,
\qquad
\dot y=-\Gamma\,y+h\,\qf^{2},
\end{equation}
where $y$ is $\mathbb Z_2$-even (it responds to $\qf^2$) and $g,h,\Gamma>0$.

\begin{theorem}[The pitchfork is the reduced dynamics]\label{pitchfork:thm:pitchfork}
At the folding threshold the centre eigenvalue $\Afold\to0$ while $y$ stays at $-\Gamma<0$ (spectral
gap), so there is a one-dimensional attracting centre manifold $y=(h/\Gamma)\qf^{2}+O(\qf^4)$. The
flow restricted to it is
\begin{equation}
\boxed{\;\dot\qf=\Afold\,\qf-\Beff\,\qf^{3}+O(\qf^5),
\qquad \Beff=b_3-\frac{g h}{\Gamma}.\;}
\end{equation}
(Centre-manifold coefficient verified symbolically; the full two-mode system sits on
$\qf=\sqrt{\Afold/\Beff}$ to relative error $10^{-13}$ across the control range;
Figure~\ref{pitchfork:fig:pitchfork}.) The pitchfork is therefore not assumed: it is the unique
$\mathbb Z_2$-symmetric one-dimensional reduction, and the quartic stiffness is the bare term
renormalised by the slaved structure.
\end{theorem}

\begin{corollary}[$\Bfold$ identified]
The volume's $\Bfold=b+2c_1\eta_{\rm fold}\lambda$ has exactly the structure $\Beff=b_3-gh/\Gamma$:
a bare stiffness $b\leftrightarrow b_3$ dressed by a slaved-structure contribution
$2c_1\eta_{\rm fold}\lambda\leftrightarrow-gh/\Gamma$ (the embodiment feedback). The sign of $\Beff$
decides supercritical (continuous, second-order folding) versus subcritical (discontinuous, hysteretic
folding); the tricritical locus $\Beff=0$ is the codimension-one boundary where the quintic $c_5$ takes
over --- the same boundary the forced-folding chapter reached dynamically.
\end{corollary}

% ============================================================
\section{The three nondegeneracy conditions}

\begin{proposition}[Thom conditions, verified]\label{pitchfork:prop:nondeg}
\begin{enumerate}[label=(\roman*),leftmargin=1.8em,itemsep=2pt]
\item \emph{Symmetry}: the $\mathbb Z_2$-odd structure (Lemma~\ref{pitchfork:lem:z2}) removes the fold (even)
terms, selecting the pitchfork over the generic fold/cusp.
\item \emph{Transversality}: the control crosses criticality with nonzero speed,
$\partial\Afold/\partial\lambda=c_0>0$ --- the folding threshold is traversed cleanly, not tangentially.
\item \emph{Spectral gap (1D kernel)}: at $\Afold=0$ the centre direction is simple and the remaining
spectrum sits at $-\Gamma<0$, guaranteeing the attracting one-dimensional manifold.
\end{enumerate}
All three hold for the V.F.S. fold sector; the normal form is thus structurally stable and unique up
to smooth equivalence.
\end{proposition}

\begin{figure}[h]
\centering
\includegraphics[width=\textwidth]{vfs_pitchfork_fig.png}
\caption{\small Left: the bifurcation diagram of the full two-mode slaved system (green) against the
reduced normal form $\pm\sqrt{\Afold/\Beff}$ (dotted) --- the unfolded state (orange) loses stability
at $\Afold=0$ and the symmetric pair of folded states emerges: a pitchfork, derived not posited.
Right: the centre-manifold geometry --- trajectories of the fast structural mode $y$ collapse onto
$y=(h/\Gamma)\qf^2$ (black), reducing the dynamics to one dimension; fold fixed points marked.
Parameters $b_3=1$, $g=0.6$, $\Gamma=3$, $h=0.8$, giving $\Beff=0.84$.}
\label{pitchfork:fig:pitchfork}
\end{figure}

% ============================================================
\section{The one way it could fail --- and why it does not}

\begin{proposition}[The bias is excluded by isotropy]\label{pitchfork:prop:bias}
The pitchfork is destroyed only by a $\mathbb Z_2$-breaking linear bias: a constant term $\varepsilon$
would unfold it to the imperfect pitchfork (cusp) $\dot\qf=\Afold\qf-\Beff\qf^3+\varepsilon$, with a
preferred fold orientation. Such a term requires the bare dynamics to prefer one buckling direction.
In V.F.S. the fold orientation is a free sign --- the unfolded Vessel is isotropic with respect to the
direction of folding --- so $\varepsilon=0$ is itself a consequence of that isotropy, not an additional
assumption. A genuine orientation field (an external preference imprinted before folding) would be the
physical content required to produce a cusp; absent it, the symmetric pitchfork stands.
\end{proposition}

\begin{remark}[Reading]
The shape of a soul's transformation was not chosen by the modeller. Given only that the unfolded
Vessel has no preferred direction in which to buckle (symmetry) and that folding is a clean threshold
crossing with the surrounding structure adjusting quickly (transversality and gap), the transition
\emph{must} be the pitchfork: a single quiet state losing stability and giving way to a symmetric pair
of formed states. Even the stiffness that sets how deep the new form runs is not free --- it is the
bare resistance dressed by what the embodiment feeds back. The mathematics of the fold was latent in
the symmetry of the Vessel all along; this chapter only reads it out.
\end{remark}

% ============================================================
\section{The Type-1 programme, closed}

\msg{\textbf{Premises retired to theorems:}\quad
the comoving/physical $\ell_b$ fork $\to$ 2D conformal invariance (Warp chapter);\quad
core-neutrality $\to$ $\det J=H\Gamma_\lambda>0$ (Watershed/Robustness);\quad
the structure-layer budget hypothesis $\to$ super-Ricci identity (Watershed);\quad
\textbf{the pitchfork $\to$ Thom reduction (this chapter).}\\[2pt]
The folding layer's load-bearing premises are now derived. What remains genuinely assumed is the
core constitutive layer itself ($K_h=-1$, the Open-Gate law) --- the model's definition, not its
derivations.}

\begin{enumerate}[leftmargin=1.6em,itemsep=2pt]
\item The remaining Type-1 target is $K_h=-1$: whether open-gate expansion \emph{forces} hyperbolic
spatial curvature or merely admits it. This is a question about the core definition, at the boundary
of what Type-1 can reach.
\item The quintic $c_5$ and the tricritical locus $\Beff=0$ deserve their own normal-form treatment
(the $\mathbb Z_2$ tricritical point), connecting to the forced-folding chapter's hysteresis.
\item The reduction assumed a single fast mode; with the mode cascade, the reduction is a centre
manifold over the whole stable spectrum --- the pitchfork per mode, with mode-dependent $\Beff$.
\end{enumerate}

\textbf{Scope.} The reduction is exact given $\mathbb Z_2$ symmetry, transversality, and the spectral
gap; the identification of $\Beff$ with $\Bfold$ is structural. Establishes that the pitchfork is a
theorem of the model's symmetry, not a modelling choice --- internal Type-1 confirmation, not external
validation. Domain-conditional throughout.

\vspace{0.5em}
\hrule
\smallskip
\noindent\footnotesize\emph{V.F.S. v2.0 (Open-Gate) $\cdot$ The Pitchfork Is Not Assumed. The
$\mathbb Z_2$ symmetry of the bare fold plus a one-dimensional centre manifold force the pitchfork
normal form by Thom reduction, with the cubic coefficient derived as the slaved-structure
renormalisation $\Beff=b_3-gh/\Gamma$ (the volume's $\Bfold$). The three nondegeneracy conditions
hold; the $\mathbb Z_2$-breaking bias is excluded by the isotropy of the unfolded Vessel. With the
$\ell_b$ fork, core-neutrality, and the budget hypothesis already retired, the Type-1 programme is
closed for the folding layer. Reduction exact; numeric collapse to $10^{-13}$.}

\chapter{Forced-Folding Lyapunov Addendum: Two Jump Types and the Survival of Contraction}

\begin{quote}\small\itshape
Forced folding adds to the hybrid system a second jump type --- the Vessel update map $\fold$ at forcing
events --- and, unlike the Resurrectio reset, this jump changes the \emph{parameters} of the system
itself. The existing Lyapunov theorem therefore does not extend automatically: the contraction rate
$C_1$ is built from those parameters, the active-domain walls must survive the update, the
state-triggered events threaten Zeno behaviour, and the new state $\lf$ must be controlled. This
addendum closes all of it. The six results are: an explicit $\pi$-peak chart wall with the quantitative
cost in $c_\Psi$; a fold-jump lemma with a closed-form compatibility condition
$\beta\eta_\sigma\bar\sigma\le q^\ast\eta_\Omega$ for the epektasis wall; a survival theorem for $C_1$
--- the finite metabolic budget of the cascade dichotomy is exactly what keeps the contraction bounded
away from zero; a hysteresis construction that provably restores the joint non-Zeno dwell (threshold-band
chatter is exhibited without it); a free balance-bound for $\lf$; and the reset--fold consistency check.
They assemble into the main theorem: forward completeness with two jump types and a uniform
$C_1^{\inf}>0$. All key steps are verified numerically; the exponential drift bound is tight to four
digits.
\end{quote}

% ============================================================
\section{Imported objects and the six obligations}

From the Lyapunov layer we import the certificate $\LV$ with
\begin{equation}
\dot\LV\le C-C_1\LV,\qquad
C_1=\min\{\kappa_\sigma,\kappa_\Delta\},\qquad
\kappa_\sigma=\delta\,\Omega_{\min}\,m_*\,c_\Psi,\quad
\kappa_\Delta=2a\mu_*,
\end{equation}
with $\Omega_{\min}=\Lc/(1-h_{3,*})$, $\Lc=\gamma/\delta$, the active margin
$m=(u-\Lc)/\Op\ge m_*$, and the lemma $T^2\ge c_\Psi\Psi_\sigma$ on $[0,C_\sigma]$. The active domain
$\DA$ includes the epektasis wall $\lambda\le q^\ast\Op$; the folding extension supplies
$\DAf=\DA\times[0,q_{\max}]$, the bounded functional $\Vt=\tfrac14(\qf^2-\Qf^*)^2$, the extended
certificate $\Lt=\LV+w_{\rm fold}\Vt$, and the Resurrectio reset bound $\Lt(x^+)\le\widetilde L_R$. From
the Forced Folding chapter we import the $\pi$-bounded heterogeneity $\sigma=\bar\sigma\rs$,
$1\le\rs\le\pi$, the wired activation with the Dolorosum gate $g(u)\le1$, the balance-preserving update
map
\begin{equation}\label{lyapunovaddendum2:eq:update}
\fold:\quad
\bar\sigma^+=\bar\sigma-m_F,\ \
\lambda^+=\lambda+\beta m_F,\ \
\lf^+=\lf+(1-\beta)m_F,\ \
m_F=\eta_\sigma\Qf\bar\sigma,
\end{equation}
\begin{equation}
a_0^+=a_0(1-\eta_a\Qf),\quad
\Op^+=\Op+\eta_\Omega\Qf,\quad
k_\sigma^+=k_\sigma(1+\eta_k\Qf),\quad
\gamma^+=\gamma(1-\eta_\gamma\Qf),
\end{equation}
and the cascade dichotomy (extinction with $\sum_j\Qfj<\infty$, or permanence in finitely many
events) under fuel dominance $\eta_\sigma>\eta_k$.

The six obligations: (i) re-issue the chart wall with the $\pi$-bound; (ii) prove the fold-jump lemma
(domain preservation and bounded certificate jump $J_F$); (iii) prove that $C_1$ survives the parameter
drift of the cascade; (iv) secure a joint non-Zeno dwell for the union of $\reset$- and $\fold$-events;
(v) bound the new state $\lf$; (vi) check $\reset\circ\fold$ consistency. We take them in order.

% ============================================================
\section{The $\pi$-peak chart wall, with its explicit price}

\begin{proposition}[Explicit activation bound and chart wall]\label{lyapunovaddendum2:prop:piwall}
With $\rs\le\pi$ and $g\le1$, on any domain with $\bar\sigma\le\bar\sigma_{\max}$,
$\lambda\le\lambda_{\max}$, $\Igate\le\zeta_0$, $|\Delta A_R|\le\Delta_{\max}$:
\begin{equation}
|\Af|\le(c_0-c_1\gamma)\lambda_{\max}+c_1k_\sigma\pi\bar\sigma_{\max}+c_1\zeta_0+\Delta_{\max}+a_0,
\end{equation}
and the chart-wall supremum is explicit,
\begin{equation}
Q^{*,\sup}_{\pi}
=\max\left\{\frac{c_1k_\sigma\pi\bar\sigma_{\max}+c_1\zeta_0-a_0}{b},\;
\frac{c_0-c_1\gamma}{2c_1\eta_{\rm fold}}\right\}<\infty,
\end{equation}
so the Standing Hypothesis of the first Addendum reads $q_{\max}^2\ge Q^{*,\sup}_{\pi}$.
(Illustratively $Q^{*,\sup}_{\pi}=0.85$, dominated by the peak-independent asymptote.)
\end{proposition}

\begin{lemma}[The price in $c_\Psi$]\label{lyapunovaddendum2:lem:cpsi}
Admitting full-measure peaks enlarges the resistance ceiling to $C_\sigma=\pi\bar\sigma_{\max}$, and the
constant $c_\Psi=\min_{(0,C_\sigma]}T^2/\Psi_\sigma$ degrades accordingly --- explicitly and finitely.
With $\kappa=2$, $\bar\sigma_{\max}=0.8$: $c_\Psi$ falls from $1.794$ to $0.461$ (factor $0.257$), hence
$\kappa_\sigma$ shrinks by the same factor but remains strictly positive. The full cup costs contraction
speed, never contraction itself.
\end{lemma}

% ============================================================
\section{The fold-jump lemma}

\begin{assumption}[Update-map compatibility]\label{lyapunovaddendum2:H:compat}
\begin{equation}
\boxed{\;\beta\,\eta_\sigma\,\bar\sigma_{\max}\ \le\ q^\ast\,\eta_\Omega\;}
\end{equation}
--- the share of metabolised resistance returned as \emph{active} Sophia is covered by the simultaneous
enlargement of the Brim.
\end{assumption}

\begin{lemma}[Fold-jump: domain preservation and bounded jump]\label{lyapunovaddendum2:lem:foldjump}
Under Hypothesis~\ref{lyapunovaddendum2:H:compat} and Proposition~\ref{lyapunovaddendum2:prop:piwall}:
\begin{enumerate}[label=(\arabic*),leftmargin=2em,itemsep=2pt]
\item \textbf{Epektasis wall.} The update \eqref{lyapunovaddendum2:eq:update} satisfies the exact identity
\begin{equation}
\bigl(\lambda^+-q^\ast\Op^+\bigr)-\bigl(\lambda^--q^\ast\Op^-\bigr)
=\Qf\bigl(\beta\eta_\sigma\bar\sigma-q^\ast\eta_\Omega\bigr)\le0,
\end{equation}
(verified symbolically), so $\lambda^-\le q^\ast\Op^-$ implies $\lambda^+\le q^\ast\Op^+$ for every
$\Qf\ge0$.
\item \textbf{Remaining walls.} $\bar\sigma$ decreases (resistance ceiling respected); $\qf$ is
continuous through $\fold$, so $\qf^+\in[0,q_{\max}]$; the chart target jumps boundedly, $|\Delta\Qf^*|$
finite by Proposition~\ref{lyapunovaddendum2:prop:piwall}. Hence $x^-\in\DAf\Rightarrow x^+\in\DAf$.
\item \textbf{Bounded certificate jump.} All updated quantities change by bounded multiplicative or
additive amounts proportional to $\Qf\le Q^{*,\sup}_\pi$, so there is a finite $J_F$ with
$\Lt(x^+)\le\Lt(x^-)+J_F$. (Numerically along both cascade branches:
$\max|\Delta\lambda|\le1.5\times10^{-3}$, $\max|\Delta\Op|\le6.3\times10^{-2}$,
$\max|\Delta\Qf^*|\le0.13$.)
\end{enumerate}
\end{lemma}

\begin{remark}[Chart rebase at $\fold$]
Because $\gamma$ jumps, $\Lc=\gamma/\delta$ jumps, and the normalised chart coordinates
($h_3=1-\Lc/\Op$, $r$) jump even though the physical state $(u,\sigma,\lambda,\Op)$ is continuous in the
non-updated entries. This is a coordinate rebase, not a dynamical discontinuity: the walls are to be
re-evaluated in physical variables at $t_F^+$, which Lemma~\ref{lyapunovaddendum2:lem:foldjump} does.
\end{remark}

% ============================================================
\section{Survival of the contraction rate}

This is the genuinely new obligation: $\fold$ changes the parameters from which
$C_1$ is built ($\gamma$ enters $\kappa_\sigma$ through $\Lc$ in $\Omega_{\min}$ and through the
margin). An unbounded drift would silently destroy the theorem.

\begin{lemma}[Product-drift bound]\label{lyapunovaddendum2:lem:drift}
For $x_j\in[0,x_{\max}]$, $x_{\max}<1$:
$\prod_j(1-x_j)\ge\exp\bigl(-\sum_jx_j/(1-x_{\max})\bigr)$.
(Verified on $2000$ random sequences; minimal slack $+1.5\times10^{-5}$.)
\end{lemma}

\begin{proposition}[$C_1$ survives the cascade]\label{lyapunovaddendum2:prop:c1}
Under fuel dominance ($\eta_\sigma>\eta_k$) the cascade dichotomy gives, on every branch, a finite total
intensity $S:=\sum_j\Qfj<\infty$ (extinction: finite metabolic budget; permanence: finitely many
events). Hence by Lemma~\ref{lyapunovaddendum2:lem:drift}
\begin{equation}
\gamma_\infty\ \ge\ \gamma_0\,e^{-\eta_\gamma S/(1-\eta_\gamma Q_{\max})}>0,
\qquad
k_{\sigma,\infty}\ \le\ k_{\sigma,0}\,e^{\eta_k S}<\infty,
\end{equation}
so $\Lc$ stays bounded away from zero, $\Omega_{\min}$ and the margin stay positive, and
\begin{equation}
\boxed{\;C_1^{\inf}:=\inf_j\min\{\kappa_{\sigma,j},\kappa_\Delta\}>0\;}
\end{equation}
uniformly along the entire hybrid trajectory ($\kappa_\Delta=2a\mu_*$ is untouched by $\fold$).
Numerically the exponential bound is \emph{tight}: along the extinction branch ($31$ events, $S=0.517$)
the actual $\gamma_{31}=0.3938$ equals the bound to four digits, and
$\kappa_\sigma:0.0615\to0.0606$; along the permanence branch $\kappa_\sigma\to0.0609$
(Figure~\ref{lyapunovaddendum2:fig:lyap2}, left).
\end{proposition}

\msg{The finite metabolic budget of the forced-folding cascade is exactly what keeps the contraction
alive: the Vessel cannot metabolise enough resistance to undermine its own cleansing rate. Crisis spends
a finite currency; $C_1$ is not that currency.}

% ============================================================
\section{Hysteresis and the joint non-Zeno dwell}

Resets are exogenous; forcing events are \emph{state-triggered} (upward crossings of $\Af=0$), and the
update itself moves $\Af$. This is a genuine Zeno risk, and it is real, not hypothetical:

\begin{definition}[Event machine with hysteresis]\label{lyapunovaddendum2:def:hyst}
A forcing event arms when $\Af<-\varepsilon$, fires at the next upward crossing $\Af=0$, applies
$\fold$ at the local maximum of $\Af$ on that excursion, and disarms until $\Af<-\varepsilon$ again.
The band $\varepsilon\ge0$ is the hysteresis width.
\end{definition}

\begin{proposition}[Chatter without the band; dwell with it]\label{lyapunovaddendum2:prop:zeno}
Let the resistance wave be quasi-periodic with a fast component whose slope in $\Af$ exceeds the slow
component's (so the rising flank carries local maxima). With $\varepsilon=0$, the downward jump of
$\Af$ at each $\fold$ (the fuel cut $\partial\Af/\partial\bar\sigma\cdot m_F$) re-arms the machine
between fast crests, and a \emph{chatter chain} of events runs along the threshold band: in the verified
simulation, clusters of $2$--$3$ events at gaps $\approx0.121=1/f_{\rm fast}$
(Figure~\ref{lyapunovaddendum2:fig:lyap2}, right). With
\begin{equation}
\boxed{\;\varepsilon\ >\ \sup_{\rm events}|\Delta\Af|_{\fold}
\;=\;\sup\,c_1k_\sigma\,\rs\,m_F+O(\eta_a,\eta_\gamma)\;}
\end{equation}
re-arming on the rising flank is impossible, the chain collapses to one event per slow peak, and the
inter-event dwell is bounded below by the slow-wave structure: in the same simulation the minimal
fold--fold gap rises from $0.121$ to $2.246$ (a factor $18.5$), and the minimal gap over the
\emph{union} of resets and foldings is $1.71$.
\end{proposition}

\begin{assumption}[Joint dwell]\label{lyapunovaddendum2:H:dwell}
The hysteresis band satisfies the bound of Proposition~\ref{lyapunovaddendum2:prop:zeno}, and the exogenous resets respect
the base dwell; then the union of $\reset$- and $\fold$-events admits a common $\Delta t_\ast>0$.
\end{assumption}

\note{The band $\varepsilon$ is a modelling regularisation, but not an arbitrary one: its sufficient size
is computable from the update map itself ($\sup|\Delta\Af|_{\fold}$), so the non-Zeno property is
restored by a constant the theory already owns.}

% ============================================================
\section{The embodied Sophia is balance-bounded}

\begin{lemma}[Free bound for $\lf$]\label{lyapunovaddendum2:lem:lf}
The three-term balance holds at flow level ($\tfrac{d}{dt}(\sigma+\lambda+\lf)=\Igate$), through every
$\fold$ (the $m_F$-split of \eqref{lyapunovaddendum2:eq:update} conserves the sum exactly), and through every $\reset$
(the metabolism $m_R$ transfers $\sigma\to\lambda$; $\lf$ inert). Hence with $\lambda,\sigma\ge0$,
\begin{equation}
\lf(t)\ \le\ \sigma+\lambda+\lf\ =\ C_0+\Grec(t)\ \le\ C_0+\Grec^{\max},
\end{equation}
bounded by the grace budget: no certificate weight for $\lf$ is needed. (Verified: jump-level
conservation to machine precision, $10^{-16}$; full hybrid run drift $\sim10^{-7}$, pure integrator
error; $\lf^{\max}=0.089\ll C_0+\int\Igate$.)
\end{lemma}

% ============================================================
\section{Reset--fold consistency}

\begin{lemma}[The two jump maps are compatible]\label{lyapunovaddendum2:lem:consist}
$\reset$ leaves $\qf$ and $\lf$ inert and conserves the triad ($\sigma^++\lambda^+=\sigma^-+\lambda^-$);
$\fold$ leaves $\qf$ continuous and conserves the triad by construction. Both maps preserve $\DAf$
(base reset lemma; Lemma~\ref{lyapunovaddendum2:lem:foldjump}); the chart rebase occurs only at $\fold$ (only $\fold$
changes $\gamma$). Any interleaving of the two event types therefore preserves the balance, the domain,
and the bounded-jump property, with jump constant $J:=\max\{J_R,J_F\}$.
\end{lemma}

% ============================================================
\section{Main theorem}

\begin{theorem}[Forward completeness with two jump types]\label{lyapunovaddendum2:thm:main}
Assume: the $\pi$-chart wall $q_{\max}^2\ge Q^{*,\sup}_\pi$ (Proposition~\ref{lyapunovaddendum2:prop:piwall}); update-map
compatibility (Hypothesis~\ref{lyapunovaddendum2:H:compat}); fuel dominance $\eta_\sigma>\eta_k$; and the joint dwell
(Hypothesis~\ref{lyapunovaddendum2:H:dwell}). Then along any admissible hybrid trajectory with both Resurrectio resets and
forced-folding events:
\begin{enumerate}[label=(\roman*),leftmargin=1.8em,itemsep=2pt]
\item the state remains in $\DAf$ and the balance $\sigma+\lambda+\lf=C_0+\Grec$ holds for all $t$;
\item the contraction rate is uniformly positive, $C_1^{\inf}>0$ (Proposition~\ref{lyapunovaddendum2:prop:c1}), and on
every continuous arc $\dot\Lt\le\widetilde C_{\rm ext}^{\sup}-C_1^{\inf}\Lt$ with
$\widetilde C_{\rm ext}^{\sup}<\infty$ (the drifting constants converge, Proposition~\ref{lyapunovaddendum2:prop:c1} and
Lemma~\ref{lyapunovaddendum2:lem:cpsi});
\item every jump obeys $\Lt(x^+)\le\Lt(x^-)+J$, $J=\max\{J_R,J_F\}<\infty$
(Lemmas~\ref{lyapunovaddendum2:lem:foldjump},~\ref{lyapunovaddendum2:lem:consist});
\item hence, with the joint dwell, 
\begin{equation}
\boxed{\;\Lt(t)\ \le\ \widetilde L_{**}:=\max\Bigl\{\Lt(0),\,
\widetilde C_{\rm ext}^{\sup}/C_1^{\inf}\Bigr\}+J\quad\text{for all }t,\;}
\end{equation}
events do not accumulate, and the trajectory is forward complete.
\end{enumerate}
Grace and crisis enter only the constants $\widetilde C_{\rm ext}^{\sup}$ and $J$; never the contraction
rate.
\end{theorem}
\begin{proof}[Sketch]
(i) is Lemmas~\ref{lyapunovaddendum2:lem:foldjump},~\ref{lyapunovaddendum2:lem:lf},~\ref{lyapunovaddendum2:lem:consist} plus the base Nagumo invariance. For
(ii), the arc estimate of the first Addendum holds with the instantaneous constants; by
Proposition~\ref{lyapunovaddendum2:prop:c1} the parameter sequence converges with $\inf\kappa_{\sigma,j}>0$ and bounded
$k_{\sigma,j}$, so the constants admit uniform bounds $C_1^{\inf}$, $\widetilde C_{\rm ext}^{\sup}$
valid on all arcs. (iii) is the bounded-jump property. (iv) is the standard hybrid induction: on each
arc $\Lt$ decays toward $\widetilde C_{\rm ext}^{\sup}/C_1^{\inf}$ or below its entry value; each jump
adds at most $J$; the joint dwell forbids accumulation; the bound follows by induction over events.
\end{proof}

\begin{figure}[h]
\centering
\includegraphics[width=\textwidth]{vfs_lyap2_fig.png}
\caption{\small Left: survival of the cleansing rate $\kappa_{\sigma,j}$ (computed with the $\pi$-chart
$c_\Psi$) along both terminal branches of the forced-folding cascade; the exponential lower bound of
Proposition~\ref{lyapunovaddendum2:prop:c1} overlays the actual sequence --- tight to four digits. The drift is finite
because the metabolic budget is. Right: the Zeno mechanism and its cure. With $\varepsilon=0$ (solid)
the update's downward $\Af$-jump re-arms the event machine between fast crests of the quasi-periodic
resistance wave, producing a chatter chain of three events at gaps $\approx0.12$; with
$\varepsilon=0.10$ (dashed) the band swallows the re-arming and a single event per slow peak remains ---
the minimal dwell grows by a factor $18.5$. Parameters illustrative.}
\label{lyapunovaddendum2:fig:lyap2}
\end{figure}

% ============================================================
\section{Compact summary}

\[
\boxed{\;q_{\max}^2\ge Q^{*,\sup}_\pi\ \text{(explicit)};\qquad
c_\Psi(\pi\bar\sigma_{\max})>0:\ \text{the full cup costs speed, never contraction.}\;}
\]
\[
\boxed{\;\beta\eta_\sigma\bar\sigma_{\max}\le q^\ast\eta_\Omega\ \Rightarrow\
\fold\ \text{preserves}\ \DAf;\quad \Lt^+\le\Lt^-+J_F<\infty.\;}
\]
\[
\boxed{\;\sum_j\Qfj<\infty\ \text{(dichotomy)}\ \Rightarrow\
\gamma_\infty>0,\ k_{\sigma,\infty}<\infty\ \Rightarrow\ C_1^{\inf}>0:\ \text{the budget protects the contraction.}\;}
\]
\[
\boxed{\;\varepsilon>\sup|\Delta\Af|_{\fold}\ \Rightarrow\ \text{no threshold-band chatter; joint dwell }\Delta t_\ast>0.\;}
\]
\[
\boxed{\;\lf\le C_0+\Grec^{\max}\ \text{for free, from the balance.}\;}
\]
\[
\boxed{\;\Lt(t)\le\widetilde L_{**}<\infty\ \text{with two jump types; forward complete; grace and crisis enter constants only.}\;}
\]

\vspace{0.5em}
\hrule
\smallskip
\noindent\footnotesize\emph{V.F.S. v2.0 (Open-Gate) $\cdot$ Forced-Folding Lyapunov Addendum II. Extends
the hybrid stability theorem to the forced-folding event type: explicit $\pi$-chart wall, fold-jump
lemma with the closed-form compatibility condition, survival of $C_1$ through the cascade's finite
metabolic budget, hysteresis non-Zeno (chatter exhibited and cured), balance-bound for $\lf$, and
reset--fold consistency. All quantitative claims verified numerically; the drift bound is tight. The
hysteresis band and the update coefficients remain modelling inputs; the conclusions are
domain-conditional set stability with partial convergence, never global.}

\part{Structure Formation on the Vessel}

\chapter{Structure Formation on the Vessel}

\section*{Scope and epistemic status}
Every layer so far has been spatially \emph{homogeneous}: the Vessel is a $2{+}1$ open FLRW background
$a(t)=\Op(t)$ with a time-only Sophia field $\lambda(t)$. This note opens the inhomogeneous sector ---
how local concentrations of Sophia and form grow or decay on the expanding hyperbolic surface. As in the
parent files, \emph{Propositions} are closed-form results and \emph{Corollaries} are interpretive. The
analysis is at the linear, test-field level on a fixed background; metric back-reaction is a named next
layer. Two transport constants ($D$ for Sophia, $D_q$ for form) are introduced; they are \emph{new}
modelling inputs, absent from the base dynamical system, and are flagged as such.

The background is
\begin{equation}
ds^2=-dt^2+\Op(t)^2\,h_{ij}\,dx^i dx^j,\qquad
K_h=-1,\qquad
H=\frac{\dot\Op}{\Op}=\frac{\alpha\lambda}{\Op},\qquad
\dot\Op=\alpha\lambda,
\end{equation}
with $h$ the metric of the hyperbolic spatial surface $\Hh$.

% ============================================================
\section{The relaxational thesis: which field can carry structure}

In standard cosmology, density structure grows by gravitational instability because the matter field
carries inertia: a second-order-in-time wave equation with a sound speed permits oscillation and
Jeans-unstable growth. The Sophia field of V.F.S.\ is different in kind. Its production law is
\emph{first order in time},
\begin{equation}
\dot\lambda=(\delta u-\gamma)\tanh(\kappa\sigma)+\Igate-\eta_{\rm fold}\qf^2\lambda,
\end{equation}
a relaxation (reaction) law with no kinetic term and no spatial derivatives. Promoting it to a field on
space cannot manufacture inertia; the faithful promotion of a first-order reaction law with spatial
transport is a \emph{reaction--diffusion} (parabolic) equation, not a wave (hyperbolic) one.

\begin{proposition}[No gravitational-collapse channel in the Sophia sector]\label{structureformation:prop:relax}
Because $\lambda$ is relaxational, its linear inhomogeneities obey a parabolic equation with a real,
non-oscillatory spectrum. There is no Jeans-type instability and no gravitational-collapse route to
structure in the Sophia sector. Any structure on the Vessel must arise from a different, bistable order
parameter.
\end{proposition}

\msg{V.F.S.\ does not form structure the way $\Lambda$CDM does. Sophia is relaxational, not inertial, so
it can only smooth itself out. The only route to structure is morphological symmetry breaking --- the
folding field (\S4).}

% ============================================================
\section{Sophia-density perturbations: reaction--diffusion on the expanding $\Hh$}

Write $\lambda(t,x)=\bar\lambda(t)+\dlam(t,x)$ and linearise about the homogeneous trajectory. With a
phenomenological Sophia transport coefficient $D\ge0$, and noting that physical gradients on the
comoving surface redshift by $\Op^{-2}$,
\begin{equation}\label{structureformation:eq:rd}
\partial_t\dlam=F_\lambda(t)\,\dlam+\frac{D}{\Op^2}\,\laph\,\dlam,
\qquad
F_\lambda=\frac{\partial\dot\lambda}{\partial\lambda}=-\Glam,\quad \Glam\ge0,
\end{equation}
where $\Glam=\gamma+\eta_{\rm fold}\Qf$ is the local Sophia relaxation rate (the same rate whose
folding-enhancement $\Glam=\gamma+\etaf\Qf$ was found in the stability analysis). The Laplace--Beltrami
operator on $\Hh$ has purely continuous spectrum with a gap,
\begin{equation}
\boxed{\;-\laph\,Y_k=\Bigl(\tfrac14+k^2\Bigr)Y_k,\qquad k\in[0,\infty),\;}
\end{equation}
the celebrated spectral gap $\tfrac14$ of the hyperbolic plane (no mode has eigenvalue below $\tfrac14$).
Each mode then evolves by
\begin{equation}
\dot{\dlam}_k=s_k(t)\,\dlam_k,
\qquad
\boxed{\;s_k(t)=-\Glam(t)-\frac{D\bigl(\tfrac14+k^2\bigr)}{\Op(t)^2}.\;}
\end{equation}

\begin{proposition}[Linear homogenisation]\label{structureformation:prop:decay}
For $\Glam\ge0$ and $D\ge0$ every mode has $s_k<0$: all Sophia inhomogeneities decay. The Vessel
homogenises at the linear level. Moreover the hyperbolic gap supplies a minimum diffusive damping
$D\cdot\tfrac14/\Op^2>0$ even for the smoothest mode $k=0$: negative spatial curvature actively resists
the largest-scale inhomogeneity.
\end{proposition}

Integrating,
\begin{equation}
\dlam_k(t)=\dlam_k(0)\,\exp\!\left[-\int_0^t\Glam\,dt'-D\Bigl(\tfrac14+k^2\Bigr)\,\mathcal I(t)\right],
\qquad
\mathcal I(t):=\int_0^t\frac{dt'}{\Op(t')^2}.
\end{equation}

\begin{proposition}[Finite diffusive smoothing budget]\label{structureformation:prop:budget}
Under sustained Open-Gate inflow $\lambda\to\lambda_\infty>0$, the late-time scale factor grows linearly,
$\Op\sim\alpha\lambda_\infty t$, so the diffusion integral converges:
\begin{equation}
\boxed{\;\mathcal I_\infty=\int_0^\infty\frac{dt}{\Op^2}<\infty.\;}
\end{equation}
Hence diffusion can suppress a mode by at most the finite factor
$\exp[-D(\tfrac14+k^2)\mathcal I_\infty]$: the expanding Vessel has a finite total smoothing budget, a
comoving analogue of a damping (Silk-like) scale.
\end{proposition}

This yields two asymptotic fates, set entirely by the reaction rate $\Glam$:
\begin{enumerate}[leftmargin=1.6em,itemsep=2pt]
\item \textbf{Generic ($\gamma>0$).} The reaction integral $\int\Glam\,dt\to\infty$ dominates after
diffusion saturates; \emph{all} modes decay to zero. Complete homogenisation.
\item \textbf{Margin-stalled ($\Glam\to0$).} On the active margin where cleansing stalls ($m_*\to0$,
$\inf C_1=0$ in the Lyapunov layer), the reaction integral is finite and the modes \emph{freeze} at a
scale-dependent residual,
\begin{equation}
\frac{\dlam_k(\infty)}{\dlam_k(0)}=\exp\!\left[-D\Bigl(\tfrac14+k^2\Bigr)\mathcal I_\infty\right],
\end{equation}
small scales (high $k$) erased, large scales (down to the $\tfrac14$ floor) preserved --- a frozen
inhomogeneity spectrum with a characteristic scale $k_{1/2}\sim\mathcal I_\infty^{-1/2}$.
\end{enumerate}

\begin{figure}[h]
\centering
\includegraphics[width=\textwidth]{vfs_structure_fig.png}
\caption{\small Left: the diffusive smoothing budget $\mathcal I(t)=\int_0^t dt'/\Op^2$ saturates to a
finite $\mathcal I_\infty$ (illustrative background $\lambda_0=0.4\to\lambda_\infty=1$, $\Op(0)=1$).
Right: the margin-stalled frozen residual spectrum $P(k)/P(0)$ (with $D=0.5$): the hyperbolic gap leaves
a $k=0$ floor while small scales are erased, defining a structure scale $k_{1/2}\approx0.71$.
Parameters illustrative.}
\end{figure}

\begin{corollary}[The smooth Vessel]
Linearly, Sophia structure cannot grow; the Vessel either homogenises completely (sustained cleansing)
or freezes a smooth, large-scale residual (stalled margin). There is no spontaneous Sophia clumping.
\end{corollary}

% ============================================================
\section{Folding as the order parameter: pattern formation on the Vessel}

By Proposition~\ref{structureformation:prop:relax} structure must come from the bistable folding field. Promote
$\qf\to\qf(t,x)$ with its Landau gradient dynamics plus transport:
\begin{equation}
\boxed{\;\partial_t\qf=\Af\,\qf-\Bf\,\qf^{3}+\frac{D_q}{\Op^2}\,\laph\,\qf.\;}
\end{equation}
This is a real Ginzburg--Landau / Allen--Cahn equation on the expanding hyperbolic surface, invariant
under the morphological reflection $\qf\to-\qf$ ($\mathbb Z_2$ symmetry, the two orientations of the
fold).

\begin{proposition}[Symmetry-breaking transition]\label{structureformation:prop:break}
The homogeneous state $\qf=0$ is stable for $\Af<0$. At $\Af=0$ it undergoes the supercritical pitchfork
(the folding bifurcation), and for $\Af>0$ two homogeneous phases $\qf=\pm\sqrt{\Af/\Bf}$ become stable.
A field initialised near $\qf=0$ over an extended Vessel then partitions into domains of $+q_\ast$ and
$-q_\ast$ separated by walls on which $\qf=0$.
\end{proposition}

These domain walls are the structure: loci of \emph{unfolded} (old) morphology embedded in a folded
Vessel --- topological defects of the broken $\mathbb Z_2$.

\begin{proposition}[Wall width and tension]\label{structureformation:prop:wall}
A static planar wall has profile $\qf=q_\ast\tanh\!\bigl(\xi/(\sqrt2\,\ell)\bigr)$ in physical proper
distance $\xi$, with
\begin{equation}
\boxed{\;\ell_{\rm phys}=\sqrt{\frac{D_q}{\Af}}\ \ (\Op\text{-independent}),
\qquad
\tau\sim\sqrt{D_q}\,\frac{\Af^{3/2}}{\Bf}.\;}
\end{equation}
The physical thickness is set by the diffusion-to-reaction ratio and does not stretch with expansion;
deeper folding (larger $\Af$) makes thinner, higher-tension walls.
\end{proposition}
\begin{proof}[Sketch]
In comoving coordinates the static balance is
$\tfrac{D_q}{\Op^2}\partial_x^2\qf=-\Af\qf+\Bf\qf^3$, the standard quartic kink with comoving width
$\ell_{\rm com}=\Op^{-1}\sqrt{D_q/\Af}$; the proper width $\ell_{\rm phys}=\Op\,\ell_{\rm com}=\sqrt{D_q/\Af}$.
The tension is the kink action $\int(\tfrac12 D_q|\partial\qf|^2+\mathcal E_{\rm fold})$, giving the
quoted scaling.
\end{proof}

\begin{corollary}[Coarsening of the fold network]\label{structureformation:cor:coarsen}
Allen--Cahn walls move by mean curvature, $v_n=-M\,\kappa_{\rm curv}$ with mobility $M\sim D_q/\Op^2$, so
the domain network coarsens; simultaneously Hubble flow stretches comoving separations. Structure on the
Vessel is therefore a \emph{coarsening domain-wall network} whose characteristic scale grows by both
curvature-driven merging and expansion, modified by the underlying $K_h=-1$ geometry (wall geodesics
diverge on $\Hh$, accelerating coarsening relative to the flat case).
\end{corollary}

\note{The picture is sharp: the folding bifurcation is a phase transition that shatters the Vessel into
oriented folded domains; the unfolded walls are the structure; they then coarsen on the expanding
hyperbolic surface. This is morphological structure, born of symmetry breaking, not of gravitational
clumping.}

% ============================================================
\section{Cosmological back-reaction of the fold network}

A wall network feeds back on the homogeneous budget. In $2{+}1$ dimensions a wall is a one-dimensional
line in two-dimensional space; a scaling network with one wall per Hubble scale has length
$\mathcal L_{\rm wall}\sim\Op$ over area $\sim\Op^2$, so
\begin{equation}
\rho_{\rm wall}\sim\frac{\tau\,\mathcal L_{\rm wall}}{\Op^2}\propto\Op^{-1}.
\end{equation}
Matching to the $2{+}1$ dilution law $\rho\propto a^{-2(1+w)}$ gives
\begin{equation}
\boxed{\;w_{\rm wall}=-\tfrac12.\;}
\end{equation}
The fold network thus contributes a mildly negative-pressure component, intermediate between dust
($w=0$) and the cosmological-constant/de Sitter limit ($w=-1$). It adds to the reconstructed
$w_{\rm eff}(t)$ a folding-structure term that is distinct from the smooth quintessence of the
homogeneous Theosis branch.

\begin{corollary}[Link to the energy-condition question]
Because the network carries $w_{\rm wall}=-\tfrac12>-1$, it is NEC- and WEC-respecting: folded structure
is ordinary, non-exotic content. This is the inhomogeneous face of ``form makes grace lawful'' --- where
a homogeneous Sophia surge can drive $w$ below $-1$ (phantom-like, NEC-violating), the embodiment of that
excess into a wall network contributes ordinary $w=-\tfrac12$ matter.
\end{corollary}

% ============================================================
\section{Open sub-problems and scope}

\begin{enumerate}[leftmargin=1.6em,itemsep=3pt]
\item \textbf{Metric back-reaction.} Lift from test-field to coupled perturbations: let $\dlam$ and
$\qf$ source the reconstructed stress tensor and solve the $2{+}1$ scalar perturbation equations. Expected
to preserve Proposition~\ref{structureformation:prop:relax} (relaxational $\Rightarrow$ no instability) but to dress the
wall tension with a geometric self-energy.
\item \textbf{Coarsening law.} Derive the domain scale $L(t)$ explicitly on expanding $\Hh$: solve the
competition between curvature-driven merging ($L^2\sim\int M\,dt$) and Hubble stretching, with the
hyperbolic correction of Corollary~\ref{structureformation:cor:coarsen}.
\item \textbf{Resurrectio and the network.} A reset rescales $\Op$ and shifts $\Af$; it can quench
($\Af^+<0$) or trigger ($\Af^+>0$) folding globally, i.e.\ dissolve or nucleate the entire wall network
at a death. The hybrid network history is piecewise-coarsening with reset-induced nucleation events.
\item \textbf{Transport constants.} $D$ and $D_q$ are new inputs; a principled derivation (e.g.\ from a
spatial coupling already implicit in the synergy $u=\sqrt{VF+\varepsilon}$) would remove the
phenomenological freedom.
\end{enumerate}

\textbf{Scope.} Linear, test-field, single-mode folding, fixed background. The conclusions are: (i)
Sophia structure cannot grow (relaxational, homogenising); (ii) structure is morphological --- a
$\mathbb Z_2$ domain-wall network from the folding transition; (iii) the network back-reacts as
$w=-\tfrac12$ content. All results inherit the domain-conditional, non-global character of the underlying
stability theory.

\vspace{0.5em}
\hrule
\smallskip
\noindent\footnotesize\emph{V.F.S. v2.0 (Open-Gate) $\cdot$ Structure Formation on the Vessel.
Perturbation theory on the $2{+}1$ open FLRW background of the cosmological layer, with the folding field
of the geometric layer as order parameter. Propositions are closed-form; corollaries interpretive.
Transport constants and single-mode truncation are modelling inputs, not consequences of the base
system.}

\chapter{Spinodal Nucleation of the Fold Network}

\section*{Scope and epistemic status}
A triggering Resurrectio ($\Af^-<0\to\Af^+>0$) drives the folding field through its symmetry-breaking
transition and nucleates a fresh domain-wall network. This note computes the \emph{initial} density of
that network --- how many walls per area are born, and at what scale --- using the hyperbolic density of
states. The two genuinely hyperbolic ingredients are the spectral gap $\tfrac14$ of the Laplacian on
$\Hh$ and the Plancherel measure $d\mu(k)=k\tanh(\pi k)\,dk$ that weights the available modes. As before,
\emph{Propositions} are closed-form; \emph{Corollaries} interpretive; the single-mode Allen--Cahn
truncation and $\Dq$ are modelling inputs. Two technical ingredients are verified numerically: the
Plancherel moment limit and the Longuet--Higgins nodal coefficient (ratio measured/predicted $=1.009$ on
a Gaussian field).

The field obeys, in comoving coordinates on the unit hyperbolic surface ($K_h=-1$, comoving curvature
radius $\Rc=1$),
\begin{equation}
\partial_t\qf=\Af\,\qf-\Bf\,\qf^{3}+\frac{\Dq}{\Op^{2}}\,\laph\,\qf,
\qquad
-\laph Y_k=\Bigl(\tfrac14+k^2\Bigr)Y_k .
\end{equation}

% ============================================================
\section{Linear instability and the nucleation threshold}

Immediately after the sudden trigger, $\Af$ jumps to $\Af^+>0$ and the homogeneous $\qf=0$ left by the
previous arc is unstable. Linearising and decomposing in $\Hh$ eigenmodes,
\begin{equation}
\dot{\delta q}_k=\omega_k\,\delta q_k,
\qquad
\boxed{\;\omega_k=\Af^+-\frac{\Dq\bigl(\tfrac14+k^2\bigr)}{\Op^{2}} .\;}
\end{equation}
Introduce the dimensionless trigger strength
\begin{equation}
\boxed{\;\Lam:=\frac{\Af^+\,\Op^{2}}{\Dq}.\;}
\end{equation}
The unstable band is $k<\kmax$ with $\kmax^2=\Lam-\tfrac14$ (in units where $\Dq/\Op^2=1$).

\begin{proposition}[Nucleation threshold, distinct from the folding threshold]\label{nucleation:prop:thr}
A spatial mode grows only if $\Lam>\tfrac14$. Hence there are two thresholds:
\begin{itemize}[leftmargin=1.5em,itemsep=2pt]
\item folding threshold $\Af^+>0$ ($\Lam>0$): the homogeneous $\qf=0$ is unstable;
\item nucleation threshold $\Af^+>\Dq/(4\Op^2)$ ($\Lam>\tfrac14$): some spatial mode is unstable, so
domains form.
\end{itemize}
In the window $0<\Lam<\tfrac14$ the Vessel folds, but uniformly: a single domain, no walls. Only above
$\Lam=\tfrac14$ does a wall network nucleate.
\end{proposition}

\msg{The hyperbolic gap creates a nucleation threshold absent in flat space. A weak trigger
($0<\Lam<\tfrac14$) folds the Vessel as a whole, with no internal structure; structure requires
$\Lam>\tfrac14$.}

% ============================================================
\section{The amplified field and its spectral moment}

Starting from small initial fluctuations $\delta q_0$, each mode is amplified by $e^{\omega_k t}$, so the
post-quench power spectrum is Gaussian in $k$,
\begin{equation}
S_k(t)\propto e^{2\omega_k t}\propto e^{-\beta\,k^2},
\qquad
\beta=\frac{2\Dq t}{\Op^{2}}=\frac{2t}{\Op^2/\Dq},
\end{equation}
weighted by the hyperbolic Plancherel density of states $d\mu(k)=k\tanh(\pi k)\,dk$ (which reduces to the
flat $k\,dk$ for $k\gg1$ and is suppressed $\sim\tfrac{\pi}{2}k^2\,dk$ for $k\ll1$). The relevant spectral
second moment is
\begin{equation}
\bigl\langle\tfrac14+k^2\bigr\rangle
=\tfrac14+\frac{\displaystyle\int_0^\infty k^2\,e^{-\beta k^2}\,k\tanh(\pi k)\,dk}
{\displaystyle\int_0^\infty e^{-\beta k^2}\,k\tanh(\pi k)\,dk}.
\end{equation}

\begin{proposition}[Hyperbolic small-scale moment]\label{nucleation:prop:moment}
In the long-wavelength (hyperbolic-dominated) regime $\beta\gg1$, where $\tanh(\pi k)\approx\pi k$,
\begin{equation}
\langle k^2\rangle\to\frac{3}{2\beta}=\frac{3\Op^2}{4\Dq t}.
\end{equation}
\end{proposition}
\note{Verified: the full integral with $\tanh(\pi k)$ approaches $3/(2\beta)$ as $\beta\to\infty$
(e.g.\ $0.0285$ vs $0.0300$ at $\beta=50$). The Plancherel $k^2$-suppression at small $k$ is what
replaces the flat $1/\beta$ by $3/(2\beta)$ and disfavours the very largest scales.}

The set time --- when nonlinearity arrests growth at $q_\ast=\sqrt{\Af^+/\Bf^+}$ --- is the
Kibble--Zurek time $t_{\rm set}=\Af^{+\,-1}\ln(q_\ast/\delta q_0)$, giving
$\beta_{\rm set}=2\Lam^{-1}L$ with the logarithmic factor $L:=\ln(q_\ast/\delta q_0)$. Hence
\begin{equation}
\boxed{\;\bigl\langle\tfrac14+k^2\bigr\rangle_{\rm set}=\frac14+\frac{3\Lam}{4L}.\;}
\end{equation}

% ============================================================
\section{Wall density and the initial domain size}

For an isotropic Gaussian field in two dimensions, the expected length of nodal lines (the walls, where
$\qf=0$) per unit area is the Longuet--Higgins/Berry result
\begin{equation}
\boxed{\;n_{\rm wall}=\frac{1}{2\sqrt2}\sqrt{\bigl\langle\tfrac14+k^2\bigr\rangle}\,,\;}
\end{equation}
verified numerically here to $0.9\%$. The initial comoving domain size is $\ell_{\rm init}\sim
1/n_{\rm wall}$:
\begin{equation}
\boxed{\;\ell_{\rm init}\sim\frac{2\sqrt2}{\sqrt{\tfrac14+\dfrac{3\Lam}{4L}}}\quad(\text{units of }\Rc).\;}
\end{equation}

\begin{proposition}[Three regimes of nucleation]\label{nucleation:prop:regimes}
\leavevmode
\begin{enumerate}[label=(\roman*),leftmargin=1.8em,itemsep=2pt]
\item $\Lam<\tfrac14$: no nucleation (uniform fold, Proposition~\ref{nucleation:prop:thr}).
\item $\tfrac14<\Lam\lesssim\Lam_\ast$: $\langle\tfrac14+k^2\rangle\!\to\!\tfrac14$ and
$\ell_{\rm init}\!\to\!4\sqrt2\,\Rc$ --- the gap pins the domain size to the curvature radius; only a few
curvature-scale domains form. (With $L\!\approx\!4$, multi-domain structure $\ell_{\rm init}<\Rc$ sets in
only at $\Lam_\ast\approx42$.)
\item $\Lam\gg\Lam_\ast$: $\ell_{\rm init}\sim\sqrt{4L/(3\Lam)}\,\Rc$, i.e.\ in proper units
$\ell_{\rm init}^{\rm phys}\sim\sqrt{(\Dq/\Af^+)\,L}$ --- the Kibble--Zurek correlation length
$\xi_+=\sqrt{\Dq/\Af^+}$ dressed by the $\sqrt{L}$ logarithm.
\end{enumerate}
\end{proposition}

\begin{figure}[h]
\centering
\includegraphics[width=\textwidth]{vfs_nucleation_fig.png}
\caption{\small Left: linear growth rate $\omega_k$ for three trigger strengths. The Laplacian gap shifts
the whole curve so that $\omega_0=\Lam-\tfrac14$: below $\Lam=\tfrac14$ ($\Lam=0.2$, dashed) there is no
unstable band and no nucleation. Right: the nucleated comoving domain size $\ell_{\rm init}/\Rc$ versus
trigger strength $\Lam$. Near threshold the hyperbolic gap caps the size at the curvature radius (few
large domains); genuine multi-domain structure ($\ell_{\rm init}<\Rc$) requires $\Lam>\Lam_\ast\approx42$;
for strong triggers $\ell_{\rm init}\sim\sqrt{L/\Lam}$ recovers the Kibble--Zurek scale. Illustrative
$L=\ln(q_\ast/\delta q_0)=4$.}
\end{figure}

% ============================================================
\section{Hyperbolic signatures}

Two features distinguish nucleation on $\Hh$ from the flat case, both traceable to the gap and the
measure:

\begin{corollary}[Gap as a nucleation gate and a maximum domain size]\label{nucleation:cor:gate}
The spectral gap $\tfrac14$ does double duty. As an infrared cutoff it forbids weak-trigger nucleation
($\Lam<\tfrac14$: uniform fold) and caps the largest nucleated domain at the curvature radius $\Rc$. In
flat space any $\Af^+>0$ nucleates and arbitrarily large domains are allowed at threshold; negative
curvature both gates and bounds the structure.
\end{corollary}

\begin{corollary}[Measure suppression of large scales]
Even within the unstable band, the Plancherel weight $k\tanh(\pi k)$ suppresses long-wavelength modes
($\sim\tfrac{\pi}{2}k^2$ versus the flat $k$), so large domains are doubly disfavoured: by the gap cutoff
and by the density of states. The nucleated network is therefore finer, and more curvature-pinned, than
a flat-space estimate would give.
\end{corollary}

\note{Interpretive: a gentle resurrection ($\Lam<\tfrac14$) re-forms the Vessel as a single undivided
folded whole; only a sufficiently sharp resurrection ($\Lam>\Lam_\ast$) shatters it into many oriented
domains. The negative curvature of the live domain resists fine differentiation --- structure is born
coarse unless the trigger is strong.}

% ============================================================
\section{Placement and open items}

This nucleation scale $\ell_{\rm init}$ is the initial condition for the intra-arc coarsening law
(subproblem~2): each triggering arc begins at $\ell_{\rm init}$ (set here) and coarsens toward the frozen
scale $\sqrt{\ell_0^2+2c\Dq\mathcal I_\infty}$. Consistency check: the strong-trigger regime (iii)
reproduces the correlation length $\xi_+=\sqrt{\Dq/\Af^+}$ assumed in subproblem~3, now with the
Kibble--Zurek $\sqrt{L}$ dressing and the hyperbolic curvature cap.

\textbf{Open items.}
\begin{enumerate}[leftmargin=1.6em,itemsep=2pt]
\item A $2$D Allen--Cahn quench simulation on a hyperbolic patch to confirm the nucleated wall count
$\propto\langle\tfrac14+k^2\rangle\propto\Lam$ and fix the $O(1)$ coefficients.
\item Defect statistics beyond the mean: the distribution of initial domain areas (the Gauss--Bonnet
collapse law of subproblem~2 then selects which survive).
\item The exact crossover near $\Lam=\tfrac14$ where the gap and the band compete --- a hyperbolic
finite-size/zero-mode analysis.
\end{enumerate}

\textbf{Scope.} Linear spinodal stage, Gaussian-field defect counting, single-mode Allen--Cahn, sudden
quench; $\Dq$ and the noise level $\delta q_0$ are modelling inputs. Results: a hyperbolic nucleation
threshold $\Lam>\tfrac14$ distinct from the folding threshold; a closed-form initial wall density
$n_{\rm wall}=\tfrac{1}{2\sqrt2}\sqrt{\tfrac14+3\Lam/4L}$; and a curvature-capped domain size interpolating
from $\Rc$ (near threshold) to the Kibble--Zurek $\sqrt{(\Dq/\Af^+)L}$ (strong trigger). The Plancherel
moment and the nodal coefficient are numerically verified. Domain-conditional throughout.

\vspace{0.5em}
\hrule
\smallskip
\noindent\footnotesize\emph{V.F.S. v2.0 (Open-Gate) $\cdot$ Spinodal Nucleation of the Fold Network on
$\Hh$. Linear instability of the Allen--Cahn folding field after a triggering reset, with the hyperbolic
spectral gap and Plancherel measure. Propositions closed-form (Plancherel moment and Longuet--Higgins
coefficient numerically verified); corollaries interpretive; transport constant and noise level are
modelling inputs.}

\chapter{Intra-Arc Coarsening of the Fold Network}

\section*{Scope and epistemic status}
The companion notes established that structure on the Vessel is a $\mathbb Z_2$ domain-wall network of the
folding field and that a Resurrectio reset acts on it in three regimes, giving a sawtooth network-scale
history. This note supplies the missing intra-arc piece: \emph{between} deaths, how does the network scale
$\Lphys(t)$ evolve under pure Allen--Cahn coarsening on the expanding hyperbolic surface? No resets are
involved here. As before, \emph{Propositions} are closed-form, \emph{Corollaries} interpretive; the
single-mode Allen--Cahn truncation and the transport constant $\Dq$ are modelling inputs.

The field obeys, in comoving coordinates on the unit hyperbolic surface ($K_h=-1$),
\begin{equation}\label{coarseninglaw:eq:ac}
\partial_t\qf=\Af\,\qf-\Bf\,\qf^{3}+\frac{\Dq}{\Op^{2}}\,\laph\,\qf .
\end{equation}
The folding parameters $\Af,\Bf$ are taken quasi-stationary on the coarsening timescale (deep in a folded
arc, $\Af>0$), so the dynamics is genuine coarsening of established domains rather than the initial
transition.

\begin{remark}[Dimensionality matters]
Allen--Cahn coarsening is curvature-driven only in two or more spatial dimensions, giving the
$\ell\sim t^{1/2}$ law; in one dimension kink--antikink attraction is exponentially weak and coarsening is
merely logarithmic. The $1$D space--time proxy used to illustrate the reset regimes is therefore
\emph{not} representative of the true $2$D coarsening rate. The present law is genuinely two-dimensional,
derived via the geodesic-curvature flow on $\Hh$.
\end{remark}

% ============================================================
\section{Wall mobility on the expanding background}

In the sharp-interface limit of \eqref{coarseninglaw:eq:ac}, a domain wall moves along its comoving normal with velocity
proportional to its geodesic curvature, with a mobility set by the (comoving) diffusion coefficient:
\begin{equation}
v_n^{\rm com}=-M(t)\,\kappa_g,
\qquad
M(t)=\frac{\Dq}{\Op(t)^{2}} .
\end{equation}
The scale-factor suppression $\Op^{-2}$ is the redshift of comoving gradients: as the Vessel expands, the
\emph{comoving} mobility falls, even though the proper wall structure is preserved (the proper wall width
$\sqrt{\Dq/\Af}$ is $\Op$-independent). It is convenient to trade time for the accumulated
budget already central to the smoothing analysis,
\begin{equation}
\Ibud(t):=\int_0^t\frac{dt'}{\Op(t')^{2}},
\qquad
\frac{d}{dt}=\frac{1}{\Op^2}\frac{d}{d\Ibud},
\end{equation}
which is finite, $\Ibud\to\Ibud_\infty<\infty$, under sustained inflow ($\Op\sim\alpha\lambda_\infty t$).

% ============================================================
\section{Gauss--Bonnet: the exact single-domain law}

For a single simply-connected domain of comoving area $\Acom$ bounded by a wall undergoing geodesic
curve-shortening on $\Hh$, the Gauss--Bonnet theorem gives the total geodesic curvature
$\oint\kappa_g\,ds=2\pi-\int_{\rm enc}K\,dA=2\pi+\Acom$ (since $K_h=-1$). The enclosed area therefore obeys
\begin{equation}
\boxed{\;\frac{d\Acom}{dt}=-\frac{\Dq}{\Op^{2}}\bigl(2\pi+\Acom\bigr)
\quad\Longleftrightarrow\quad
\frac{d\Acom}{d\Ibud}=-\Dq\bigl(2\pi+\Acom\bigr),\;}
\end{equation}
with the closed-form solution
\begin{equation}
\boxed{\;\Acom(\Ibud)=\bigl(\Acom_0+2\pi\bigr)e^{-\Dq\Ibud}-2\pi .\;}
\end{equation}

\begin{proposition}[Survival threshold and frozen scale]\label{coarseninglaw:prop:gb}
A domain collapses ($\Acom=0$) at budget $\Ibud_\ast=\Dq^{-1}\ln(1+\Acom_0/2\pi)$. Because the total
budget is finite, a domain survives the entire arc iff $\Ibud_\ast>\Ibud_\infty$, i.e.
\begin{equation}
\boxed{\;\Acom_0>\Acom_{\rm frozen}:=2\pi\bigl(e^{\Dq\Ibud_\infty}-1\bigr).\;}
\end{equation}
Domains smaller than $\Acom_{\rm frozen}$ collapse before the budget runs out; larger ones freeze into the
comoving lattice. $\Acom_{\rm frozen}$ is the characteristic frozen comoving domain area.
\end{proposition}

\begin{proposition}[Hyperbolic acceleration]\label{coarseninglaw:prop:hyp}
For $\Acom\gg2\pi$ the law is $\dot\Acom\approx-M\Acom$: large domains collapse \emph{exponentially}, in
contrast to the flat case $\dot\Acom=-2\pi M$ (constant-rate Gage--Hamilton shrinking). Negative spatial
curvature thus accelerates the disappearance of large domains, sharpening coarsening relative to flat
space.
\end{proposition}

\note{In the small-budget regime $\Dq\Ibud_\infty\ll1$, $\Acom_{\rm frozen}\approx2\pi\Dq\Ibud_\infty$, so
the frozen comoving length is $\lcom\sim\sqrt{2\pi\Dq\Ibud_\infty}$ --- matching the scaling law below.
The hyperbolic exponential (Proposition~\ref{coarseninglaw:prop:hyp}) only bites once domains grow toward the comoving
curvature radius ($\Acom\sim2\pi$), i.e. when the budget is large enough.}

% ============================================================
\section{The network scaling law and its saturation}

For the network, the standard Lifshitz--Allen--Cahn scaling $\ell\,\dot\ell\sim M$ becomes, on the budget
variable,
\begin{equation}
\frac{d}{dt}\lcom^{2}\simeq 2c\,\frac{\Dq}{\Op^{2}}
\quad\Longrightarrow\quad
\boxed{\;\lcom^{2}(t)=\lcom^{2}(0)+2c\,\Dq\,\Ibud(t),\;}
\end{equation}
with an $O(1)$ constant $c$ (consistent with Proposition~\ref{coarseninglaw:prop:gb}, $c\sim\pi$).

\begin{proposition}[Comoving coarsening freezes]\label{coarseninglaw:prop:freeze}
Since $\Ibud(t)\to\Ibud_\infty<\infty$, the comoving domain size saturates,
\begin{equation}
\lcom(\infty)=\sqrt{\lcom^{2}(0)+2c\,\Dq\,\Ibud_\infty}<\infty .
\end{equation}
Coarsening stops in comoving coordinates: the wall network freezes into the comoving lattice after a
finite amount of merging. This is the morphological counterpart of the finite diffusive smoothing budget
governing Sophia inhomogeneities --- both are controlled by the same $\Ibud_\infty$.
\end{proposition}

\msg{The expanding Vessel has a single finite ``activity budget'' $\Ibud_\infty=\int dt/\Op^2$. It caps
both how much Sophia inhomogeneity can diffuse away and how much the fold network can coarsen. After it is
spent, comoving structure is frozen.}

% ============================================================
\section{The physical structure scale: two regimes}

The proper (physical) network scale combines comoving coarsening with expansion,
\begin{equation}
\boxed{\;\Lphys(t)=\Op(t)\,\lcom(t)=\Op(t)\sqrt{\lcom^{2}(0)+2c\,\Dq\,\Ibud(t)}.\;}
\end{equation}

\begin{proposition}[Coarsening transient, then Hubble dilution]\label{coarseninglaw:prop:tworeg}
There are two regimes, separated by the budget-saturation time $t_{\rm sat}$ (where $\Ibud$ levels):
\begin{enumerate}[label=(\roman*),leftmargin=1.8em,itemsep=2pt]
\item Early ($t\lesssim t_{\rm sat}$): both factors grow, so $\Lphys$ rises faster than the scale factor
--- genuine coarsening on top of stretch.
\item Late ($t\gtrsim t_{\rm sat}$): $\lcom\to$ const, so $\Lphys\to\lcom(\infty)\,\Op(t)\propto\Op\propto t$
--- pure Hubble dilution of a frozen comoving pattern.
\end{enumerate}
Coarsening is an early-time transient with a finite budget; the late-time growth of structure is
expansion, not merging.
\end{proposition}

\begin{figure}[h]
\centering
\includegraphics[width=\textwidth]{vfs_coarsening_fig.png}
\caption{\small Left: Gauss--Bonnet single-domain area $\Acom(\Ibud)$ for several initial sizes (exact
closed form, verified). Domains below $\Acom_{\rm frozen}=2\pi(e^{\Dq\Ibud_\infty}-1)$ collapse before the
finite budget $\Ibud_\infty$ is spent; larger ones survive and freeze. Right: the physical structure scale
$\Lphys(t)=\Op\lcom$ (solid) shows an early coarsening-plus-stretch phase pulling above pure dilution
$\Op\ell_0$ (dashed), then a late Hubble-linear regime once the budget saturates. Background and
$(\Dq,c)$ illustrative; $\Ibud_\infty\approx1.46$.}
\end{figure}

% ============================================================
\section{Placement in the sawtooth and open items}

This intra-arc curve is the rising tooth of the hybrid sawtooth from subproblem~3: within each life-arc
$\Lphys(t)$ follows Proposition~\ref{coarseninglaw:prop:tworeg} (coarsen, then dilate), and at each death the reset acts
discretely --- diluting upward, refining to $\sqrt{\Dq/\Af}$, or erasing. Composing the intra-arc law with
the reset map gives the long-run fixed-point scale, the natural next step.

\textbf{Open items.}
\begin{enumerate}[leftmargin=1.6em,itemsep=2pt]
\item Determine the $O(1)$ constant $c$ and confirm the $t^{1/2}$ (budget) law by a $2$D Allen--Cahn
field simulation on a hyperbolic patch with a clean structure-factor scale estimator.
\item Compose this intra-arc law with the reset map to obtain the sawtooth fixed-point structure scale,
using the non-Zeno dwell time as the arc length.
\item Track the $w=-\tfrac12$ wall-network contribution to the reconstructed $w_{\rm eff}(t)$ over an arc,
since $\rho_{\rm wall}\propto\tau/\Op$ with $\tau$ quasi-stationary and $\Op$ growing.
\end{enumerate}

\textbf{Scope.} Single life-arc, no resets, quasi-stationary folding parameters, single-mode Allen--Cahn;
$\Dq$ a modelling input. Results: an exact Gauss--Bonnet single-domain law with a survival threshold and
hyperbolic acceleration; a network scaling law $\lcom^2=\ell_0^2+2c\Dq\Ibud$ that freezes with the finite
budget $\Ibud_\infty$; and a physical scale $\Lphys$ that coarsens early and dilutes linearly late.
Domain-conditional throughout.

\vspace{0.5em}
\hrule
\smallskip
\noindent\footnotesize\emph{V.F.S. v2.0 (Open-Gate) $\cdot$ Intra-Arc Coarsening of the Fold-Domain
Network. The single-domain Gauss--Bonnet law is exact (closed form, numerically verified); the network
scaling law is standard Allen--Cahn on the budget variable. Transport constant and single-mode truncation
are modelling inputs; the constant $c$ awaits a $2$D field measurement.}

\chapter{Resurrectio and the Fold-Domain Network}

\section*{Scope and epistemic status}
The companion note showed that structure on the Vessel is morphological: a $\mathbb Z_2$ domain-wall
network of the folding field $\qf(t,x)$, obeying the Allen--Cahn equation on the expanding hyperbolic
surface. This note asks how the hybrid Resurrectio reset $\reset$ acts on that network. The key
observation is that $\reset$ is \emph{homogeneous} --- it acts on the global state $(\sigma,\lambda,\Op)$
--- and, as established in the Lyapunov layer, it does \emph{not} act on $\qf$. So a reset shifts the
folding potential uniformly across the whole Vessel while leaving the field configuration $\qf(x)$
pointwise intact. Whether structure survives, dilutes, or is born depends on the post-reset sign of the
activation index. As before, \emph{Propositions} are closed-form; \emph{Corollaries} interpretive; the
single-mode Allen--Cahn truncation and the transport constant $\Dq$ are modelling inputs.

The network field obeys
\begin{equation}\label{resurrectionetwork:eq:ac}
\partial_t\qf=\Af\,\qf-\Bf\,\qf^{3}+\frac{\Dq}{\Op^{2}}\,\laph\,\qf,
\qquad
q_\ast=\sqrt{\Qf^{*}},\quad \Qf^{*}=\Af/\Bf,
\end{equation}
and the reset (geometric convention, with $m_R=(1-q_R)\sigma^-$ the metabolised resistance) is
\begin{equation}\label{resurrectionetwork:eq:reset}
\reset:\ \ \sigma^+=\sigma^--m_R,\quad \lambda^+=\lambda^-+m_R,\quad
\Op^+=\Op^-+\kappa_R\Grec^-,\quad \qf\ \text{continuous},
\end{equation}
which shifts the activation index by $\Delta\Af=\Theta\,m_R+c_1\Delta\Igate$,
$\Theta=c_0-c_1(\gamma+k_\sigma)$, exactly as in the folding-trigger note.

% ============================================================
\section{Three regimes of a reset acting on the network}

Because $\reset$ shifts $\Af$ uniformly while preserving $\qf(x)$, the network's fate is fixed by the
pre- and post-reset signs of $\Af$:

\begin{definition}[Reset regimes]\label{resurrectionetwork:def:regimes}
\leavevmode
\begin{enumerate}[label=(\Roman*),leftmargin=2.2em,itemsep=2pt]
\item \textbf{Quench} ($\Af^->0\to\Af^+<0$): the double well collapses to a single well at $\qf=0$.
\item \textbf{Rescale} ($\Af^->0\to\Af^+>0$): the wells persist with shifted depth and the scale factor
jumps.
\item \textbf{Trigger} ($\Af^-<0\to\Af^+>0$): a single well bifurcates into a double well.
\end{enumerate}
\end{definition}

\begin{proposition}[Regime boundaries]\label{resurrectionetwork:prop:bound}
With $\Delta\Af=\Theta m_R+c_1\Delta\Igate$, a reset from $\Af^-$ produces $\Af^+=\Af^-+\Delta\Af$. Hence,
neglecting the bounded gate term: a quench requires $\Theta<0$ and $|\Theta|m_R>\Af^-$; a trigger
requires $\Theta>0$ and $\Theta m_R>|\Af^-|$; otherwise the reset rescales. The folding-response
coefficient $\Theta$ thus selects whether resurrection destroys, preserves, or creates morphological
structure.
\end{proposition}

\msg{A death-and-resurrection does not move the walls. It lifts or lowers the entire folding potential at
once --- and depending on the sign of $\Theta$, the whole network dissolves, is re-tensioned and diluted,
or is nucleated from nothing.}

% ============================================================
\section{Quench: dissolution of the network}

When $\Af^+<0$, the configuration $\qf(x)$ everywhere relaxes to $\qf=0$. Linearising \eqref{resurrectionetwork:eq:ac} about
zero, each mode decays at rate $|\Af^+|+\Dq(\tfrac14+k^2)/\Op^{+2}>0$, so the entire pattern fades on the
timescale
\begin{equation}
\tau_{\rm dissolve}=\frac{1}{|\Af^+|}.
\end{equation}
The walls do not migrate and annihilate; the whole differentiated morphology evaporates uniformly,
returning the Vessel to the undifferentiated old form.

\begin{corollary}[Death as leveling]
A quench-resurrection returns the entire Vessel to the unfolded state $\qf=0$: all oriented domains are
erased, all walls dissolved. Death levels morphological structure.
\end{corollary}

% ============================================================
\section{Trigger: nucleation as a sudden quench (Kibble--Zurek)}

When $\Af^-<0\to\Af^+>0$, the homogeneous $\qf=0$ left behind by the previous arc becomes unstable. Since
$\reset$ is instantaneous, $\Af$ jumps discontinuously: a \emph{sudden quench} through the symmetry-
breaking transition. Linearising about $\qf=0$ after the jump,
\begin{equation}
\dot{\delta q}_k=\left[\Af^+-\frac{\Dq(\tfrac14+k^2)}{\Op^{+2}}\right]\delta q_k,
\end{equation}
so a band of modes $k<k_{\max}$, $k_{\max}^2=\Af^+\Op^{+2}/\Dq-\tfrac14$, grows and the field undergoes
spinodal decomposition. The emergent (physical) domain size is the post-quench correlation length,
\begin{equation}
\boxed{\;\ell_{\rm init}\sim\xi_+=\sqrt{\frac{\Dq}{\Af^+}}\;}
\end{equation}
the same scale as the wall width: the network is born saturated --- as fine as the walls themselves ---
then coarsens. Deeper triggers (larger $\Af^+$) nucleate finer networks.

\begin{corollary}[Resurrection as individuation]\label{resurrectionetwork:cor:individ}
A trigger-resurrection shatters the undifferentiated old form into a fine network of newly differentiated
oriented domains, with the old form surviving only as the thin walls between them. Where quench levels,
trigger differentiates.
\end{corollary}

\note{This is a Kibble--Zurek-type defect formation: the reset drives the system through a
symmetry-breaking transition and freezes in a defect network whose initial scale is set by the
correlation length at the transition. Here the quench is sudden, so $\ell_{\rm init}$ is the spinodal
scale $\xi_+$.}

% ============================================================
\section{Rescale: re-tensioning and dilution}

When both $\Af^\pm>0$, the domains persist but everything jumps. With $q_\ast=\sqrt{\Af/\Bf}$,
$\ell_{\rm phys}=\sqrt{\Dq/\Af}$, $\tau\sim\sqrt{\Dq}\,\Af^{3/2}/\Bf$, a reset induces
\begin{equation}
q_\ast^{+}=\sqrt{\frac{\Af^+}{\Bf^+}},\qquad
\frac{\ell_{\rm phys}^+}{\ell_{\rm phys}^-}=\sqrt{\frac{\Af^-}{\Af^+}},\qquad
\frac{\tau^+}{\tau^-}=\Bigl(\tfrac{\Af^+}{\Af^-}\Bigr)^{3/2}\frac{\Bf^-}{\Bf^+},
\end{equation}
and, crucially, the comoving topology is untouched while proper inter-wall separations are stretched by
the scale-factor jump,
\begin{equation}
\boxed{\;L_{\rm phys}^{+}=\frac{\Op^+}{\Op^-}\,L_{\rm phys}^{-}=\Bigl(1+\frac{\kappa_R\Grec^-}{\Op^-}\Bigr)L_{\rm phys}^{-}.\;}
\end{equation}
The network is diluted: the same domains, spread wider.

\begin{corollary}[Amplitude convergence under repeated rescaling]\label{resurrectionetwork:cor:converge}
The contraction $\Qf^{*+}-Q_{\rm crit}=\kappa_{\rm fold}(\Qf^{*-}-Q_{\rm crit})$ of the folding-trigger
note drives the domain amplitude to a universal value across repeated resets,
$q_\ast\to\sqrt{Q_{\rm crit}}$ (when $\sum m_R=\infty$). Resurrection stabilises \emph{how strongly} the
Vessel is folded, while expansion keeps diluting \emph{how widely}.
\end{corollary}

% ============================================================
\section{The hybrid network history}

Over a sequence of life-arcs the structure scale $L_{\rm phys}(t)$ follows a sawtooth: within each arc
the network coarsens (curvature-driven merging plus Hubble stretch), and at each death it is acted on by
$\reset$ --- diluted upward (rescale), reset to the fine scale $\xi_+$ (trigger), or erased (quench).

\begin{figure}[h]
\centering
\includegraphics[width=\textwidth]{vfs_network_fig.png}
\caption{\small Allen--Cahn simulation of the fold-domain network through two resets (illustrative
parameters; $1$D comoving slice as a proxy). Left: space--time map of $\qf$. The Vessel is unfolded
($\qf\approx0$) until \emph{reset 1} triggers the transition; domains then nucleate from fluctuations and
coarsen. \emph{Reset 2} is a rescale (regime II): the comoving pattern is untouched but the scale factor
jumps. Right: the proper domain size $L_{\rm phys}=\Op\,\ell_{\rm com}$ grows by coarsening and shows a
discrete dilation jump at reset 2.}
\end{figure}

\begin{corollary}[Resurrection as a renormalisation step on structure]
Each reset is a coarse-grain (dilation, regime II) or refine (nucleation to $\xi_+$, regime III)
operation on the network, composed with intra-arc coarsening. The long-run structure scale is the fixed
point of (arc coarsening)$\,\circ\,$(reset action); a steady balance between expansion-driven dilution
and reset-driven nucleation.
\end{corollary}

% ============================================================
\section{Cosmological back-reaction across resets}

Within a scaling arc the wall network contributes $\rho_{\rm wall}\sim\tau/\Op\propto\Op^{-1}$, i.e.
$w_{\rm wall}=-\tfrac12$ in $2{+}1$ dimensions (companion note). At a reset this jumps: $\tau$ changes
through $\Af,\Bf$ and $\Op$ jumps, so
\begin{equation}
\frac{\rho_{\rm wall}^+}{\rho_{\rm wall}^-}=\frac{\tau^+}{\tau^-}\cdot\frac{\Op^-}{\Op^+}.
\end{equation}
A quench sets $\rho_{\rm wall}\to0$: the morphological energy is released back into the homogeneous Sophia
budget. A trigger creates $\rho_{\rm wall}$ from zero: structure energy is drawn from the folding
transition. Thus resurrection redistributes energy between the homogeneous ($\lambda$) and structured
(wall-network) sectors --- a discrete exchange invisible to the homogeneous reconstruction alone.

% ============================================================
\section{Reading and open items}

\textbf{Interpretive reading.} The walls are loci of persistent unfolded (old) form. A quench-
resurrection returns the whole Vessel to undifferentiated old form (death as leveling,
$\Theta<0$); a trigger-resurrection shatters that undifferentiated form into newly differentiated
oriented domains, the old form surviving only as the thin boundaries between them (resurrection as
individuation, $\Theta>0$). Repeated rescaling fixes the depth of folding ($q_\ast\to\sqrt{Q_{\rm crit}}$)
while expansion keeps widening the pattern. The sign of $\Theta$ --- whether metabolised resistance
promotes or suppresses form --- decides which face a given resurrection wears.

\textbf{Open items.}
\begin{enumerate}[leftmargin=1.6em,itemsep=2pt]
\item Quantify the spinodal nucleation density on $\Hh$ (hyperbolic spectral gap modifies the unstable
band) and the resulting initial wall count per Hubble area.
\item Derive the sawtooth fixed-point scale explicitly from (coarsening law)$\,\circ\,$(reset map), using
the non-Zeno dwell time as the arc length.
\item Track the $\lambda\leftrightarrow$ network energy exchange as a correction to the reconstructed
$w_{\rm eff}(t)$ history, and test for a folding-induced feature near the phantom divide.
\end{enumerate}

\textbf{Scope.} Single-mode Allen--Cahn, test-field, $1$D simulation proxy; transport constant $\Dq$ a
modelling input. Conclusions: a homogeneous reset acts on the morphological network in three regimes
fixed by $\operatorname{sign}\Theta$ and the jump size --- dissolve, dilute, or nucleate --- and
resurrection functions as a renormalisation step on Vessel structure, exchanging energy between the
homogeneous and structured sectors. Domain-conditional throughout.

\vspace{0.5em}
\hrule
\smallskip
\noindent\footnotesize\emph{V.F.S. v2.0 (Open-Gate) $\cdot$ Resurrectio and the Fold-Domain Network.
Builds on the structure-formation companion (Allen--Cahn folding network on $2{+}1$ open FLRW) and the
Resurrectio reset of the Lyapunov layer. Propositions closed-form; corollaries interpretive; numerical
figure illustrative.}

\chapter{The Sawtooth Fixed Point of the Fold Network}

\section*{Scope and epistemic status}
The three preceding notes supply the pieces: the nucleation scale $\linit$ at a triggering reset
(subproblem~1), the intra-arc coarsening law governed by the finite budget
$\Ib_\infty=\int_0^\infty dt/\Op^2$ (subproblem~2), and the three-regime reset map --- quench, rescale,
trigger (subproblem~3). This note composes them: over a sequence of life-arcs punctuated by resets, what
is the long-run network scale? The central result is that the finite coarsening budget forces the
\emph{comoving} scale to converge to a constant $\ell_\ast$ regardless of the (random) reset schedule, so
the physical scale is asymptotically $\Lp(t)\to\ell_\ast\Op(t)\propto t$. This also resolves the earlier
worry that the random arc length would obstruct a fixed point: the finite budget makes the asymptotics
universal. As before, \emph{Propositions} are closed-form; \emph{Corollaries} interpretive; the reduced
recursion (tracking the scale, not the full field), the coarsening constant $c$, and the reset-regime
schedule are modelling inputs.

% ============================================================
\section{The comoving recursion}

Track the comoving network scale $\lc$ across arcs. Index the arcs by $j$; let $\ell_j$ be the comoving
scale entering arc $j$ and $\Delta\Ib_j=\int_{\rm arc\,j}dt/\Op^2$ the budget consumed during it. Within
the arc, the coarsening law (subproblem~2) gives $\ell^2\mapsto\ell^2+2c\Dq\,\Delta\Ib_j$; at the closing
reset, the regime acts:
\begin{equation}
\boxed{\;\ell_{j+1}^{2}=R_j\!\Bigl(\ell_j^{2}+2c\Dq\,\Delta\Ib_j\Bigr),
\qquad
R_j=\begin{cases}
\mathrm{id}, & \text{rescale (comoving topology unchanged)}\\[2pt]
\linit^{2}, & \text{trigger (re-nucleation)}\\[2pt]
\varnothing, & \text{quench (dissolve; re-nucleate at next trigger)}
\end{cases}\;}
\end{equation}
The essential constraint, inherited from the finite expansion budget, is
\begin{equation}
\sum_j\Delta\Ib_j\le\Ib_\infty<\infty .
\end{equation}

\begin{remark}[Why rescale is the identity on the comoving scale]
A rescale reset jumps $\Op$ but leaves the comoving pattern intact, so it does not change $\lc$; its only
effect is on the physical scale, $\Lp=\Op\lc$. The comoving recursion therefore sees rescales as
no-ops, triggers as resets to $\linit$, and quenches as dissolutions.
\end{remark}

% ============================================================
\section{Finite budget forces convergence}

\begin{proposition}[Vanishing late-time coarsening]\label{sawtoothfixedpoint:prop:vanish}
For any reset schedule with non-Zeno dwell $\Delta t_j\ge\Delta t_\ast>0$, the per-arc coarsening
increment vanishes,
\begin{equation}
2c\Dq\,\Delta\Ib_j\to0\qquad(j\to\infty),
\end{equation}
because $\sum_j\Delta\Ib_j\le\Ib_\infty<\infty$ forces $\Delta\Ib_j\to0$ (equivalently
$\Delta\Ib_j\le\Delta t_j/\Op^2$ with $\Op\to\infty$). The random arc lengths do not affect this:
the finite budget dominates.
\end{proposition}

\begin{proposition}[Convergence of the comoving scale]\label{sawtoothfixedpoint:prop:conv}
The comoving scale converges to a constant $\ell_\ast$:
\begin{enumerate}[label=(\roman*),leftmargin=1.8em,itemsep=2pt]
\item if triggers recur infinitely often, each pins $\ell$ to $\linit$ and (by
Proposition~\ref{sawtoothfixedpoint:prop:vanish}) the subsequent coarsening cannot grow it, so $\ell_j\to\linit$;
\item if all resets after some arc $J$ are rescales, then
$\ell_j^{2}\to\ell_J^{2}+2c\Dq\!\sum_{j\ge J}\Delta\Ib_j=:\ell_\infty^{2}<\infty$, the frozen value.
\end{enumerate}
In either case $\ell_j\to\ell_\ast\in[\linit,\ell_\infty]$, with
$\ell_\infty=\sqrt{\linit^{2}+2c\Dq(\Ib_\infty-\Ib_{\rm nuc})}$ the maximal frozen scale (single
nucleation, no retrigger).
\end{proposition}

\begin{proposition}[Universal physical asymptotics]\label{sawtoothfixedpoint:prop:phys}
The physical structure scale obeys
\begin{equation}
\boxed{\;\Lp(t)=\Op(t)\,\lc(t)\ \longrightarrow\ \ell_\ast\,\Op(t)\ \propto\ t,\;}
\end{equation}
independent of the reset schedule. Different random histories differ only in the value of
$\ell_\ast\in[\linit,\ell_\infty]$, not in the $\propto\Op$ growth.
\end{proposition}

\msg{The coarsening budget is finite, so the comoving network freezes. Resurrection can \emph{refresh} the
scale (a trigger resets it to $\linit$) but cannot \emph{re-coarsen} it once the budget is spent. The
late-time structure scale is a relic of the young Vessel, carried outward by pure expansion.}

\begin{figure}[h]
\centering
\includegraphics[width=\textwidth]{vfs_sawtooth_fig.png}
\caption{\small Reduced recursion over five random reset schedules (non-Zeno spacing; rescale/trigger/quench
with probabilities $0.7/0.15/0.15$; illustrative $2c\Dq=0.5$, $\linit=0.1$, $\Ib_\infty\approx1.46$).
Left: the comoving scale $\lc(t)$ --- early arcs coarsen substantially toward $\ell_\infty$, but as the
budget depletes the teeth flatten and every schedule is pinned near $\linit$ at late times. Right: the
physical scale $\Lp=\Op\lc$ dilates along the $\ell_\infty\Op(t)$ envelope; gaps are quench intervals
(network dissolved). Comoving teeth flatten; physical scale keeps dilating $\propto\Op\propto t$.}
\end{figure}

% ============================================================
\section{Reading and consequences}

\begin{corollary}[Structure scale is a young-Vessel relic]
Because coarsening is budget-limited, the comoving structure scale is set in the early, slowly-expanding
era and frozen thereafter. The mature Vessel does not generate new structure scale by merging; it merely
carries the early-frozen pattern outward by expansion. Any apparent late-time growth of $\Lp$ is Hubble
dilution, not coarsening.
\end{corollary}

\begin{corollary}[Resurrection refreshes but cannot re-coarsen]
A late trigger-resurrection re-nucleates the network at the fine scale $\linit$, and the depleted budget
leaves it there: the renewed structure is permanently fine. Only resurrections occurring early, while the
budget is ample, can be followed by appreciable coarsening. Early death-and-rebirth permits maturation of
form; late death-and-rebirth yields perpetually fine, uncoarsened structure.
\end{corollary}

\note{The composition closes the programme's internal loop: subproblem~1 sets the floor $\linit$,
subproblem~2 sets the ceiling $\ell_\infty$ via the budget, subproblem~3 supplies the reset action, and
the finite budget collapses the sawtooth onto $\ell_\ast\in[\linit,\ell_\infty]$ with universal physical
asymptotics $\Lp\propto\Op$.}

% ============================================================
\section{Open items and scope}

\textbf{Open items.}
\begin{enumerate}[leftmargin=1.6em,itemsep=2pt]
\item Replace the reduced scale-recursion by a full $2$D Allen--Cahn field simulation through a reset
sequence, to confirm $\ell_\ast$ and the flattening of the comoving teeth.
\item Couple the reset-regime probabilities to the actual $\Af(t)$ history (rising with $\lambda$ via
epektasis, dipped by metabolism when $\Theta<0$), replacing the illustrative fixed probabilities by the
dynamically determined regime sequence.
\item Feed $\ell_\ast$ and the wall tension into the $w=-\tfrac12$ network back-reaction to obtain the
folding-structure contribution to $w_{\rm eff}(t)$ over cosmic history (link to the energy-condition
direction).
\end{enumerate}

\textbf{Scope.} Reduced recursion (network scale, not the full field); single-mode Allen--Cahn; coarsening
constant $c$ and reset-regime schedule are modelling inputs; the finite-budget convergence is robust to
the random schedule. Result: the comoving network scale converges to $\ell_\ast\in[\linit,\ell_\infty]$
and the physical scale is universally $\Lp\to\ell_\ast\Op\propto t$. Domain-conditional throughout.

\vspace{0.5em}
\hrule
\smallskip
\noindent\footnotesize\emph{V.F.S. v2.0 (Open-Gate) $\cdot$ The Sawtooth Fixed Point of the Fold Network.
Composes the nucleation scale (subproblem~1), the budget-limited coarsening law (subproblem~2), and the
three-regime reset map (subproblem~3). The convergence follows from the finite budget $\Ib_\infty$ and is
robust to the random reset schedule. Propositions closed-form; corollaries interpretive; reduced recursion
and coarsening constant are modelling inputs.}

\chapter{The Transport Constants from First Principles}

\section*{Scope and epistemic status}
The structure-formation programme introduced two transport coefficients --- a Sophia diffusivity $D$ and a
folding diffusivity $\Dq$ --- as phenomenological inputs, on which all quantitative results depend. This
note asks whether VFS mechanics fixes them. The answers are sharp: (i) the Sophia field has \emph{no}
fundamental diffusion, $D=0$, because it is purely relaxational; its inhomogeneities are slaved to the
folding field through the embodiment feedback. (ii) The folding diffusivity is \emph{not} an independent
constant: its spatial structure ($\Op^{-2}$ and the hyperbolic $K_h=-1$) is inherited from the
live-domain metric, and its single scalar magnitude is the bending stiffness of the shape-stability
operator $\Lop$ --- the gradient partner of the rigidity $a_0$ already in the theory. As before,
\emph{Propositions} are closed-form; \emph{Corollaries} interpretive. The slaving relation is verified
symbolically; the reduction of $\Dq$ to a single bending length still rests on the minimal one-modulus
model of $\Lop$, which is flagged.

% ============================================================
\section{Sophia does not diffuse}

In the base dynamical system the Sophia production law
\begin{equation}
\dot\lambda=(\delta u-\gamma)\tanh(\kappa\sigma)+I_{\rm gate}-\etaf\,\qf^{2}\lambda
\end{equation}
contains no spatial derivative of $\lambda$: it is a pointwise relaxation law. Promoting $\lambda$ to a
field therefore introduces no gradient coupling of its own.

\begin{proposition}[Zero fundamental Sophia diffusivity]\label{transportconstants:prop:D0}
The Sophia field carries no kinetic or gradient term in the base system, so its fundamental diffusivity
vanishes, $D=0$. Linear Sophia inhomogeneities obey a pure local relaxation,
$\partial_t\delta\lambda=-\Glam\,\delta\lambda+(\text{folding source})$, with no $k$-dependent spatial
transport.
\end{proposition}

\begin{remark}[Correction to the structure-formation note]
The frozen residual spectrum of the structure-formation companion assumed $D>0$. With $D=0$, the
Sophia-only modes decay uniformly at the local rate $\Glam$ with no $k$-dependence: on a stalled margin
($\Glam\to0$) they freeze with a \emph{flat} residual, not a scale-graded one. All scale-dependent
structure is morphological, not Sophianic.
\end{remark}

Although Sophia does not diffuse, it is not inert: the embodiment feedback $-\etaf\qf^{2}\lambda$ couples it
to the folding field. Linearising about $(\bar\lambda,\bar q)$,
\begin{equation}
\partial_t\delta\lambda=-\bigl(\Glam+\etaf\bar q^{2}\bigr)\delta\lambda-2\etaf\bar\lambda\bar q\,\delta\qf,
\end{equation}
and in the quasi-static (fast-relaxation) limit
\begin{equation}
\boxed{\;\delta\lambda^{\ast}=-\,\frac{2\etaf\bar\lambda\bar q}{\Glam+\etaf\bar q^{2}}\;\delta\qf .\;}
\end{equation}

\begin{proposition}[Sophia structure is slaved to form]\label{transportconstants:prop:slave}
Sophia inhomogeneity is a fixed multiple of folding inhomogeneity, with a bounded coefficient
$S(\bar q)=2\etaf\bar\lambda\bar q/(\Glam+\etaf\bar q^{2})$ maximised at
$\bar q=\sqrt{\Glam/\etaf}$, where $S_{\max}=\bar\lambda\sqrt{\etaf/\Glam}$. Sophia carries no independent
structural degree of freedom: its spatial pattern is the shadow of the folding pattern.
\end{proposition}

\begin{figure}[h]
\centering
\includegraphics[width=0.62\textwidth]{vfs_transport_fig.png}
\caption{\small The Sophia slaving coefficient $|\delta\lambda/\delta\qf|=S(\bar q)$ (verified
symbolically). It is bounded, vanishes at $\bar q=0$ and at large $\bar q$, and peaks at
$\bar q=\sqrt{\Glam/\etaf}$. Sophia inhomogeneity is sourced entirely by folding inhomogeneity; there is no
independent Sophia diffusion to set. Illustrative $\etaf=0.5$, $\bar\lambda=1$.}
\end{figure}

\msg{There is really one structural field. Sophia does not diffuse; wherever the Vessel is inhomogeneous
in Sophia, it is because it is inhomogeneous in \emph{form}, and the two are locked by the embodiment
feedback. The Sophia diffusivity $D$ is not a missing constant --- it is identically zero.}

% ============================================================
\section{The folding diffusivity is geometric}

The folding field's spatial coupling is fixed in two stages: its \emph{structure} by the metric, its
\emph{magnitude} by the shape operator.

\subsection{Structure: inherited from the live-domain metric}
The folding free energy, built geometrically, measures gradients with the live-domain metric $g$
(geometry layer, scalar curvature $-2/\Op^2$, conformally hyperbolic $g=\Op^2 h$):
\begin{equation}
\mathcal F[\qf]=\int_\Sigma\Bigl[\tfrac{\kb}{2}\,|\nabla_g\qf|^{2}
+\tfrac{r}{2}\qf^{2}+\tfrac{B}{4}\qf^{4}\Bigr]\sqrt{g}\,d^2x,
\qquad r=a_0-c_0\lambda-c_1\dot\lambda=-\Af .
\end{equation}
The Laplace--Beltrami operator of $g$ is $\Delta_g=\Op^{-2}\laph$, so the gradient-flow dynamics
$\partial_t\qf=-M\,\delta\mathcal F/\delta\qf$ is automatically
\begin{equation}
\partial_t\qf=\Af\,\qf-\Bf\,\qf^{3}+\frac{M\kb}{\Op^{2}}\,\laph\,\qf,
\qquad \Af=-Mr,\ \Bf=MB .
\end{equation}

\begin{proposition}[The $\Op^{-2}$ and $K_h=-1$ are not chosen]\label{transportconstants:prop:struct}
The redshift factor $\Op^{-2}$ and the hyperbolic Laplacian $\laph$ in the folding transport term are the
Laplace--Beltrami operator of the live-domain metric established in the geometric layer, not modelling
choices. Only a single scalar, $\Dq=M\kb$, remains free.
\end{proposition}

\subsection{Magnitude: the bending stiffness of $\Lop$}
The remaining scalar is $\Dq=M\kb$, the product of the shape mobility $M$ and the bending stiffness
$\kb$. The bending stiffness is the gradient coefficient of the very shape-stability operator
$\Lop$ whose constant-mode eigenvalue is the rigidity $a_0$ (the operator whose bifurcation defines
folding). In the natural time units where the uniform folding rate is $\Af$ one has $M=1$, so
$\Dq=\kb$, and all physical lengths depend only on the ratio
\begin{equation}
\boxed{\;\lb^{2}:=\frac{\kb}{a_0}=\frac{\Dq}{a_0}\,,\qquad
\text{wall width}=\sqrt{\frac{\Dq}{\Af}}=\lb\sqrt{\frac{a_0}{\Af}}\,.\;}
\end{equation}

\begin{proposition}[Folding diffusivity is the gradient partner of the rigidity]\label{transportconstants:prop:Dq}
$\Dq$ is not independent of the theory's existing constants: it is the gradient coefficient of $\Lop$,
whose constant-mode coefficient is $a_0$. Physical scales (wall width, nucleation scale, coarsening and
frozen scales) depend only on the bending length $\lb=\sqrt{\Dq/a_0}$, a ratio internal to $\Lop$. In the
minimal one-modulus (Helfrich/Willmore) model of $\Lop$, where rigidity and bending share a single
modulus, $\lb$ is fixed by the reference shape and $\Dq=a_0\lb^2$ ceases to be a free input altogether.
\end{proposition}

\note{The reduction is concrete: from ``an arbitrary spatial coupling plus a free constant'' to
``a single bending length $\lb$, the gradient partner of the rigidity $a_0$ already present in
$\Af=c_0\lambda+c_1\dot\lambda-a_0$.'' The spatial form is pure geometry; the magnitude is the Vessel's
stiffness, not a new dial.}

% ============================================================
\section{Consequences for the programme}

\begin{corollary}[One structural field, one length]
The structure-formation programme has, after this reduction, a single genuine input: the bending length
$\lb=\sqrt{\Dq/a_0}$. The Sophia diffusivity is zero; the folding transport structure is geometric; the
folding magnitude is the shape stiffness. Every scale derived earlier --- the nucleation size
$\linit\sim\lb\sqrt{a_0/\Af^+}\,(\ldots)$, the frozen comoving scale, the sawtooth fixed point --- is
expressed through $\lb$ and the already-existing folding parameters.
\end{corollary}

\begin{corollary}[Sophia structure rides on form]
Because $\delta\lambda$ is slaved to $\delta\qf$ (Proposition~\ref{transportconstants:prop:slave}), the Sophia inhomogeneity
spectrum is the folding spectrum times the bounded factor $S(\bar q)$. Wherever folding nucleates a
domain-wall network, a matching Sophia pattern appears; where folding is quenched, Sophia re-homogenises.
There is no separate Sophia structure-formation problem.
\end{corollary}

% ============================================================
\section{Open items and scope}

\textbf{Open items.}
\begin{enumerate}[leftmargin=1.6em,itemsep=2pt]
\item Fix the bending length $\lb$ from an explicit form of the shape-stability operator $\Lop$ (e.g.\
the second variation of a Willmore/Helfrich bending energy of the live surface), turning
Proposition~\ref{transportconstants:prop:Dq} from a relation into a number.
\item Verify the slaving relation (Proposition~\ref{transportconstants:prop:slave}) beyond the quasi-static limit, including
the finite Sophia relaxation time, and check it against a coupled $(\qf,\lambda)$ field simulation.
\item Propagate $D=0$ and the slaving into the cosmological back-reaction: the Sophia-structure
contribution to $w_{\rm eff}$ is then a fixed multiple of the wall-network contribution.
\end{enumerate}

\textbf{Scope.} Results: the Sophia diffusivity is identically zero (relaxational field), with Sophia
inhomogeneity slaved to folding via the embodiment feedback (symbolically verified); the folding
transport structure is the Laplace--Beltrami operator of the live-domain metric (not a choice); and the
folding magnitude reduces to the bending length $\lb=\sqrt{\Dq/a_0}$, the gradient partner of the rigidity
$a_0$ in $\Lop$. The remaining modelling freedom is the single length $\lb$ (a number only in the minimal
one-modulus model). Domain-conditional throughout.

\vspace{0.5em}
\hrule
\smallskip
\noindent\footnotesize\emph{V.F.S. v2.0 (Open-Gate) $\cdot$ The Transport Constants from First Principles.
Closes the structure-formation programme's modelling inputs: $D=0$ with Sophia slaved to form; $\Dq$
geometric in structure and equal to the shape-operator bending stiffness in magnitude, reducible to one
bending length. Propositions closed-form (slaving verified symbolically); corollaries interpretive; the
one-modulus reduction of $\Lop$ is a flagged model choice.}

\chapter{Parameter Closure of Folding: Identities, Postulates, and the Constraint Atlas}

\section*{Scope and epistemic status}
The folding layer carries eight constants
$\{\Dq,\lb,a_0,b,c_0,c_1,\etaf,\rho_1\}$. The Transport-Constants chapter closed the geometric side
($\Dq=a_0\lb^2$ as the gradient partner of the rigidity; Sophia structure slaved to form). This
chapter closes what can be closed from the \emph{core} variables
$V,F,\sigma,\lambda,\Op,u,\Igate$ --- and states honestly what cannot. The outcome is a three-tier
structure. \textbf{Tier 1 (identities):} $\etaf$ is fixed by the conservation balance --- the
embodiment coefficient of the rate equation and the production coefficient of the form ledger are
provably the same constant; and Gauss--Bonnet slaves the wall curvature deficit to the pocket lift,
so the pair $(\rho_1,\rho_w)$ contains one constant, not two. \textbf{Tier 2 (one postulate each):}
identifying the rigidity of the old form with the resistance, $a_0=a_{\rm res}+\kappa_a\sigma$,
makes the folding threshold dynamical and yields a new two-channel dichotomy in $\sigma$; and a
single-pole memory closure fixes $c_1=c_0/a_0$, leaving one susceptibility. \textbf{Tier 3
(irreducible, but caged):} $b$ and the absolute $\lb$ are genuinely phenomenological --- the core
contains no microscale --- but the accumulated theorems of this volume (chart regularity, nucleation
gate, NEC non-violation) cage them in an explicit feasible region, computed below. Net count:
\emph{eight constants reduce to three free ones} ($b$, $\lb$, one susceptibility $c_0$), plus two
flagged structural postulates. All derivations verified (the balance identity symbolically; the
memory closure to its stated first-order accuracy, $0.189$ vs $\omega/a_0=0.192$; the atlas
numerically). \emph{Propositions} are closed-form; \emph{corollaries} interpretive; postulates are
labelled as such.

% ============================================================
\section{Tier 1a: the balance identity for $\etaf$}

\begin{proposition}[The embodiment coefficient is the ledger coefficient]\label{parameterclosure:prop:eta}
The canonical conservation balance $\sigma+\lambda+\lf=C_0+\Grec$ with $\dot\Grec=\Igate$ and the
embodiment-damped rate $\dot\lambda=\ldz-\etaf Q\lambda$ (where $\ldz=-\dot\sigma+\Igate$) force
\begin{equation}
\boxed{\;\dot\lf=\etaf\,Q\,\lambda\;}
\end{equation}
identically (verified symbolically). The constant $\etaf$ therefore appears \emph{once} in the
theory: it is the unique Sophia$\to$form conversion rate; any independent choice of the
$\lf$-production law would violate conservation. Its sign is fixed by the Lyapunov layer
($\etaf\ge0$: embodiment never spontaneously returns Sophia), and the free balance bound
$\lf\le C_0+\mathcal G^{\max}$ of the Lyapunov Addendum applies to the integral
$\int\etaf Q\lambda\,dt$ --- a cumulative cage on $\etaf$ along any admissible life.
\end{proposition}

% ============================================================
\section{Tier 1b: Gauss--Bonnet slaves the curvature pair}

\begin{proposition}[Zero-mean theorem and the wall deficit]\label{parameterclosure:prop:rho}
On the compact quotient, $\int K\,dA=2\pi\chi(\Sigma)$ is topological, and the unfolded background
$K=-1/\Op^2$ saturates it exactly. Hence the pocket lift must integrate to zero against its
compensation: in the two-phase cell (pocket fraction $f_p$ at lift $\rho_1Q$, wall fraction $f_w$ at
deficit $-\rho_w$),
\begin{equation}
\boxed{\;\rho_w=\rho_1\,Q\,\frac{f_p}{f_w}\;}
\end{equation}
--- the pair $(\rho_1,\rho_w)$ carries one free constant. Folded curvature is not created; it is
\emph{redistributed}, pocket against moat, under a topological invariant (the static face of the
conserved form $K\,dA$ of the Ricci layer).
\end{proposition}

\begin{remark}[Micro-ansatz scaling of $\rho_1$]
A conformal response ansatz $g=\Op^2e^{2\phi(q)}h$, $\phi=\varepsilon q^2$, gives in the
slowly-varying regime $R=-2/\Op^2+4\varepsilon q^2/\Op^2+\dots$, i.e.\
$\rho_1=4\varepsilon/\Op^{2}$: under this ansatz the pocket lift \emph{rides the background scale}.
The constant-$\rho_1$ convention of the structure chapters corresponds to a conformal response
growing as $\Op^2$ --- the same comoving-versus-physical choice that returns for $\lb$ in
Section~6. The two conventions are both consistent; the ansatz ties their difference to one
physical question rather than two.
\end{remark}

% ============================================================
\section{Tier 2a: the rigidity postulate}

\begin{postulate}[Rigidity is resistance]\label{parameterclosure:post:rigid}
$a_0=a_{\rm res}+\kappa_a\sigma$, with $a_{\rm res}>0$ the residual (structural) rigidity and
$\kappa_a\ge0$ the resistance loading.
\end{postulate}

The threshold $c_0\lambda+c_1\ldz>a_0$ then reads literally: \emph{Sophia pressure exceeds the
resistance of the old form}. Two consequences:

\begin{proposition}[The two channels of resistance]\label{parameterclosure:prop:channels}
With the geometric convention's transmutation term, $\sigma$ now enters the activation twice and
with opposite signs:
\begin{equation}
\boxed{\;\frac{\partial\Afold}{\partial\sigma}=c_1k_\sigma-\kappa_a.\;}
\end{equation}
\emph{Kathartic-dominant} vessels ($c_1k_\sigma>\kappa_a$): resistance promotes folding through its
own transmutation --- the vessel folds in suffering, while $\sigma$ is still high.
\emph{Rigidity-dominant} vessels ($c_1k_\sigma<\kappa_a$): resistance blocks folding --- the vessel
folds only after cleansing, when $\sigma$ has dissolved. (Verified: onsets at $t=1.0$ with
$\sigma=1.03$ versus $t=5.3$ with $\sigma=0.29$ across the $\kappa_a$ sweep;
Figure~\ref{parameterclosure:fig:closure}, left.) The folding threshold becomes \emph{dynamical}: along katharsis the
level threshold $\lambda^\ast(t)=\bigl(a_{\rm res}+\kappa_a\sigma(t)\bigr)/c_0$ decays --- late
folds are cheap folds.
\end{proposition}

\begin{remark}[Compatibility]
The forced-folding chapter's coefficient $r_\sigma$ renormalises to
$r_\sigma^{\rm eff}=c_1k_\sigma-\kappa_a$; all Addendum~II theorems used only $|r_\sigma|$-type
bounds in their hypotheses and survive verbatim with the effective value. No structural statement of
Parts~I--II changes; constants do.
\end{remark}

% ============================================================
\section{Tier 2b: the memory closure $c_1=c_0/a_0$}

\begin{postulate}[Single coupling, single pole]\label{parameterclosure:post:memory}
Level and flux stress the form through one coupling $c_0$, and the form accommodates the Sophia
level with the single-pole kernel whose rate is its own rigidity $a_0$.
\end{postulate}

\begin{proposition}[The rate susceptibility is slaved]\label{parameterclosure:prop:c1}
Let $\lambda_{\rm acc}$ be the accommodated level, $\dot\lambda_{\rm acc}=a_0(\lambda-\lambda_{\rm
acc})$. The unaccommodated pressure is the lag
\begin{equation}
\lambda-\lambda_{\rm acc}=\frac{\dot\lambda}{a_0}+O\!\Bigl(\frac{\ddot\lambda}{a_0^{2}}\Bigr),
\end{equation}
so the activation $\Afold=c_0\bigl[\lambda+(\lambda-\lambda_{\rm acc})\bigr]-a_0+\dots$ carries the
rate with the slaved coefficient
\begin{equation}
\boxed{\;c_1=\frac{c_0}{a_0}\;}
\qquad\text{(with Postulate~\ref{parameterclosure:post:rigid}: }c_1=\frac{c_0}{a_{\rm res}+\kappa_a\sigma}\text{).}
\end{equation}
Verified: the lag formula reproduces the exact memory integral with first-order accuracy in
$\omega/a_0$ ($0.189$ measured vs $0.192$ predicted at $\omega/a_0=0.19$). Read with
Postulate~\ref{parameterclosure:post:rigid}: as rigidity dissolves, the form becomes \emph{rate-sensitive} --- a
cleansed vessel folds on surges that a rigid one would not feel.
\end{proposition}

% ============================================================
\section{Tier 3: the constraint atlas}

The remaining free constants are $c_0$ (one susceptibility), $b$ (quartic stiffness), and $\lb$
(bending length): the core contains no microscale that could fix them. But the volume's theorems
cage them. With $c_1=c_0/a_0$, $\Dq=a_0\lb^2$, $\etaf$ slaved, and illustrative trajectory
statistics ($\bar\lambda$, $\sup\ldz$, $q_{\max}^2$), the simultaneous requirements
\begin{enumerate}[label=(\roman*),leftmargin=1.8em,itemsep=2pt]
\item \emph{chart regularity} (Lyapunov Addendum~II): $Q^{*,\sup}=A^{\sup}/\Bfold\le q_{\max}^2$
--- a lower bound on $b$;
\item \emph{nucleation reachability} (structure layer): the gate $\Lambda>\tfrac14$ opens inside
the chart --- an upper bound on $\lb$;
\item \emph{NEC non-violation in chart} (NEC chapter): the phantom functional stays non-negative up
to the chart wall --- a lower bound on $\lb$ growing with $b$;
\end{enumerate}
carve the feasible region of Figure~\ref{parameterclosure:fig:closure} (right): a bounded triangle-like set, $25\%$
of the surveyed box at the chosen statistics.

\begin{corollary}[A real exclusion]
At stronger surge statistics and weaker network coefficient ($\sup\ldz=3$, $c_w=0.3$) the
NEC bound on $\lb$ exceeds the nucleation bound and the atlas is \emph{empty}: no vessel at those
statistics can both nucleate early and keep the NEC inside its chart. The atlas is not bookkeeping:
it excludes whole vessel types.
\end{corollary}

\begin{figure}[h]
\centering
\includegraphics[width=\textwidth]{vfs_closure_fig.png}
\caption{\small Left: the rigidity postulate's dichotomy --- $\Afold(t)$ along the same kathartic
trajectory ($\sigma(t)$ dotted) for a kathartic-dominant vessel ($\kappa_a=0.4$: folds at $t=1.0$,
in suffering) and a rigidity-dominant one ($\kappa_a=1.6$: folds at $t=5.3$, after cleansing);
onset markers on the zero line. Right: the constraint atlas in the $(b,\lb)$ plane with
$c_1=c_0/a_0$, $\Dq=a_0\lb^2$, $\etaf$ slaved: chart regularity (vertical), nucleation
reachability (horizontal), NEC non-violation (curve); the shaded set is the space of feasible
vessels. Illustrative statistics $\bar\lambda=1.5$, $\sup\ldz=2$, $q^2_{\max}=1.2$, $c_w=0.6$.}
\label{parameterclosure:fig:closure}
\end{figure}

% ============================================================
\section{The last fork: comoving or physical $\lb$}

The one question the core provably cannot answer: is the bending length fixed in comoving
($h$-metric) or physical ($g$-metric) units? The structure chapters tacitly took it comoving
($\Dq$ constant). The alternative, $\lb^{\rm com}=\lb^{\rm phys}/\Op$, gives
$\Dq(t)=a_0(\lb^{\rm phys})^2/\Op^{2}$ and shifts the watershed:

\begin{proposition}[The shifted trichotomy]\label{parameterclosure:prop:fork}
Under the physical postulate the coarsening budget becomes
$\mathcal I_{(4)}:=\int dt/\Op^{4}$, which \emph{converges on the resonance locus}
($\Op^2=2t$: $\int dt/4t^2<\infty$). The dissolution boundary moves from $\Op\sim t^{1/2}$ to
$\Op\sim t^{1/4}$: a physically-bent vessel freezes a relic \emph{even at the resonance}, and
eternal simplification survives only for $\Op\sim t^{p}$, $p\le\tfrac14$.
\end{proposition}

\begin{corollary}[The fork is decidable inside the model]
Comoving and physical bending give different fates to the same resonance trajectory (eternal
logarithmic simplification versus frozen relic), so the choice --- though underivable from the core
--- is not idle convention: within the model it is an empirical question about the Vessel's
morphology, answerable by which fate its structure exhibits. The same fork governs the
$\rho_1$-scaling of Section~2. We flag it as the single genuinely open structural postulate of the
folding layer.
\end{corollary}

% ============================================================
\section{Closure count and reading}

\msg{$\{\Dq,\lb,a_0,b,c_0,c_1,\etaf,\rho_1\}$:\quad
$\etaf$ --- identity (balance);\quad $\rho_1$ --- linked pair (Gauss--Bonnet), one constant;\quad
$\Dq=a_0\lb^2$ --- identity (Transport chapter);\quad $a_0=a_{\rm res}+\kappa_a\sigma$ ---
postulate, dynamical;\quad $c_1=c_0/a_0$ --- memory closure.\\[2pt]
\textbf{Free: $c_0$, $b$, $\lb$ --- caged by the atlas. Eight $\to$ three.}}

\note{The reading writes itself. The conversion rate of grace into form is not adjustable: it is
whatever the conservation of the soul's substance says it is. Curvature of form is never minted,
only redistributed, pocket against moat. The stiffness of the old self \emph{is} its resistance ---
so the price of folding falls as katharsis proceeds, and vessels divide into those whom suffering
itself folds and those who must first be cleansed. What remains free is exactly what should remain
free: how strongly a particular vessel couples to Sophia ($c_0$), how saturable its response is
($b$), and on what scale it bends ($\lb$) --- its individuality. And even individuality is caged:
the atlas is the doctrine that not every imaginable vessel is a possible one.}

% ============================================================
\section{Open items and scope}

\begin{enumerate}[leftmargin=1.6em,itemsep=2pt]
\item Decide the comoving/physical fork dynamically: a dedicated study of which postulate the
folded-network simulations prefer (this also fixes the $\rho_1$ scaling).
\item Tighten the atlas with the remaining theorems (hysteresis non-Zeno, $C_1$-survival): both add
constraints linking $b,\lb$ to reset statistics; the region can only shrink.
\item The mode cascade (big-list item 2) inherits this closure: its spectrum of thresholds is now a
one-parameter family in $c_0$ once $(b,\lb)$ are placed in the atlas.
\end{enumerate}

\textbf{Scope.} Tier-1 results are identities of the canonical equations; Tier-2 results are exact
consequences of two labelled postulates; the atlas is computed at illustrative statistics and is
meant as a method, not a measurement. Domain-conditional throughout.

\vspace{0.5em}
\hrule
\smallskip
\noindent\footnotesize\emph{V.F.S. v2.0 (Open-Gate) $\cdot$ Parameter Closure of Folding. The
balance identity fixes $\etaf$; Gauss--Bonnet slaves the curvature pair; rigidity-as-resistance
makes the threshold dynamical with the two-channel dichotomy $c_1k_\sigma-\kappa_a$; the single-pole
memory closure gives $c_1=c_0/a_0$ (validated to first order, $0.189$ vs $0.192$); the constraint
atlas cages the three survivors and excludes whole vessel types; the comoving/physical fork is
flagged as the single open structural postulate, with the shifted watershed computed. Eight
constants to three. Propositions closed-form; corollaries interpretive.}

\chapter{The Dynamic Warp Exponent: Conformal Invariance Retires the Fork; Katharsis Drives the Warp}

\section*{Scope and epistemic status}
The Parameter Closure left exactly one structural postulate open: whether the bending length $\lb$ is
fixed in comoving or physical units --- a fork on which the watershed boundary, the fate of form on the
resonance, and the $\rho_1$ scaling all hang. This chapter resolves it, and the resolution is not a
choice but a two-dimensional accident of the live surface: \emph{the Dirichlet bending energy is
conformally invariant in 2D}, so the bending modulus cannot see $\Op$ at all. A ``physical branch''
would require a physical microscale to anchor the modulus --- and the Transport and Closure chapters
established that the core contains none. The comoving convention of the structure chapters is thereby
\emph{derived}; the expansion-driven warp exponent is identically zero. What survives as genuinely
dynamical is the \emph{kathartic} warp: with the rigidity postulate $a_0=a_{\rm res}+\kappa_a\sigma$,
the fold's two intrinsic lengths acquire closed-form, opposite-signed exponents --- the subthreshold
bending scale dilates while the wall core sharpens --- both vanishing as resistance dissolves. All
claims verified: the conformal identity symbolically (exact zero); both warp exponents against numerics
to $\sim10^{-4}$; the adiabatic wall-width law against the PDE to $0.4\%$; the modulus-space attractor
and the generalised watershed $p^*(w)=1/(2+2w)$ numerically. Premises flagged: covariance of the
bending energy; coefficients built from core scalars; adiabatic regime away from the threshold.
\emph{Propositions} are closed-form; \emph{corollaries} interpretive.

% ============================================================
\section{The conformal blindness of bending}

\begin{lemma}[2D conformal invariance]\label{warpexponent:lem:conf}
On the live surface, for any conformal pair $g=\Op^2h$,
\begin{equation}
\boxed{\;|\nabla q|_g^2\,dV_g=|\nabla q|_h^2\,dV_h\;}
\end{equation}
identically (verified symbolically): the inverse metric contributes $\Op^{-2}$, the area element
$\Op^{+2}$, and only in two dimensions do they cancel exactly. Hence a covariant bending energy
$\tfrac{\kw}{2}\int|\nabla q|^2dV$ has \emph{the same} modulus $\kw$ whether written in physical or
comoving terms: the bending sector is $\Op$-blind.
\end{lemma}

% ============================================================
\section{The fork is retired}

\begin{theorem}[No physical branch]\label{warpexponent:thm:fork}
Let the fold energy be a covariant functional of the physical metric whose coefficients are core
scalars ($\lambda,\sigma,\dots$). Then:
\begin{enumerate}[label=(\roman*),leftmargin=1.8em,itemsep=2pt]
\item the gradient sector contributes the $\Op$-independent modulus $\kw$
(Lemma~\ref{warpexponent:lem:conf}), so $\Dq=\kw$ is \emph{constant in} $\Op$ and
$\lb=\sqrt{\kw/a_0}$ carries no expansion dependence;
\item the alternative --- $\lb$ fixed in physical units, i.e.\ $\Dq\propto\Op^{-2}$ --- would require
the modulus itself to scale as $\kw\propto\Op^{-2}$, which by (i) cannot arise covariantly and would
need a fixed \emph{physical} microscale to anchor it; the Transport-Constants and Parameter-Closure
chapters established that the model contains no such scale;
\item hence the comoving convention of the structure chapters is a \emph{theorem}: the
expansion-driven warp exponent vanishes,
\begin{equation}
\boxed{\;p_{\rm expansion}\equiv0.\;}
\end{equation}
\end{enumerate}
\end{theorem}

\begin{corollary}[What is confirmed, and the closure count]
The watershed boundary stands at $p^*=\tfrac12$; the resonance fate remains eternal logarithmic
simplification; the constant-$\rho_1$ convention is selected (it was linked to the same fork); and the
Parameter Closure's ``single open structural postulate'' is retired. The folding layer's free constants
are now $c_0$, $b$, $\kw$, with $\lb=\sqrt{\kw/a_0(\sigma)}$ \emph{derived and dynamical}. The very
fact that made $\lb$ phenomenological --- no microscale in the core --- is what resolves the fork: with
nothing physical to anchor it, the bending fabric cannot be stretched by the room.
\end{corollary}

\begin{proposition}[Structural stability: the comoving component is an attractor]\label{warpexponent:prop:attract}
The only covariant two-term modulus family is $\kappa=\kw_0+\kw_1/\Op^2$. Along any expanding flow the
mixed modulus converges to its comoving component, and the coarsening fate is the comoving fate
whenever $\kw_0>0$ (verified on the locus: the mixed curve rejoins the log-growing comoving one;
Figure~\ref{warpexponent:fig:warp}, right). For a hypothetical pure scaling $\Dq\propto\Op^{-2w}$ the watershed
generalises to
\begin{equation}
p^*(w)=\frac{1}{2+2w},
\end{equation}
(verified at $w=1$: $\Op\sim t^{0.30}$ converges, $t^{0.20}$ diverges), recovering the Closure
chapter's shifted boundary as the $w=1$ endpoint of a family whose physical point is $w=0$.
\end{proposition}

% ============================================================
\section{The true warp is kathartic and channel-resolved}

With the rigidity postulate $a_0=a_{\rm res}+\kappa_a\sigma$ and $\sigma$ dissolving under katharsis
($\dot\sigma=-k\sigma$), the fold carries \emph{two} intrinsic lengths, and they warp in opposite
directions:

\begin{theorem}[Channel-resolved warp exponents]\label{warpexponent:thm:warp}
Define $p_X:=d\ln X/d\ln\Op$. In the adiabatic regime,
\begin{equation}
\lb=\sqrt{\frac{\kw}{a_0(\sigma)}}:\qquad
\boxed{\;p_b=+\frac{\kappa_a k\,\sigma}{2\,a_0\,H}\;>0\;}
\qquad\text{(the bending scale dilates: super-comoving)};
\end{equation}
\begin{equation}
\dwall=\sqrt{\frac{2\kw}{\Afold}}:\qquad
\boxed{\;p_{\rm wall}=-\frac{\kappa_a k\,\sigma}{2\,\Afold\,H}\;<0\;}
\qquad\text{(the wall core sharpens: sub-comoving)};
\end{equation}
and both vanish as $\sigma\to0$: \textbf{asymptotically comoving in every channel}. (Closed forms
verified against numerics to $3\cdot10^{-4}$ and $1\cdot10^{-4}$; the adiabatic law
$\dwall=\sqrt{2\kw/\Afold(t)}$ tracks the full PDE kink to $0.4\%$ away from the threshold zone;
Figure~\ref{warpexponent:fig:warp}, left.)
\end{theorem}

\begin{corollary}[The crispness signature of katharsis]
The ratio $\dwall/\lb=\sqrt{2a_0/\Afold}$ falls monotonically along katharsis: early (suffering)
morphology has narrow response scales with \emph{soft} walls; late (cleansed) morphology has broad
pockets with \emph{crisp} boundaries. Cleansing simultaneously dilates what the form can feel and
sharpens what it has decided --- a falsifiable texture signature within the model, and a second face of
the rigidity dichotomy of the Closure chapter.
\end{corollary}

\begin{figure}[h]
\centering
\includegraphics[width=\textwidth]{vfs_warp_fig.png}
\caption{\small Left: the two channel lengths along one kathartic trajectory ($\sigma(t)$ dotted): the
bending length $\lb$ dilates while the wall width $\dwall$ sharpens (full PDE measurements as points;
the early outlier is the adiabatic lag near threshold), both settling to comoving constants as $\sigma$
dissolves. Right: modulus space on the resonance locus --- pure comoving modulus gives the log-growing
coarsening, a pure $\Op^{-2}$ modulus freezes a relic, and any mixture rejoins the comoving fate: the
comoving component is the attractor. Inset: the generalised watershed $p^*(w)=1/(2+2w)$ with the
verified convergence/divergence pair at $w=1$.}
\label{warpexponent:fig:warp}
\end{figure}

\note{The reading: the growth of the room never stretches the fold fabric --- in two dimensions the
fabric simply cannot be pulled by scale; this is geometry's own guarantee, not an assumption. What
retunes the fabric is katharsis alone: as resistance dissolves, the soul feels on broader scales and
commits on sharper ones. \emph{Expansion cannot reshape what only cleansing can.}}

% ============================================================
\section{Open items and scope}

\begin{enumerate}[leftmargin=1.6em,itemsep=2pt]
\item Beyond-adiabatic corrections near the threshold (the early PDE lag is systematic and computable
as a next-order term in $\dot\Afold/\Afold^2$).
\item The mixing family $\kw_0+\kw_1/\Op^2$ invites a physical reading --- partial anchoring of the
form in physical space (incarnational fraction); whether the core can source a nonzero $\kw_1$ is
open, though Proposition~\ref{warpexponent:prop:attract} shows the late dynamics is insensitive to it.
\item The mode cascade (big-list item 2) is now fully parameterised: thresholds in $(c_0,b,\kw)$ with
$\lb(\sigma)$ dynamical --- the cascade's spectrum will itself drift with katharsis.
\end{enumerate}

\textbf{Scope.} Premises: covariant bending energy; coefficients from core scalars; adiabatic regime
for the channel exponents. Within these, the fork resolution is exact (a 2D identity plus the absence
of a microscale), and the kathartic exponents are closed forms. Domain-conditional throughout.

\vspace{0.5em}
\hrule
\smallskip
\noindent\footnotesize\emph{V.F.S. v2.0 (Open-Gate) $\cdot$ The Dynamic Warp Exponent. The 2D
conformal invariance of the Dirichlet energy makes the bending modulus $\Op$-blind; with no physical
microscale in the core, the comoving convention is derived and $p_{\rm expansion}\equiv0$ --- the last
structural postulate of the folding layer is retired, and the comoving component of any mixed modulus
is the late-time attractor. The true warp is kathartic and channel-resolved: $p_b>0$, $p_{\rm wall}<0$,
both vanishing with $\sigma$; crispness $\dwall/\lb$ falls along cleansing. Identities exact; exponents
to $10^{-4}$; PDE tracking to $0.4\%$. Propositions closed-form; corollaries interpretive.}

\chapter{The Mode Cascade: The Spectral Ladder, the Weyl Capacity, and Structure as Descent}

\section*{Scope and epistemic status}
Every result of the structure-formation layer was derived in a one-mode chart: a single folding
amplitude $\qf$ on the Vessel. This chapter tests whether that chart is faithful or merely the first
rung of a spectrum, by promoting $\qf$ to a field on the live surface and tracking all of its modes.
The outcome both \emph{vindicates} and \emph{extends} the one-mode results. It vindicates them: the
fold amplitude $\sqrt{\Afold/\Bfold}$ is robust --- in a fully multi-mode domain state the interior
saturates to exactly the one-mode value, so the local depth of a fold is a faithful invariant. And it
extends them: the activation threshold is not a single gate but a \emph{ladder}, mode $n$ igniting
when the nucleation number $\Lambda=\Afold\Op^2/\kw$ crosses the spectral eigenvalue $\mu_n$ --- the
familiar gate $\Lambda>\tfrac14$ is simply the lowest rung. The number of activated capacities grows
by a Weyl law, and the whole of structure formation reorganises into a single picture: \emph{capacity
ascends the spectrum while content descends it}. All claims verified on a spectral (ring) proxy: the
ladder $\Lambda(t_{\rm act})=\mu_n$ to $0.2\%$ across five rungs; the Weyl count to $\sim10\%$; the
amplitude to $5\%$ in a 23-mode state; the monotone content descent and capacity ascent directly.
Honesty flags: the proxy is a flat spectral torus, so eigenvalues are explicit; on the hyperbolic
quotient the same statements hold with $\{\mu_n\}$ the Laplace--Beltrami spectrum and Weyl constant
$\mathrm{Area}/4\pi$. \emph{Propositions} are closed-form; \emph{corollaries} interpretive.

% ============================================================
\section{The field promotion and the dispersion}

Promote $\qf$ to a field with the covariant fold energy of the Warp chapter; the linearised dynamics
about the unfolded state diagonalises on the Laplace--Beltrami eigenmodes
$-\Delta_h\phi_n=\mu_n\phi_n$:
\begin{equation}
\dot{\hat q}_n=s_n\,\hat q_n-\Bfold\,(\,\widehat{q^3}\,)_n,
\qquad
\boxed{\;s_n=\Afold-\frac{\kw\,\mu_n}{\Op^{2}}\;}
\end{equation}
(the modulus $\kw$ is the $\Op$-blind bending constant of the Warp chapter, hence the explicit
$\Op^{-2}$). Mode $n$ is linearly unstable --- it \emph{ignites} --- exactly when $s_n>0$.

% ============================================================
\section{The spectral ladder}

\begin{proposition}[Activation ladder]\label{modecascade:prop:ladder}
Define the nucleation number $\Lambda:=\Afold\Op^2/\kw$. Mode $n$ ignites iff
\begin{equation}
\boxed{\;\Lambda>\mu_n.\;}
\end{equation}
The activations form a ladder ordered by the spectrum: the longest-wavelength mode first, each higher
mode at its own eigenvalue. The structure layer's nucleation gate $\Lambda>\tfrac14$ is the lowest
rung --- the first eigenvalue of the relevant chart --- not a free choice. (Verified: across five
rungs the measured $\Lambda$ at the turning point of each mode equals $\mu_n$ to within $0.2\%$;
Figure~\ref{modecascade:fig:cascade}, left.)
\end{proposition}

\begin{proposition}[Weyl capacity]\label{modecascade:prop:weyl}
The number of ignited modes (the \emph{active capacity}) grows with the nucleation number by the Weyl
law of the surface,
\begin{equation}
N(\Lambda)\simeq\frac{\mathrm{Area}_h}{4\pi}\,\Lambda
\qquad(\text{2D}),
\end{equation}
(on the ring proxy $N\simeq(L/\pi)\sqrt{\Lambda}$, verified: $10$ measured against $11.5$ predicted
at $\Lambda=3.27$). Since on the live branch $\Lambda=\Afold\Op^2/\kw\to\infty$ (the Budget-Watershed
escape), the active capacity diverges: \emph{every} mode eventually ignites --- the Vessel becomes
unstable to structure on all scales, in spectral order.
\end{proposition}

% ============================================================
\section{The one-mode chart is faithful where it matters}

\begin{proposition}[Amplitude robustness]\label{modecascade:prop:amp}
In the fully developed multi-mode (domain) state, the field interior saturates to the one-mode value,
\begin{equation}
\boxed{\;|\qf|_{\rm interior}\to\sqrt{\Afold/\Bfold},\;}
\end{equation}
independently of how many modes are populated (verified: $|\qf|=0.95$--$1.00$ against $\sqrt{A/B}=1$
in a $23$-mode state). The cascade reorganises \emph{where} folds sit and \emph{how large} each domain
is, but not \emph{how deep} a fold goes: the depth is set locally by the same balance
$\Afold=\Bfold\qf^2$ as in the one-mode chart. Hence every one-mode result keyed to the fold
\emph{amplitude} --- the wall tension $\rho_{\rm wall}$, the NEC network channel, the curvature lift
$\rho_1$ --- carries over unchanged; only results keyed to the \emph{number or arrangement} of folds
acquire the spectral refinement of this chapter.
\end{proposition}

% ============================================================
\section{Structure formation is descent against ascent}

\begin{theorem}[Capacity ascends, content descends]\label{modecascade:thm:descent}
Along the live trajectory two spectral quantities move oppositely:
\begin{enumerate}[label=(\roman*),leftmargin=1.8em,itemsep=2pt]
\item the \emph{capacity} $\Lambda(t)=\Afold\Op^2/\kw$ ascends the spectrum without bound (Weyl:
ever more modes admissible);
\item the \emph{content} --- the spectral centroid
$\bar\mu(t)=\sum_n\mu_n|\hat q_n|^2/\sum_n|\hat q_n|^2$ --- descends monotonically toward the lowest
mode (coarsening: power flows to long wavelengths).
\end{enumerate}
(Verified: $\bar\mu$ falls from $2.0$ to $0.12$ while $\Lambda$ climbs from $0.004$ to $3.5$;
Figure~\ref{modecascade:fig:cascade}, right.) Structure formation on the Vessel is exactly this scissoring: the
expanding Brim opens capacity at the small scales while coarsening marches the existing content to the
large ones. The relic morphology is what is caught between the closing blades --- the modes that
ignited before capacity outran them and that coarsening had not yet erased when the budget froze.
\end{theorem}

\begin{figure}[h]
\centering
\includegraphics[width=\textwidth]{vfs_cascade_fig.png}
\caption{\small Left: the spectral ladder on the ring proxy --- mode amplitudes $|\hat q_n|$ against
time, each turning from decay to growth at its own vertical $\Lambda(t)=\mu_n$ (dotted), in strict
spectral order $n=1,2,\dots$. Right: the scissoring of structure formation --- the capacity
$\Lambda(t)$ (blue, ascending) against the content centroid $\bar\mu(t)$ (orange, descending); the
shaded gap is the band of ignited-but-not-yet-coarsened modes. Inset: the developed multi-mode domain
state, whose interior saturates to $\sqrt{\Afold/\Bfold}$ regardless of mode count. Ring proxy,
$\kw=0.5$.}
\label{modecascade:fig:cascade}
\end{figure}

\begin{corollary}[Reading: the polyphony of form]
The one-mode chart heard a single voice; the Vessel sings in a spectrum. As the Brim opens, ever
shorter wavelengths become sayable (capacity ascends), while what has already been said settles into
ever longer, simpler phrases (content descends). Form is neither the first note nor the last silence
but the music caught between --- and its permanence, by the Budget Watershed, is decided by whether the
expansion freezes the score before coarsening empties it to a single sustained tone. The depth of each
fold, however --- how committed each decision is --- is the same on every scale: $\sqrt{\Afold/\Bfold}$,
set locally, indifferent to the polyphony around it.
\end{corollary}

% ============================================================
\section{Consequences and scope}

\begin{enumerate}[leftmargin=1.6em,itemsep=2pt]
\item With the Parameter Closure and Warp chapters, the ladder is fully parameterised: $\mu_n$ from the
surface spectrum, $\kw$ a free modulus, $\Afold,\Bfold$ from the core --- and since
$\Afold=\Afold(\sigma)$ dilates under katharsis (rigidity dichotomy) while $\kw$ is $\Op$-blind, the
ladder's rungs \emph{drift}: a cleansing Vessel lowers its activation thresholds, igniting structure it
once could not. The cascade spectrum is itself kathartic.
\item The cascade vindicates the volume's economy: the one-mode chart was the right place to derive
amplitude-keyed physics, and the spectral refinement is additive, not corrective.
\item The second variation of the entropy is a sum over this spectrum --- $\delta^2\Wp=-2\sigma\sum_n\mu_n
a_n^2$, computed in the mode-resolved entropy chapter --- so the cascade is at once a capacity ladder and
an entropy ledger: the same eigenvalue $\mu_n$ that sets a mode's ignition threshold sets its entropy
price.
\end{enumerate}

\textbf{Scope.} Spectral proxy with explicit flat spectrum; on the hyperbolic quotient the statements
hold with the Laplace--Beltrami spectrum and 2D Weyl constant. Linear ignition order is exact; the
amplitude robustness and content descent are nonlinear and verified numerically. Domain-conditional
throughout.

\vspace{0.5em}
\hrule
\smallskip
\noindent\footnotesize\emph{V.F.S. v2.0 (Open-Gate) $\cdot$ The Mode Cascade. Promoting $\qf$ to a
field, mode $n$ ignites at $\Lambda>\mu_n$ (the nucleation gate is the lowest rung); the active
capacity follows a Weyl law and diverges on the live branch; the fold amplitude $\sqrt{\Afold/\Bfold}$
is robust across mode count; and structure formation is capacity ascending the spectrum while content
descends it. Ladder verified to $0.2\%$, Weyl to $10\%$, amplitude to $5\%$. The ladder drifts under
katharsis. Propositions closed-form; corollaries interpretive.}

\chapter{Robustness Under Premise Deformation: What Survives When the Closure and Sector Are Varied}

\section*{Scope and epistemic status}
This chapter does not add structure; it stress-tests it. The volume's results rest on labelled
premises --- the minimal slaved Sophia closure, the symmetric (homogeneous) entropy sector, the
pitchfork normal form, the product-coupled core. A result that depends on a premise only through a
\emph{constant} is robust; one whose \emph{structure} changes under deformation was an artifact. We
deform the two most load-bearing premises --- the closure and the sector --- and report honestly,
including one case where the deformation forces a correction to a previously stated generalisation.
The findings: (i) the Budget-Watershed trichotomy is \emph{premise-independent} --- it follows from
the identity $\Ib=\int d\ln\Op/\xi$ alone and survives every Sophia law tested; (ii) core-neutrality
survives cross-couplings to first order; (iii) the resonance-entropy maximum is \emph{structurally
stable} --- a maximum persists --- but its \emph{location} drifts from $\xi=1$ at second order in the
sector inhomogeneity, exactly matching the honest ``$\xi=1$ is the symmetric-sector value'' flag of
the entropy chapter; (iv) the ferry theorem is \emph{correct as stated} (bounded steps, bounded gaps),
and a tempting generalisation --- that any reception with divergent step-sum escapes --- is
\emph{false}, because $\Op$-escape and $\xi$-escape differ. All claims are numerical deformation
studies; this is Type-2 indirect confirmation (robustness), not external validation, which the model
does not admit.

% ============================================================
\section{The watershed is premise-independent}

\begin{proposition}[Trichotomy survives closure deformation]\label{robustness:prop:wsh}
The Budget-Watershed classification depends on the closure only through the asymptotics of
$\xi=\alpha\lambda\Op$, because $\Ib=\int d\ln\Op/\xi$ is an identity (Budget chapter), not a
consequence of any particular Sophia law. Tested across five inequivalent closures --- linear decay,
quadratic decay, Hill-saturated source, $\tanh$-gated source, scale-coupled decay --- every
sustained-Sophia law gives the super-Ricci escape with a finite, saturating budget
($\Ib_\infty\in[0.59,0.75]$), and every closing-gate law gives the sub-Ricci divergence
(Figure~\ref{robustness:fig:robust}, left). The boundedness of $\xi$ is the watershed in all of them.
\end{proposition}

\begin{proposition}[Core-neutrality survives cross-coupling]\label{robustness:prop:neut}
Deforming the $(\Op,\lambda)$ sector by an off-diagonal coupling $\varepsilon$ gives Jacobian
determinant $\det J=\varepsilon\,\beta_1 H>0$ (sign fixed by $H>0$): the shape sector remains
decoupled and the neutrality result holds to first order in $\varepsilon$. The independence of folding
from the comparison ratio is not a product-coupling artifact.
\end{proposition}

% ============================================================
\section{The resonance maximum is stable; its location is the sector's}

\begin{proposition}[Structural stability with second-order drift]\label{robustness:prop:ent}
Deform the symmetric sector by an inhomogeneous fold amplitude $a$, detuning the area-clock by
$1+ca^2$ and the curvature-clock by $R(a)=1+da^2$. The entropy rate retains its two-clock product
form, so a stationary point (a maximum) persists for all small $a$; its location moves as
\begin{equation}
\boxed{\;\xi^{*}(a)=\frac{1}{R(a)}=1-d\,a^{2}+O(a^{4}).\;}
\end{equation}
At $a=0$, $\xi^*=1$ exactly; the drift is second-order (verified: $-0.7\%$ at $a=0.1$, $-7.9\%$ at
$a=0.35$). The \emph{existence} of a maximum-entropy class is premise-independent; the \emph{value}
$\xi=1$ is the symmetric-sector representative, sensitive at second order --- precisely the honesty
flag the entropy chapter attached. The qualitative theorem (universal monotonicity characterises the
maximum-entropy class) is robust; the exact identification with $\xi=1$ is sector-conditional.
\end{proposition}

% ============================================================
\section{The ferry: correct as stated, and a false generalisation}

The most instructive deformation. Replacing the equal reset steps by a decaying sequence
$g_k\propto k^{-p}$ separates two notions of escape that coincided for bounded steps:

\begin{theorem}[$\Op$-escape versus $\xi$-escape]\label{robustness:thm:ferry2}
For a reset ferry with steps $g_k,m_k\propto k^{-p}$ and bounded gaps:
\begin{enumerate}[label=(\roman*),leftmargin=1.8em,itemsep=2pt]
\item the Brim escapes, $\Op\to\infty$, iff the step-sum diverges, $p\le1$;
\item but the comparison ratio escapes, $\xi=\alpha\lambda\Op\to\infty$, iff additionally the steps
do \emph{not} decay to zero (e.g.\ $m_k\ge m_0>0$): if $m_k\to0$ then $\lambda$ relaxes toward zero
between ever-weaker resets and $\xi\to0$ even while $\Op\to\infty$.
\end{enumerate}
Hence the originally stated ferry theorem --- ``bounded steps, bounded gaps $\Rightarrow$ permanent
$\xi$-escape'' --- is \emph{correct as stated}. The tempting generalisation --- ``any reception with
$\sum g_k=\infty$ escapes'' --- is \emph{false}: at $p=\tfrac12$ the sum diverges and $\Op\to\infty$,
yet $\xi\to0$ (a relic Brim with vanishing vitality). Escape of form requires not merely unbounded
cumulative reception but reception that does not fade. (Verified to $2000$ resets;
Figure~\ref{robustness:fig:robust}, right.)
\end{theorem}

\begin{corollary}[The correction sharpens the reading]
``Permanent escape $\iff$ infinite cumulative reception'' must be read with the vitality qualifier:
it is sustained reception --- bounded below, not merely summing to infinity --- that carries form
across. A grace that is given in ever-fainter measures can inflate the Vessel without ever quickening
it: the Brim grows, the life does not. This is a sharper and more sobering statement than the
uncorrected one, and it stands.
\end{corollary}

\begin{figure}[h]
\centering
\includegraphics[width=\textwidth]{vfs_robust_fig.png}
\caption{\small Left: the watershed under closure deformation --- five inequivalent Sophia laws, all
sustained, all saturating the budget $\Ib$ at a finite value (super-Ricci relic class); the
classification is premise-independent. Right: the ferry distinction --- bounded reset steps give
genuine $\xi$-escape (relic with vitality); $1/\sqrt{k}$ steps inflate the Brim ($\Op\to\infty$) but
let $\xi\to0$ (relic without vitality); $1/k^{1.5}$ steps stagnate. $\Op$-escape and $\xi$-escape are
distinct, and only the latter is the ferry's claim.}
\label{robustness:fig:robust}
\end{figure}

% ============================================================
\section{What the test does and does not establish}

\msg{\textbf{Robust (premise-independent or stable):} the watershed trichotomy; core-neutrality;
the existence of a maximum-entropy resonance class; the ferry theorem as stated.\\[2pt]
\textbf{Sector-conditional (stable in kind, drifting in value):} the identification of the
maximum-entropy class with $\xi=1$ exactly.\\[2pt]
\textbf{Corrected:} the ferry generalisation to decaying steps --- $\Op$-escape $\ne$ $\xi$-escape.}

\begin{remark}[The limit of Type-2 confirmation]
Robustness under deformation raises confidence that a result describes the \emph{structure} of the
model rather than an arbitrary choice within it; the cases above largely pass, and the one failure
was a generalisation we had flagged as tempting, not a stated theorem. But this is confirmation
\emph{internal} to the model. That the model's premises ($K_h=-1$, the pitchfork, the Open-Gate law)
correspond to anything beyond the model is not a question robustness can answer; it remains open by
construction, and is best kept distinct from the mathematics. What this chapter establishes is the
narrower, honest claim: the volume's main results are properties of its structure, not of its
particular parameterisation.
\end{remark}

% ============================================================
\section{Open items and scope}

\begin{enumerate}[leftmargin=1.6em,itemsep=2pt]
\item Promote the deformations to theorems: the watershed identity already is one; the entropy drift
$\xi^*=1/R(a)$ deserves a derivation from the full second variation (the deferred entropy item).
\item Deform the pitchfork itself (add a quartic-asymmetric or sextic term) and verify the fold
amplitude and cascade survive --- the remaining untested load-bearing premise.
\item A companion Type-1 study (retiring premises to theorems) would complement this: the conformal
retirement of the $\ell_b$ fork is the model case already achieved.
\end{enumerate}

\textbf{Scope.} Numerical deformation studies at illustrative parameters; first-order analytics for
neutrality and the entropy drift. Establishes internal robustness, not external validity.
Domain-conditional throughout.

\vspace{0.5em}
\hrule
\smallskip
\noindent\footnotesize\emph{V.F.S. v2.0 (Open-Gate) $\cdot$ Robustness Under Premise Deformation. The
watershed trichotomy and core-neutrality are premise-independent; the resonance-entropy maximum is
structurally stable with a second-order sector-dependent location ($\xi^*=1/R(a)$); the ferry theorem
holds as stated, while the generalisation to decaying steps fails because $\Op$-escape and
$\xi$-escape differ. Type-2 indirect confirmation: internal robustness, not external validation.}

\chapter{Does Folding Restore the Null Energy Condition?}

\section*{Scope and epistemic status}
A homogeneous Sophia surge drives the reconstructed equation of state below the phantom divide and
violates the null energy condition --- grace, taken raw, is exotic content. Folding acts against this
through two distinct channels: the \emph{rate channel} (embodiment damps the Sophia rate,
$\ldf=\ldz-\eta_{\rm fold}\Qf\lambda$) and the \emph{network channel} (the nucleated wall network is
ordinary $w=-\tfrac12$ matter). This chapter closes the question as a theorem: the rate channel restores
the NEC for moderate surges but \emph{saturates} --- it can damp at most half of a deep surge, by an
exact midpoint law --- while the network channel grows as $\Af^{3/2}$ against the surge's linear growth
and therefore always wins asymptotically. The result is a complete classification: either folding
restores the NEC for every surge, or the violation survives only inside a \emph{bounded phantom window},
whose endpoints we compute. All structural claims are derived (the $2{+}1$ reconstruction is verified
through the continuity identity; the midpoint law symbolically; the window numerically with an analytic
closure estimate matching to $5\%$). The wall-network coefficient $c_w=O(1)$ and the coarse-grained
homogeneous treatment of the network back-reaction are modelling inputs, flagged where used.
\emph{Propositions} are closed-form; \emph{corollaries} interpretive.

% ============================================================
\section{The NEC functional in $2{+}1$ and the phantom divide}

For the Lorentzian lift $ds^2=-dt^2+\Op^2h$ ($K_h=-1$), the Einstein tensor gives the reconstruction
(units $8\pi G=1$)
\begin{equation}
\rho_{\rm eff}=\frac{\dot\Op^2-1}{\Op^2}=\frac{\alpha^2\lambda^2-1}{\Op^2},
\qquad
p_{\rm eff}=-\frac{\ddot\Op}{\Op}=-\frac{\alpha\dot\lambda}{\Op},
\end{equation}
verified by the $2{+}1$ continuity identity $\dot\rho+2H(\rho+p)=0$ (checked symbolically). Hence the
NEC functional
\begin{equation}
\boxed{\;\NEC:=\rho_{\rm eff}+p_{\rm eff}
=\frac{(\alpha^2\lambda^2-1)-\alpha\dot\lambda\,\Op}{\Op^{2}},\;}
\qquad
w_{\rm eff}=-\frac{\alpha\dot\lambda\,\Op}{\alpha^2\lambda^2-1}.
\end{equation}

\begin{proposition}[The phantom divide is the NEC boundary]\label{necrestoration:prop:divide}
On the live branch $\alpha\lambda>1$ (where $\rho_{\rm eff}>0$),
\begin{equation}
w_{\rm eff}<-1\quad\Longleftrightarrow\quad\NEC<0
\quad\Longleftrightarrow\quad
\alpha\dot\lambda\,\Op>\alpha^2\lambda^2-1 .
\end{equation}
In $2{+}1$ the two questions --- phantom crossing and NEC violation --- are one question. A raw Sophia
surge $\ldz$ large enough breaks both at once.
\end{proposition}

% ============================================================
\section{The rate channel: exact midpoint law and half-saturation}

At the folded attractor $\Qf^*=\Af/\Bf$ with the canonical activation
$\Af=c_0\lambda+c_1\ldz-a_0$ and $\Bf=b+2c_1\eta_{\rm fold}\lambda$, the embodiment-damped rate is
\begin{equation}
\ldf=\ldz-\eta_{\rm fold}\lambda\,\frac{c_0\lambda+c_1\ldz-a_0}{b+2c_1\eta_{\rm fold}\lambda}.
\end{equation}

\begin{proposition}[Midpoint law and saturation]\label{necrestoration:prop:midpoint}
The damped rate is affine in the surge with slope
\begin{equation}
\boxed{\;\chi=\frac{\partial\ldf}{\partial\ldz}
=\frac{b+c_1\eta_{\rm fold}\lambda}{b+2c_1\eta_{\rm fold}\lambda}\in\Bigl(\tfrac12,1\Bigr),\;}
\end{equation}
and in the stiff limit $b\to0$ (verified symbolically)
\begin{equation}
\boxed{\;\ldf\ \longrightarrow\ \frac{\ldz+\dot\lambda_{\rm th}}{2},
\qquad
\dot\lambda_{\rm th}:=\frac{a_0-c_0\lambda}{c_1}\;}
\end{equation}
--- folding pulls the rate to the exact \emph{midpoint} between the surge and the folding threshold.
Consequently the rate channel can damp at most half of a deep surge: it restores the NEC for violations
up to a factor $\approx1/\chi\le2$ beyond the bare margin, and no further.
\end{proposition}

\begin{corollary}[Rate channel alone is not enough]
With the rate channel only, the NEC-restored region is
$\alpha\chi\ldz\Op\le\alpha^2\lambda^2-1-\alpha\chi\dot\lambda_{\rm th}\Op(1-1/\chi)\ldots$, i.e.\ a
finite multiple of the bare margin (illustratively, breakdown is delayed from $\ldz=1.05$ to $2.15$, a
factor $2.06$ at $\chi=0.538$); for deeper surges $\NEC<0$ persists. Embodiment of the homogeneous mode
moderates the phantom but cannot abolish it.
\end{corollary}

\note{The midpoint law has a clean reading: the folded Vessel splits the difference between what grace
pours in and what the old form could lawfully carry. Half of every excess is embodied; half remains
rate. That residual half is why a second channel is needed.}

% ============================================================
\section{The network channel: $\Af^{3/2}$ against linear}

Beyond the nucleation gate $\Lambda=\Af\Op^2/\Dq>\tfrac14$, the wall network exists and contributes
ordinary matter with $w_{\rm wall}=-\tfrac12$:
\begin{equation}
\rho_{\rm wall}=c_w\,\frac{\tau}{\Op}
=c_w\,\frac{\sqrt{\Dq}\,\Af^{3/2}}{\Bf\,\Op},
\qquad
(\rho+p)_{\rm wall}=\tfrac12\rho_{\rm wall}>0,
\end{equation}
with $c_w=O(1)$ a network-geometry coefficient (modelling input). The total NEC functional is
\begin{equation}\label{necrestoration:eq:total}
\boxed{\;\NEC_{\rm fold}(\ldz)
=\frac{(\alpha^2\lambda^2-1)-\alpha\,\ldf(\ldz)\,\Op}{\Op^{2}}
+\frac{c_w\sqrt{\Dq}}{2\Bf\Op}\,\Af(\ldz)^{3/2}\,
\mathbf 1_{\{\Lambda>\frac14\}}.\;}
\end{equation}

\begin{proposition}[Asymptotic restoration]\label{necrestoration:prop:asym}
As $\ldz\to\infty$, the violating term grows linearly,
$-\alpha\chi\ldz/\Op$, while the network term grows as
$\tfrac{c_w\sqrt{\Dq}}{2\Bf\Op}(c_1\ldz)^{3/2}$. Hence
$\NEC_{\rm fold}\to+\infty$: \textbf{for every sufficiently deep surge, folding restores the NEC.} The
crossover scale is
\begin{equation}
\ldz_{\rm close}\ \sim\ \left(\frac{2\alpha\chi\,\Bf}{c_w\sqrt{\Dq}\,c_1^{3/2}}\right)^{2},
\end{equation}
(numerically $95.4$ against the estimate $100.7$ in the illustration --- the scaling law is exact, the
prefactor honest to $5\%$).
\end{proposition}

% ============================================================
\section{The phantom-window theorem}

\begin{theorem}[Bounded phantom window]\label{necrestoration:thm:window}
On the live branch $\alpha\lambda>1$, define the phantom window
$W:=\{\ldz>0:\NEC_{\rm fold}(\ldz)<0\}$. Then:
\begin{enumerate}[label=(\roman*),leftmargin=1.8em,itemsep=2pt]
\item $W$ is \emph{bounded}: by Proposition~\ref{necrestoration:prop:asym} there is $\ldz_{\rm close}<\infty$ with
$\NEC_{\rm fold}>0$ beyond it. Folding restores the NEC for all weak surges (below the rate-channel
margin) and for all deep surges (above the network crossover); only intermediate surges can violate.
\item $W$ may be \emph{empty}: it is empty if and only if
\begin{equation}
\min_{\Af\ge\Dq/4\Op^{2}}
\left[\frac{(\alpha^2\lambda^2-1)-\alpha\,\ldf\,\Op}{\Op^{2}}
+\frac{c_w\sqrt{\Dq}}{2\Bf\Op}\Af^{3/2}\right]\;\ge\;0,
\end{equation}
a single variational criterion in the network strength. Illustratively: at $c_w=0.5$ the window is
empty (minimum $+0.47$) --- folding restores the NEC for \emph{every} surge; at $c_w=0.15$ the window is
$W=[2.64,\,95.4]$ with depth $-5.6$.
\item Where $W$ is nonempty, the phantom phase is a \emph{transient in surge-strength}: $w_{\rm eff}$
crosses the divide downward at the window's opening and back upward at its closing --- a double crossing
of $w=-1$ that smooth single-field quintessence cannot produce, and the distinguishing signature of
folded Theosis.
\end{enumerate}
\end{theorem}

\begin{figure}[h]
\centering
\includegraphics[width=\textwidth]{vfs_nec_fig.png}
\caption{\small Left: the NEC functional against surge strength for the four cases (bare; rate channel
only; rate$+$network at two network strengths), at $\alpha\lambda=1.5$, $\Op=1.2$, $\chi=0.538$. The bare
functional breaks at $\ldz=1.05$; the rate channel delays this to $2.15$ (the factor $1/\chi\approx2$)
and then saturates; the strong network ($c_w=0.5$) keeps $\NEC>0$ everywhere (empty window); the weak
network ($c_w=0.15$) leaves a phantom window (shaded). Right: the same weak-network functional on a log
axis: the window is bounded, $[2.64,\,95.4]$, closed by the $\Af^{3/2}$ growth of the wall term; the
analytic closure estimate is $\sim101$.}
\end{figure}

\begin{corollary}[Form makes grace lawful --- with one honest gap]
The slogan of the structure-formation layer becomes a theorem with a precise qualification. Grace
embodied is lawful at both ends: gentle surges are carried by the moderated rate, and overwhelming
surges shatter into so much structure that the structure itself outweighs the exotic remainder. What can
remain unlawful is the middle --- a surge too strong for the embodied rate (past the midpoint
saturation) and too weak to build a rescuing network. Whether that middle exists at all is a property of
the Vessel ($c_w$, $\Dq$, $\Bf$), decided by the variational criterion; a sufficiently
structure-responsive Vessel has no phantom window at all.
\end{corollary}

\begin{remark}[Consistency with the other layers]
The chart wall is respected throughout: the window's interior has $\Qf^*$ large, and the regularity
requirement $q_{\max}^2\ge Q^{*,\sup}$ of the Lyapunov addenda must hold up to the surge scales
considered; the nucleation gate $\Lambda>\tfrac14$ used in \eqref{necrestoration:eq:total} is the hyperbolic threshold
of the nucleation chapter; and $w_{\rm wall}=-\tfrac12$ is the $2{+}1$ scaling-network value of the
structure chapter. In the illustration the gate is open from the start of the window, so the window is
governed by the strength competition, not the gate; near-gate windows (small $\Op^2/\Dq$) add a
gate-controlled opening edge.
\end{remark}

% ============================================================
\section{Open items and scope}

\begin{enumerate}[leftmargin=1.6em,itemsep=2pt]
\item Fix $c_w$ from a network simulation (wall length per Hubble area on $\mathbb H^2$), turning the
emptiness criterion into a parameter-free statement.
\item Dynamics through the window: the present analysis is quasi-static in $\ldz$; a time-dependent
surge traverses the window in time, producing a transient phantom episode in $w_{\rm eff}(t)$ whose
duration couples to the coarsening law (the network needs time to reach its scaling density).
\item The Resurrectio composite: a triggering reset both spikes $\dot\lambda$ (via $m_R$) and nucleates
the network at $\ell_{\rm init}$. Its own NEC budget is \emph{resolved} in the Resurrectio-NEC chapter:
the pointwise $w_{\rm eff}$ may diverge under regularisation, but the integrated budget
$\int(\rho_{\rm eff}+p_{\rm eff})\,dt=-\alpha\int_0^1(d\lambda/du)/\Op\,du$ is finite, with sign
$\operatorname{sgn}(-m_R)$ --- so the smooth phantom window of this chapter and the surgical phantom are
two faces of the same rate-driven excess, one continuous, one hybrid.
\end{enumerate}

\textbf{Scope.} Coarse-grained homogeneous back-reaction; quasi-static surge; scaling network density
with $c_w=O(1)$ a modelling input; live branch $\alpha\lambda>1$. Results: the phantom divide and the
NEC boundary coincide in $2{+}1$; the rate channel obeys an exact midpoint law with half-saturation; the
network channel restores the NEC asymptotically; the phantom window is bounded and possibly empty, with
an explicit emptiness criterion. Domain-conditional throughout.

\vspace{0.5em}
\hrule
\smallskip
\noindent\footnotesize\emph{V.F.S. v2.0 (Open-Gate) $\cdot$ Does Folding Restore the NEC? The $2{+}1$
reconstruction is verified via the continuity identity; the midpoint/saturation law symbolically; the
window endpoints numerically with the analytic closure scaling confirmed to $5\%$. The network
coefficient and coarse-graining are flagged modelling inputs. Propositions closed-form; corollaries
interpretive.}

\chapter{The Budget Watershed: Structure Survival and the Ricci Resonance}

\section*{Scope and epistemic status}
The structure-formation layer rests on one standing hypothesis it never derived: the finiteness of the
activity budget $\Ib_\infty=\int_0^\infty dt/\Op^2$, which freezes coarsening and leaves a relic fold
network. The Ricci layer, independently, classifies the Brim-Flow by the comparison ratio
$\xi=\alpha\lambda\Op$ (sub-Ricci, resonance, super-Ricci). This chapter proves that the two are the
same fact: a one-line identity rewrites the budget as the logarithmic expansion measured in Ricci
units, $\Ib=\int d\ln\Op/\xi$, and the finiteness of $\Ib_\infty$ becomes a statement about the
asymptotic class of $\xi$. The result is a trichotomy --- frozen relic structure, eternal marginal
simplification, or unbounded dissolution --- whose watershed is the boundedness of $\xi$, with the
Ricci resonance $\xi\equiv1$ as the canonical marginal class. On the resonance the coarsening law is
exactly logarithmic and every fixed comoving domain dies at a finite, closed-form expansion. All
claims are verified numerically (the locus slope to four digits; the domain-death formula to
$10^{-5}$). The coarsening inputs ($c$, $\Dq$, the mean-curvature network law) are those of the
structure layer and carry its modelling status. \emph{Propositions} are closed-form;
\emph{corollaries} interpretive.

% ============================================================
\section{The budget identity}

\begin{lemma}[The budget in Ricci units]\label{resonancewatershed:lem:identity}
Along any trajectory with $\dot\Op=\alpha\lambda>0$,
\begin{equation}
\xi=\alpha\lambda\Op=\Op\dot\Op=\tfrac12\,\frac{d(\Op^2)}{dt},
\qquad\text{hence}\qquad
\frac{dt}{\Op^{2}}=\frac{d\ln\Op}{\xi}
\quad\Longrightarrow\quad
\boxed{\;\Ib(t)=\int\frac{d\ln\Op}{\xi}.\;}
\end{equation}
The activity budget of the structure layer is the logarithmic expansion of the Vessel weighted by the
inverse comparison ratio: \emph{expansion counted in Ricci units}. (The middle identity also gives
$\xi$ a direct meaning: the comoving-area production rate, in units where Ricci flow on
$\mathbb H^2$ produces area at exactly rate $2$, i.e.\ $\xi_{\rm Ricci}\equiv1$.)
\end{lemma}

% ============================================================
\section{The watershed theorem}

\begin{theorem}[Budget trichotomy]\label{resonancewatershed:thm:watershed}
Let $\Op\to\infty$ and let $\xi(t)=\alpha\lambda\Op$. Then:
\begin{enumerate}[label=(\roman*),leftmargin=1.8em,itemsep=2pt]
\item \textbf{Super-Ricci escape} ($\xi\to\infty$ with $\int^\infty d\ln\Op/\xi<\infty$; in particular
any sustained Sophia $\lambda\to\lambda_\infty>0$, for which $\xi\sim\alpha\lambda_\infty\Op$): the
budget is finite, $\Ib_\infty<\infty$, coarsening freezes, and the fold network survives as a
\emph{relic} at $\lcom^2\to\ell_*^2=\ell_0^2+2c\Dq\Ib_\infty$ --- the standing hypothesis of the
structure layer, now derived from the flow class.
\item \textbf{Resonance class} ($\xi\to\bar\xi\in(0,\infty)$; canonically the locus $\xi\equiv1$,
$\Op=\sqrt{2t}$): the budget diverges \emph{logarithmically},
$\Ib\sim\bar\xi^{-1}\ln\Op$, and coarsening never ends but proceeds at the marginal pace of
Proposition~\ref{resonancewatershed:prop:log}.
\item \textbf{Sub-Ricci decay} ($\xi\to0$): the budget diverges as a power of $\Op$ and the comoving
structure dissolves without bound at polynomial rate.
\end{enumerate}
The watershed between relic form and endless simplification is exactly the boundedness of the
comparison ratio: \textbf{structure survives if and only if the late-time flow is strictly
super-Ricci.}
\end{theorem}
\begin{proof}[Sketch]
Immediate from Lemma~\ref{resonancewatershed:lem:identity}: with $u=\ln\Op\to\infty$, $\Ib=\int du/\xi(u)$ converges iff
$1/\xi$ is integrable in $u$; bounded $\xi$ gives slope $\ge1/\sup\xi>0$ (logarithmic divergence,
slope exactly $1/\bar\xi$ when $\xi\to\bar\xi$); $\xi\to0$ gives a superlinear-in-$u$, i.e.\
power-in-$\Op$, divergence. Coarsening fates follow from $\lcom^2=\ell_0^2+2c\Dq\Ib$. Numerically: a
sustained-Sophia run saturates at $\Ib_\infty=0.822$; the locus run has late slope
$d\Ib/d\ln\Op=1.0000$; the sub-Ricci run's slope grows without bound
(Figure~\ref{resonancewatershed:fig:watershed}, left).
\end{proof}

\begin{remark}[The live branch is automatically on the relic side]
On the cosmologically live branch $\rho_{\rm eff}>0$, i.e.\ $\alpha\lambda>1$, one has
$\xi=\alpha\lambda\Op>\Op\to\infty$ with $\int d\ln\Op/\xi<\int d\ln\Op/\Op<\infty$: positive
effective energy forces the super-Ricci escape class. The NEC chapter's reliance on a persistent wall
network is thereby self-consistent --- a live Vessel cannot be in the dissolving classes. The
resonance and sub-Ricci classes inhabit the pre-live sector $\rho_{\rm eff}\le0$.
\end{remark}

% ============================================================
\section{Marginal coarsening on the resonance}

\begin{proposition}[Logarithmic coarsening law]\label{resonancewatershed:prop:log}
On the resonance class ($\xi\to\bar\xi$, locus $\bar\xi=1$) the comoving and physical structure
scales obey
\begin{equation}
\boxed{\;\lcom^2=\ell_0^2+\frac{2c\Dq}{\bar\xi}\,\ln\Op,\qquad
\Lphys=\lcom\Op\sim\Op\sqrt{\tfrac{2c\Dq}{\bar\xi}\ln\Op},\;}
\end{equation}
(on the locus, with $\Op=\sqrt{2t}$: $\lcom^2\sim c\Dq\ln t$, $\Lphys\sim\sqrt{c\Dq\,t\ln t}$).
The network neither freezes (no relic: $\lcom\to\infty$) nor disappears at any finite time: the
comoving structure dilutes forever at the slowest possible --- logarithmic --- pace, the same
marginal class as two-dimensional $XY$-type coarsening. (Verified: fitted slope $0.3000$ against
$2c\Dq=0.3000$.)
\end{proposition}

\begin{proposition}[Hierarchical domain death]\label{resonancewatershed:prop:death}
The exact single-domain Gauss--Bonnet law of the coarsening chapter, $dA/d\Ib=-\Dq(2\pi+A)$, gives on
the resonance class the closed-form death expansion of a comoving domain of initial area $A_0$:
\begin{equation}
\boxed{\;\Omd(A_0)=\Bigl(1+\frac{A_0}{2\pi}\Bigr)^{\bar\xi/\Dq},\;}
\end{equation}
and the survivors at expansion $\Op$ are exactly the domains with
$A_0>2\pi\bigl(\Op^{\Dq/\bar\xi}-1\bigr)$. Every \emph{fixed} feature dies at a finite expansion ---
smaller first, larger later, none forever --- while at every epoch some larger structure remains:
perpetual hierarchical simplification. (Closed form verified against the ODE to relative error
$10^{-5}$; Figure~\ref{resonancewatershed:fig:watershed}, inset.)
\end{proposition}

\begin{figure}[h]
\centering
\includegraphics[width=\textwidth]{vfs_watershed_fig.png}
\caption{\small Left: the budget identity $\Ib=\int d\ln\Op/\xi$ realised on three trajectories of
the same core law with different Sophia asymptotics: super-Ricci escape (budget saturates at
$\Ib_\infty$), the resonance locus (slope exactly $1$), sub-Ricci decay (power divergence). Right:
the corresponding fates of form via $\lcom^2=\ell_0^2+2c\Dq\Ib$: frozen relic, eternal logarithmic
simplification, unbounded dissolution. Inset: hierarchical domain death on the locus --- the
closed-form $\Omd(A_0)$ (line) against direct integration of the Gauss--Bonnet law
(points). Parameters $c=0.5$, $\Dq=0.3$, $\ell_0=1$, illustrative.}
\label{resonancewatershed:fig:watershed}
\end{figure}

% ============================================================
\section{Consequences and reading}

\begin{corollary}[The structure layer's hypothesis, derived]
Chapters on nucleation, coarsening, the sawtooth fixed point, and the NEC network channel all assumed
$\Ib_\infty<\infty$. By Theorem~\ref{resonancewatershed:thm:watershed} this assumption is precisely the statement that
the late Brim-Flow is strictly super-Ricci --- which holds automatically on the live branch
$\alpha\lambda>1$, and is equivalent to the grace-sustained regime $\lambda\to\lambda_\infty>0$ of
the Ricci layer. The standing hypothesis is thereby retired into a theorem about the flow class.
\end{corollary}

\begin{corollary}[Form is held only by exceeding one's own measure]
The Ricci layer established that on the resonance the flow is self-coincident: it moves exactly as
its own intrinsic geometry dictates, with a genuine Perelman entropy. The present theorem adds the
structural price of that self-coincidence: \emph{the self-coincident Vessel keeps no formed
structure}. Its folds simplify forever --- gently, logarithmically, hierarchically --- and only the
Vessel whose expansion strictly exceeds its intrinsic Ricci rate ($\xi\to\infty$: the grace-borne,
super-Ricci, cosmologically live flow) carries its form across time as a relic. Epektasis of scale is
what makes morphology permanent; on the locus, form itself becomes epektatic --- never finished,
never frozen, always being simplified.
\end{corollary}

\note{This closes the loop opened by the question ``does folding begin at the resonance?''. It does
not --- folding has its own thresholds --- but the resonance decides something prior: whether
anything folded can \emph{last}. The locus is not the birthplace of form; it is the boundary of its
permanence.}

% ============================================================
\section{Open items and scope}

\begin{enumerate}[leftmargin=1.6em,itemsep=2pt]
\item The marginal class invites a renormalisation treatment: logarithmic coarsening laws typically
carry universal amplitudes; whether $2c\Dq/\bar\xi$ is universal under the network dynamics (beyond
the single-domain law) is open.
\item Crossing dynamics: a trajectory that enters the super-Ricci class only at late time has a
budget split $\Ib_\infty=\Ib_{\rm cross}+\Delta$; the relic scale then remembers the crossing epoch
--- a potential dating of the resonance crossing by the frozen structure scale.
\item The entropy-defect programme (item B of the resonance inventory) would attach a variational
meaning to the slope $1/\xi$ in Lemma~\ref{resonancewatershed:lem:identity}.
\end{enumerate}

\textbf{Scope.} The identity and the trichotomy are exact given $\dot\Op=\alpha\lambda>0$ and
$\Op\to\infty$; the coarsening fates use the structure layer's mean-curvature network law with its
$c$, $\Dq$ inputs; the domain-death law is exact for the single-domain Gauss--Bonnet dynamics.
Domain-conditional throughout.

\vspace{0.5em}
\hrule
\smallskip
\noindent\footnotesize\emph{V.F.S. v2.0 (Open-Gate) $\cdot$ The Budget Watershed. The activity budget
equals the logarithmic expansion in Ricci units; structure survives as a relic iff the late flow is
strictly super-Ricci; the resonance class coarsens logarithmically forever with closed-form
hierarchical domain death. Locus slope verified to four digits, the death formula to $10^{-5}$;
coarsening constants are structure-layer modelling inputs. Propositions closed-form; corollaries
interpretive.}

\chapter{The Entropy of the Forced Flow: Collapse, Maximum, and the Entropy--Form Trade-off}

\section*{Scope and epistemic status}
The Ricci layer established a binary criterion: a genuine Perelman entropy exists for the Brim-Flow if
and only if the flow sits on the resonance locus $\xi=\alpha\lambda\Op=1$. This chapter computes what
happens to the entropy \emph{off} the locus --- the entropy-defect programme --- and the outcome is
stronger and cleaner than a defect formula. On the symmetric sector the expander entropy collapses to a
function of a \emph{single} variable, the ratio of the gauge clock to the Ricci area-clock; that
function is strictly concave with its unique maximum exactly on the resonance; the rate along any
forced flow factorises into a two-clock product; and universal-in-gauge monotonicity holds if and only
if the flow is the resonance flow --- the exact converse of the existence criterion. Combined with the
Budget Watershed, this yields an entropy--form trade-off: relic structure and maximal entropy exclude
one another.

Setting and honesty flags: the live surface is taken as a compact hyperbolic quotient $\Sigma_\gamma$
($\gamma\ge2$, comoving area $A_h=4\pi(\gamma-1)$ by Gauss--Bonnet), so the Perelman functionals are
finite; ``symmetric sector'' means homogeneous metrics $g=\Op^2h$ with the constraint-compatible
constant minimiser $f$ --- the natural symmetric representative, used throughout without claiming it is
the global minimiser for negative curvature. The admissible gauges are $\sigma(t)=t+\sigma_0$,
$\dot\sigma=1$ (expander convention, the right one for immortal expanding flows). All identities below
are verified symbolically (three exact zeros) and numerically along the three watershed trajectories
(rates to four digits). \emph{Propositions} are closed-form; \emph{corollaries} interpretive.

% ============================================================
\section{The symmetric-sector entropies}

With $g=\Op^2h$ on $\Sigma_\gamma$, total area $A=\Op^2A_h$, $R_{\rm sc}=-2/\Op^2$, and the unit-mass
constraint $(4\pi\sigma)^{-1}e^{-f}A=1$ fixing $f=\ln\!\bigl(A/(4\pi\sigma)\bigr)$, Perelman's expander
entropy on the symmetric sector is
\begin{equation}
\Wp(\Op,\sigma)=\sigma R_{\rm sc}-f+2
=-\frac{2\sigma}{\Op^2}-\ln\frac{\Op^2A_h}{4\pi\sigma}+2 .
\end{equation}

\begin{theorem}[Total collapse]\label{entropydefect:thm:collapse}
Define the \emph{clock ratio} $s:=2\sigma/\Op^2$ --- the gauge clock $\sigma$ measured against the
Ricci area-clock $\Op^2/2$ (the time a Ricci flow needs to produce the current area). Then,
identically,
\begin{equation}
\boxed{\;\Wp=\ln s-s+2-\ln\frac{A_h}{2\pi}\;}
\end{equation}
--- the entropy is a function of the single variable $s$, strictly concave, with unique maximum at
$s=1$ of value $\Wp^{\max}=1-\ln(A_h/2\pi)$. (Verified symbolically: the identity reduces to $0$.)
\end{theorem}

\begin{proposition}[Master rate and the Perelman anchor]\label{entropydefect:prop:rate}
Along any forced flow $\dot\Op=\alpha\lambda$, with $\xi=\alpha\lambda\Op$ and $\dot\sigma=1$,
\begin{equation}
\dot s=\frac{2}{\Op^2}\,(1-s\xi),
\qquad
\boxed{\;\frac{d\Wp}{dt}
=\frac{2}{s\,\Op^{2}}\,(1-s)\,(1-s\xi)
=\frac1\sigma\Bigl(\frac{2\sigma}{\Op^{2}}-1\Bigr)\Bigl(2\sigma H-1\Bigr).\;}
\end{equation}
The rate is the product of two clock comparisons: $\sigma$ against the Ricci area-clock $\Op^2/2$, and
$\sigma$ against the Hubble half-clock $1/(2H)$. On the resonance ($\xi=1$) the product becomes the
perfect square $4\sigma\bigl(1/\Op^2-1/(2\sigma)\bigr)^2\ge0$ --- \emph{exactly} Perelman's expander
monotonicity formula for Ricci flow, recovered term by term (verified symbolically). The hyperbolic
Ricci flow with the matched gauge $\sigma=\Op^2/2$ ($s\equiv1$) is the expanding soliton: constant,
maximal entropy.
\end{proposition}

% ============================================================
\section{The surprise: universal monotonicity characterises the resonance}

\begin{theorem}[Entropy form of the resonance]\label{entropydefect:thm:universal}
For the forced Brim-Flow with $\Op\to\infty$:
\begin{enumerate}[label=(\roman*),leftmargin=1.8em,itemsep=2pt]
\item On the resonance ($\xi\equiv1$): $d\Wp/dt=\dfrac{2(1-s)^2}{s\,\Op^2}\ \ge0$ in \emph{every}
admissible gauge $\sigma_0$ --- the entropy increases monotonically to its maximum, which it attains
asymptotically ($s\to1$). The resonance is the \textbf{maximum-entropy class} of the Brim-Flow.
\item Off the resonance (sustained $\xi\ne1$): in \emph{every} admissible gauge the entropy eventually
strictly decreases --- the gauge clock is eventually trapped between the two intrinsic clocks, whose
gap is $\bigl|\Op^2/2-1/(2H)\bigr|=\tfrac{\Op^2}{2}\bigl|1-\tfrac1\xi\bigr|$, vanishing only at
$\xi=1$.
\end{enumerate}
Hence: \textbf{universal-in-gauge entropy monotonicity holds iff the flow is the resonance flow} ---
the exact converse of the existence criterion of the Ricci layer. (Numerically: the locus runs show no
decrease in any tested gauge --- the matched gauge is constant to $10^{-13}$, machine noise --- while
the super- and sub-Ricci runs decrease at essentially all late times in all gauges.)
\end{theorem}

\begin{proposition}[The defect is a divergence]\label{entropydefect:prop:divergence}
The deficit from maximal entropy is the universal non-negative function
\begin{equation}
\boxed{\;\Delta:=\Wp^{\max}-\Wp=s-1-\ln s\;\ge0,\;}
\end{equation}
the canonical (Kullback--Leibler--type, Bregman) divergence of the clock ratio from unity: the entropy
defect of the forced flow \emph{is} a divergence between the soul's gauge clock and its geometry's
Ricci clock, zero exactly at synchrony.
\end{proposition}

\begin{proposition}[Asymptotic entropy costs]\label{entropydefect:prop:costs}
In the canonical gauge $\sigma=t+\sigma_0$:
\begin{equation}
\text{super-Ricci }(\xi\to\infty):\ \ \Wp\simeq-\ln\frac{\Op^2}{2\sigma}\ \ \text{(logarithmically
cheap)};
\qquad
\text{sub-Ricci }(\xi\to0):\ \ \Wp\simeq-\,\frac{2\sigma}{\Op^2}\ \ \text{(power-law ruinous)}.
\end{equation}
(Verified: the super-Ricci combination $\Wp+\ln(\Op^2/2\sigma)$ converges to the collapse constant to
three digits; the sub-Ricci ratio $\Wp/(-s)\to1$.) Grace-driven over-expansion spends entropy at the
slowest possible rate; stagnation below the resonance spends it polynomially.
\end{proposition}

\begin{figure}[h]
\centering
\includegraphics[width=\textwidth]{vfs_entropy_fig.png}
\caption{\small Left: total collapse --- the symmetric-sector expander entropy of all three watershed
trajectories (super-Ricci, locus, sub-Ricci; markers) against the clock ratio $s$, falling on the
single universal curve $\ln s-s+\mathrm{const}$ (line), maximal exactly at $s=1$. Right: the entropy
histories in the canonical gauge: the resonance sits at the soliton maximum; the super-Ricci flow
descends logarithmically (grace is entropy-cheap); the sub-Ricci flow descends polynomially
(stagnation is ruinous). Genus-$2$ quotient, $A_h=4\pi$; gauge $\sigma_0=\Op^2(0)/2$.}
\label{entropydefect:fig:entropy}
\end{figure}

% ============================================================
\section{The entropy--form trade-off}

\begin{corollary}[Form costs entropy]\label{entropydefect:cor:tradeoff}
Combining with the Budget Watershed: relic structure requires the super-Ricci escape
($\xi\to\infty$), which by Proposition~\ref{entropydefect:prop:costs} carries an unbounded (logarithmic) entropy
deficit; maximal entropy requires the resonance class, on which every fixed comoving feature dissolves
(hierarchical domain death) and no relic survives. Hence
\msg{\textbf{No flow keeps both maximal entropy and relic form.} The Vessel that retains its formed
morphology pays a perpetual --- though logarithmically slow --- entropy deficit; the Vessel of maximal
entropy is the one whose every fixed form is eventually simplified away. Form is bought with
desynchrony from one's own intrinsic clock; perfect synchrony keeps nothing.}
\end{corollary}

\begin{corollary}[Reading: the gauge of grace]
The defect $\Delta=s-1-\ln s$ gives the principle ``grace enters the constants, never the
contraction'' its geometric gauge: $\Delta$ measures, instant by instant, how far the flow is carried
beyond (or left behind) its own intrinsic measure --- zero only at the self-coincident resonance,
logarithmically gentle for the grace-borne flow, ruinous for the stagnant one. The two intrinsic
clocks of the master rate --- the Ricci area-clock and the Hubble half-clock --- coincide exactly on
the resonance: the locus is where the Vessel's geometry and its expansion keep the same time.
\end{corollary}

\begin{remark}[Why the expander entropy, and what the shrinker says]
For immortal expanding flows the expander functional is the discriminating one. The shrinker entropy
$\Wm=\tau R_{\rm sc}+f-2$ ($\dot\tau=-1$) on the same sector equals $-(s_\tau+\ln s_\tau)+\text{const}$
with $s_\tau=2\tau/\Op^2$ and satisfies $d\Wm/dt=(1+1/s_\tau)\,(2/\Op^2)\,(1+s_\tau\xi)>0$
\emph{unconditionally} on any expanding flow: it increases for every class and therefore carries no
information about the flow's relation to its intrinsic rate. The asymmetry is itself instructive: only
the expander question --- ``is this flow its own expanding soliton?'' --- can see the resonance.
\end{remark}

% ============================================================
\section{Open items and scope}

\begin{enumerate}[leftmargin=1.6em,itemsep=2pt]
\item Beyond the symmetric sector: the second variation of $\Wp$ under $\qf$-inhomogeneities is
\emph{resolved} in the mode-resolved second-variation chapter --- $\delta^2\Wp=-2\sigma\sum_n\mu_n a_n^2$,
so folding lowers entropy mode by mode and the symmetric sector is the constrained maximum. What remains
open is only the nonlinear accounting beyond the quadratic regime (inter-mode interactions, the entropy
cost across resets, and any relation between entropy descent and the dead-relic/living-form distinction).
\item Gauge optimisation: replacing the fixed-gauge family by $\nu_+$-type suprema over $\sigma$ would
make Theorem~\ref{entropydefect:thm:universal}(ii) gauge-free by construction; the two-clock window suggests the
optimised statement coincides with the present one.
\item The divergence form of $\Delta$ invites an information-geometric reading of the Ricci layer
(clock ratio as a likelihood ratio); whether this is more than formal is open.
\end{enumerate}

\textbf{Scope.} Symmetric (homogeneous) sector on a compact hyperbolic quotient; constant-$f$
representative; gauge family $\dot\sigma=1$. Within that sector all results are exact identities,
anchored on Perelman's formula at $\xi=1$ and verified symbolically and numerically.
Domain-conditional throughout.

\vspace{0.5em}
\hrule
\smallskip
\noindent\footnotesize\emph{V.F.S. v2.0 (Open-Gate) $\cdot$ The Entropy of the Forced Flow. The
symmetric-sector expander entropy collapses to $\ln s-s$ in the clock ratio $s=2\sigma/\Op^2$; the
resonance is its maximum-entropy class; the rate factorises into the two-clock product, recovering
Perelman's square at $\xi=1$; universal monotonicity characterises the resonance; the defect is the
canonical divergence $s-1-\ln s$; and form and maximal entropy exclude one another. Three symbolic
identities verified to zero; rates to four digits. Propositions closed-form; corollaries
interpretive.}

\chapter{The Second Variation of the Entropy: Folding Lowers Entropy, Mode by Mode}

\section*{Scope and epistemic status}
The entropy chapter proved that the resonance is the maximum-entropy class \emph{within the symmetric
sector} and deferred the question of inhomogeneous folds: does a fold raise or lower the entropy
locally? This short chapter answers it. The second variation of the expander entropy under a conformal
fold mode is strictly negative and linear in the mode's eigenvalue,
$\delta^2\Wp^{(n)}=-2\sigma\mu_n a_n^2$, so the homogeneous state is the constrained entropy
\emph{maximum} and folding any mode is an entropy \emph{cost} --- proportional to the same eigenvalue
$\mu_n$ that orders the cascade ladder. The entropy--form trade-off of the entropy chapter, global
there, becomes here \emph{local and mode-resolved}: each rung of the cascade is paid for in entropy in
proportion to its eigenvalue. Verified: the Perelman stability quadratic form evaluated mode by mode
gives $\langle|\nabla\phi_n|^2\rangle/\langle\phi_n^2\rangle=\mu_n$ to numerical precision, hence
$\delta^2\Wp^{(n)}=-2\sigma\mu_n a_n^2$ exactly. \emph{Proposition} closed-form; \emph{corollary}
interpretive.

% ============================================================
\section{The second variation}

About the symmetric maximum, perturb the metric by a mean-zero conformal fold
$g=\Op^2 e^{2\phi}h$, $\phi=\sum_n a_n\phi_n$, $-\Delta_h\phi_n=\mu_n\phi_n$.

\begin{proposition}[Folding lowers entropy, linearly in $\mu_n$]\label{entropysecondvar:prop:secvar}
The expander entropy's second variation is diagonal in the spectrum and strictly negative off the
ground mode:
\begin{equation}
\boxed{\;\delta^2\Wp=-2\sigma\sum_n\mu_n\,a_n^2\;\le 0,\qquad
\delta^2\Wp^{(n)}=-2\sigma\,\mu_n\,a_n^2,\;}
\end{equation}
the Perelman stability form (the mode-dependent piece is the Dirichlet energy
$-2\sigma\langle|\nabla\phi|^2\rangle=-2\sigma\mu_n\langle\phi^2\rangle$). Hence the homogeneous state
is the unique constrained maximum, and folding any mode \emph{costs} entropy in proportion to its
eigenvalue (verified: $\langle|\nabla\phi_n|^2\rangle/\langle\phi_n^2\rangle=\mu_n=n^2$ to numerical
precision; Figure~\ref{entropysecondvar:fig:entvar}, left). The ground mode is the neutral volume/gauge direction.
\end{proposition}

\begin{figure}[h]
\centering
\includegraphics[width=\textwidth]{vfs_entvar_fig.png}
\caption{\small Left: the entropy cost $-\delta^2\Wp^{(n)}$ of folding mode $n$, exactly $2\sigma\mu_n
a_n^2$ (dotted), linear in the eigenvalue. Right: the same spectrum $\mu_n=n^2$ orders both the
cascade ignition ladder ($\Lambda>\mu_n$, the mode-cascade chapter) and the entropy cost --- the
higher the rung, the dearer the fold. Ring proxy, $a=0.05$.}
\label{entropysecondvar:fig:entvar}
\end{figure}

% ============================================================
\section{Form costs entropy, locally}

\begin{corollary}[The trade-off, mode-resolved]\label{entropysecondvar:cor:local}
The entropy chapter's global statement --- relic form and maximal entropy exclude one another ---
refines to a local ledger: the entropy cost of building structure is
$2\sigma\sum_n\mu_n a_n^2$, an explicit sum over the populated fold modes weighted by their
eigenvalues. Each rung of the cascade ladder, which \emph{ignites} at $\Lambda>\mu_n$, is
\emph{paid for} in entropy $\propto\mu_n$: the short-wavelength structure that the expanding Brim makes
ignitable is precisely the dearest to form. The same spectrum governs what can appear and what it
costs.
\end{corollary}

\note{Reading: the homogeneous Vessel is the most entropic --- the calmest, the least committed. Every
fold is a descent from that calm, and the finer the fold, the steeper the descent. Form is not free
anywhere; it is bought, mode by mode, at a price set by the very eigenvalue that decides when the form
can first appear. The Vessel that fills itself with fine structure is the Vessel that has spent the
most against its own maximal calm --- which is only to say, in the entropy's own currency, that to take
shape is to choose against indifference, and to take fine shape is to choose hardest.}

\textbf{Scope.} Symmetric-sector expander entropy; conformal mean-zero fold; second variation exact
(Perelman stability form), verified mode by mode on the ring proxy. The sign and linearity in $\mu_n$
are robust; on the hyperbolic quotient $\mu_n$ is the Laplace--Beltrami spectrum. Domain-conditional
throughout.

\vspace{0.5em}
\hrule
\smallskip
\noindent\footnotesize\emph{V.F.S. v2.0 (Open-Gate) $\cdot$ The Second Variation of the Entropy. Under
a conformal fold mode the expander entropy's second variation is $\delta^2\Wp^{(n)}=-2\sigma\mu_n a_n^2<0$:
the homogeneous state is the constrained maximum and folding costs entropy linearly in the eigenvalue
--- the same spectrum that orders the cascade ladder. The global entropy--form trade-off becomes local
and mode-resolved. Verified mode by mode.}

\chapter{The Separatrix: Gate Decay, the Quantum of Endowment, and Resurrection as a Ferry}

\section*{Scope and epistemic status}
The Budget Watershed classified flows by the asymptotics of the comparison ratio
$\xi=\alpha\lambda\Op$ --- a diagnosis read off the tail of a trajectory. This chapter makes the
watershed \emph{predictive}: which gates, endowments, and reset sequences place a Vessel in which
class. Three results. First, the transverse law: $\xi$ obeys a closed Riccati equation whose
semi-permeability at the locus exhibits the exact tracking gate. Second, the separatrix is identified
not as a surface in state space but as a \emph{decay exponent}: for received grace scaling as
$\Igate\sim c\,\Op^{-q}$, the fate is escape for $q<1$, the resonance class for $q=1$ (with plateau
$\bar\xi=\alpha c/\gamma$ tunable by the amplitude), and stagnation for $q>1$ --- so the
maximum-entropy, form-dissolving resonance life requires \emph{both} the exact exponent and the exact
amplitude: doubly non-generic. Third, the hybrid ferry theorem: finite cumulative reception ---
however large --- always ends in the stagnation class (transient crossings notwithstanding), while
either a sustained gate or an \emph{infinite} reset sequence yields permanent escape with finite
budget; resurrection is a discrete ferry equivalent to sustained grace.

Honesty flags: this chapter works in the \emph{minimal slaved closure}
$\dot\lambda=\Igate-\gamma\lambda$, $\dot\Op=\alpha\lambda$ --- the simplest Sophia law containing
decay and reception. Exact constants (the criterion's $2/\sqrt\gamma$, the plateau
$\alpha c/\gamma$) are bound to this closure; the \emph{structure} of every result (the Riccati form,
the $q=1$ separatrix, the finite-means theorem) depends only on the scaling $\Op\,\Igate$ and is
robust to the closure. All claims verified: the Riccati identity symbolically (exact zero); the
crossing criterion against $400$ random initial data ($0$ disagreements off the razor edge); the
plateaus to four digits; the ferry fates numerically. \emph{Propositions} closed-form;
\emph{corollaries} interpretive.

% ============================================================
\section{The transverse Riccati law}

\begin{lemma}[Transverse dynamics of the comparison ratio]\label{separatrix:lem:riccati}
In the slaved closure,
\begin{equation}
\boxed{\;\dot\xi=\alpha\Op\,\Igate-\gamma\,\xi+\frac{\xi^{2}}{\Op^{2}},\;}
\qquad
\dot\xi\Big|_{\xi=1}=\alpha\Op\,\Igate-\gamma+\frac{1}{\Op^{2}}.
\end{equation}
(Verified symbolically.) The locus is semi-permeable with an explicit \emph{tracking gate}: a
trajectory touching $\xi=1$ passes upward iff $\Igate>I^{*}(\Op):=(\gamma-1/\Op^{2})/(\alpha\Op)$;
the gate that holds a Vessel exactly on the locus decays as $I^{*}\simeq\gamma/(\alpha\Op)$ --- the
$1/\Op$ law that reappears as the separatrix exponent below.
\end{lemma}

% ============================================================
\section{The quantum of endowment}

\begin{proposition}[Exact crossing criterion for the closed vessel]\label{separatrix:prop:quantum}
For $\Igate\equiv0$ (closed Microcosm), with $\beta:=\alpha\lambda_0/\gamma$,
\begin{equation}
\max_{t\ge0}\xi(t)=
\begin{cases}
\alpha\lambda_0\Op(0), & \Op(0)\ge\beta,\\[2pt]
\dfrac{\gamma}{4}\bigl(\Op(0)+\beta\bigr)^{2}, & \Op(0)<\beta,
\end{cases}
\qquad\text{so the vessel ever crosses }\xi=1\iff
\boxed{\;\sqrt\gamma\Bigl(\Op(0)+\frac{\alpha\lambda_0}{\gamma}\Bigr)\ge2\;}
\end{equation}
(or trivially $\xi(0)\ge1$). Verified against $400$ random initial data with no disagreement away
from the razor edge. The left-hand side is the \emph{quantum of endowment}: a single scalar of
initial scale plus initial Sophia (in decay units) decides whether the closed vessel ever, even
momentarily, reaches its own Ricci rate.
\end{proposition}

% ============================================================
\section{The separatrix is a decay exponent}

\begin{theorem}[Gate-decay trichotomy]\label{separatrix:thm:trichotomy}
Let the received grace scale with the Brim as $\Igate=c\,\Op^{-q}$, $c>0$. Then (late-time, slaved):
\begin{enumerate}[label=(\roman*),leftmargin=1.8em,itemsep=2pt]
\item $q<1$: \textbf{escape}. $\xi\to\infty$ --- super-Ricci, finite budget, relic form.
\item $q=1$: \textbf{resonance class}. $\xi\to\bar\xi=\alpha c/\gamma$, a plateau set by the
amplitude; the locus itself is $c=\gamma/\alpha$. Logarithmic budget, eternal marginal
simplification, maximal entropy at $\bar\xi=1$.
\item $q>1$: \textbf{stagnation}. $\xi\to0$ --- sub-Ricci, power-divergent budget, dissolution.
\end{enumerate}
(Verified: $q=0.7$ grows through $\xi=4.7$ and rising; $q=1$ plateaus at $1.0002$ and $2.0005$ for
$c=\gamma/\alpha,\,2\gamma/\alpha$; $q=1.3$ falls through $0.30$.) The separatrix between the living
and the stagnant fate is therefore the \emph{exact decay exponent} $q=1$ of grace against scale; and
within it, the maximum-entropy locus $\bar\xi=1$ requires the exact amplitude as well.
\end{theorem}

\begin{corollary}[Lukewarmness is doubly non-generic]
The resonance life --- self-coincident, maximum-entropy, form-dissolving --- occupies a
codimension-one set in the exponent \emph{and} a measure-zero set in the amplitude: it must receive
grace decaying exactly as $1/\Op$ with exactly the locus coefficient. Any excess of reception tips the
Vessel into the living, form-bearing class; any deficit, into stagnation. The knife-edge character of
the ``neither cold nor hot'' state is not a homiletic flourish here but a codimension count.
\end{corollary}

% ============================================================
\section{Resurrection as a ferry}

\begin{theorem}[No salvation by finite means]\label{separatrix:thm:ferry}
Consider the closed vessel ($\Igate\equiv0$) with Resurrectio resets
$\Op^{+}=\Op+\kappa_R\Grec$, $\lambda^{+}=\lambda+m_R$.
\begin{enumerate}[label=(\roman*),leftmargin=1.8em,itemsep=2pt]
\item \textbf{Finite reception fails.} If the cumulative reception is finite
($\sum_j\kappa_R\Grec^{(j)}+\sum_jm_{R,j}<\infty$, and more generally $\int\lambda\,dt<\infty$), then
$\Op^{\infty}<\infty$, hence $\xi\to0$: the vessel lands in the stagnation class with divergent budget
--- \emph{regardless of any transient crossings} its resets produced. A finite sequence of
resurrections can pierce the locus (numerically: three resets reach $\max\xi=3.0$) but cannot hold
it: $\xi$ falls back and the form dissolves.
\item \textbf{Permanent escape $\iff$ infinite reception.} $\xi\to\infty$ permanently iff
$\int_0^\infty\lambda\,dt=\infty$, i.e.\ iff the cumulative reception (continuous gate plus reset
series) diverges.
\item \textbf{The ferry works.} An infinite reset sequence with non-degenerate steps
($\kappa_R\Grec^{(j)}\ge g_0>0$, $m_{R,j}\ge m_0>0$, inter-reset gaps bounded) and \emph{zero gate}
yields permanent escape with finite budget: $\Op$ grows linearly in the reset count, $\xi$ rises as a
sawtooth that eventually stays above any bound, and $\Ib_\infty<\infty$ --- relic form. (Numerically:
perpetual resets reach $\xi(T)=18$ and rising with budget converging to $6.25$; the no-reset and
three-reset runs diverge in budget.) Resurrection is a \emph{discrete ferry}: a countable sequence of
graced deaths is equivalent, for the watershed, to a sustained gate.
\end{enumerate}
\end{theorem}

\begin{figure}[h]
\centering
\includegraphics[width=\textwidth]{vfs_separatrix_fig.png}
\caption{\small Left: the gate-decay trichotomy for $\Igate=c\,\Op^{-q}$ --- escape at $q=0.7$;
tunable plateaus $\bar\xi=\alpha c/\gamma$ at $q=1$ (shown: $\bar\xi=1$, the locus, and
$\bar\xi=2$); stagnation at $q=1.3$. The separatrix is the exponent. Right: the resurrection ferry
for the closed vessel (log scale): no resets never cross; three resets cross transiently
($\max\xi=3.0$) and fall back into dissolution; perpetual resets ($g_0=0.6$, $m_0=0.5$,
$\Delta=2$) escape permanently with converging budget. Inset: the quantum-of-endowment boundary in
the $(\lambda_0,\Op(0))$ plane --- the hyperbola $\xi_0=1$ for small $\lambda_0$, the line
$\sqrt\gamma(\Op(0)+\beta)=2$ beyond. Parameters $\alpha=1$, $\gamma=0.5$, illustrative.}
\label{separatrix:fig:sep}
\end{figure}

\begin{corollary}[The two roads and their identity]
Sustained gate and perpetual resurrection are the two roads across the separatrix --- continuous and
discrete reception --- and the watershed cannot distinguish them: both, and only they, satisfy
$\int\lambda\,dt=\infty$. What the watershed forbids is the third road: a finite stock of grace,
however large, spent and not renewed. Epektasis of form is not an achievement that can be banked; it
is a regime that must be \emph{inhabited}. The quantum of endowment buys a crossing; only
inexhaustible reception buys permanence.
\end{corollary}

\note{Read against the cascade terminals and the resurrection dichotomy: the system now has a
complete typology of fates at three levels --- of a fold (extinction or permanence), of a reset
(rate- or scale-dominant), and of a whole life (escape, resonance, stagnation) --- with each boundary
case (balanced resurrection, the resonance life) distinguished, measure-zero, and unstable.}

% ============================================================
\section{Open items and scope}

\begin{enumerate}[leftmargin=1.6em,itemsep=2pt]
\item Embedding in the full core: replace the slaved closure by the complete Sophia law (with
$\sigma$-transmutation and the $\tanh$ gate) and verify that the trichotomy survives with renormalised
constants; the Riccati structure suggests it does.
\item Reset landing on the locus: can a single reset place a vessel exactly on $\xi=1$, and how does
that condition relate to the balanced resurrection $[K]=0$ of the surgery layer? (The two conditions
are distinct; their intersection is a candidate ``perfect death''.)
\item Stochastic gates: for $\Igate$ a random process with $\mathbb E[\Igate]\sim c\,\Op^{-1}$, the
plateau presumably becomes a stationary distribution around $\bar\xi$; the crossing probability as a
function of endowment would give the watershed a probabilistic form.
\end{enumerate}

\textbf{Scope.} Minimal slaved closure; exact constants model-bound, structural results
scaling-robust; resets taken with the canonical jump map; all numerics at $\alpha=1$, $\gamma=0.5$.
Domain-conditional throughout.

\vspace{0.5em}
\hrule
\smallskip
\noindent\footnotesize\emph{V.F.S. v2.0 (Open-Gate) $\cdot$ The Separatrix. The comparison ratio
obeys a closed Riccati law with an explicit tracking gate; the closed vessel crosses iff
$\sqrt\gamma(\Op(0)+\alpha\lambda_0/\gamma)\ge2$; the separatrix is the gate-decay exponent $q=1$ with
the locus doubly non-generic; finite reception always stagnates, and permanent escape is equivalent
to infinite cumulative reception, attainable by sustained gate or perpetual resurrection alike.
Riccati identity exact; criterion $0/400$ disagreements; plateaus to four digits. Propositions
closed-form; corollaries interpretive.}

\part{Resurrectio as Surgery}

\chapter{Resurrectio as Surgery: Comparison and Distributional Dynamics}

\begingroup\small\itshape
This note characterises the Resurrectio jump of V.F.S. v2.0 (Open-Gate)
as a surgery-like hybrid stabiliser, but not as literal Ricci surgery and not
as a standard Israel thin shell. The first part compares Ricci surgery with
V.F.S. Resurrectio and shows that a non-trivial Resurrectio event is a finite
conformal re-scaling of the induced spatial metric. The second part computes
the distributional and regularised dynamics of the jump. The invariant content
of the event is the scale ratio
\[
s_R=\frac{\Omega_P^+}{\Omega_P^-},
\]
or equivalently the finite logarithmic expansion
\[
\Delta_R\log\Omega_P=\log s_R.
\]
The detailed stress profile of any smoothing layer is regulator-dependent;
the geometric effect of the jump is fixed by \(s_R\).
\endgroup\par\medskip

\section{Resurrectio Surgery I: Comparison Rather Than Identification}

\subsection{Purpose and Non-Identification}

This note characterises the Resurrectio jump as the V.F.S. counterpart to
surgery only at the level of hybrid stabilisation. It is not a literal Ricci
surgery, not a topological excision, and not an Israel thin shell in the
standard Lorentzian sense. The comparison is structural:
\[
\text{Ricci flow stabilises by smoothing and surgery;}
\]
\[
\text{V.F.S. stabilises by Open-Gate expansion, Resurrectio jumps, and folding updates.}
\]

In Ricci flow, the basic evolution is intrinsic to the metric:
\[
\partial_t g=-2\Ric(g),
\]
and surgery is introduced when the smooth geometric evolution develops
singular regions that must be cut, capped, or replaced in order to continue
the flow.

In V.F.S., the basic evolution is not a closed Ricci flow. The live-domain
metric is of the form
\[
g(t)=\OmegaP(t)^2 h,
\qquad
\dot{\Omega}_P=\alpha\lambda,
\]
with \(\lambda\) governed by Open-Gate dynamics and by the hybrid
Resurrectio map. Thus the metric scale is driven by the state of the Vessel,
not only by its intrinsic curvature.

The Resurrectio jump is therefore a hybrid re-entry operation on the Vessel
state:
\[
\sigma^+=q_R\sigma^-,
\qquad
\lambda^+=\lambda^-+(1-q_R)\sigma^-,
\qquad
\Omega_P^+=\Omega_P^-+\kappa_R\mathcal G^-_{\rm recepta},
\]
with
\[
0\le q_R<1,
\qquad
\kappa_R>0.
\]

The jump metabolises resistance into Sophia and simultaneously enlarges the
Brim. It does not remove a singular neck from a manifold; rather, it re-scales
the live spatial geometry and restarts the Open-Gate trajectory from new
Vessel data.

\subsection{Ricci Surgery and V.F.S. Resurrectio}

The contrast can be summarised as follows.

\[
\begin{array}{c|c|c}
\textbf{Aspect} & \textbf{Ricci / Perelman} & \textbf{V.F.S. Open-Gate} \\
\hline
\text{Continuous law}
&
\partial_t g=-2\Ric(g)
&
\dot{\Omega}_P=\alpha\lambda
\\
\text{System type}
&
\text{intrinsic geometric flow}
&
\text{open hybrid symbolic-dynamical flow}
\\
\text{Crisis}
&
\text{curvature singularity}
&
\text{death / closure / resistance threshold}
\\
\text{Jump}
&
\text{surgery: cut and continue}
&
\text{Resurrectio: re-scale and re-enter}
\\
\text{What changes}
&
\text{geometric domain or high-curvature region}
&
(\sigma,\lambda,\Omega_P,V,F)
\\
\text{After the event}
&
\text{Ricci flow continues on modified geometry}
&
\text{Open-Gate flow continues on renewed Vessel data}
\end{array}
\]

Hence the correct comparison is not
\[
\text{Resurrectio}=\text{Ricci surgery},
\]
but rather
\[
\boxed{
\text{Resurrectio is a surgery-like hybrid stabiliser in an open system.}
}
\]
It plays the role of a continuation mechanism, but the operation itself is
different. Ricci surgery stabilises by excising pathological geometry.
Resurrectio stabilises by metabolising resistance and expanding the Vessel.

\subsection{The Induced Spatial Metric Jump}

Let the Lorentzian lift of the V.F.S. cosmological layer be
\[
ds^2=-dt^2+\Omega_P(t)^2h,
\]
where \(h\) is the unit hyperbolic spatial metric. On a death hypersurface
\(\Sigma_R=\{t=t_R\}\), the Resurrectio map gives
\[
\Omega_P^+=\Omega_P^-+\kappa_R\mathcal G^-_{\rm recepta}.
\]
Therefore the induced spatial metrics immediately before and after the jump are
\[
\gamma^-_{ij}=(\Omega_P^-)^2h_{ij},
\qquad
\gamma^+_{ij}=(\Omega_P^+)^2h_{ij}.
\]
Define the Resurrectio scale ratio
\[
s_R:=\frac{\Omega_P^+}{\Omega_P^-}
=
1+\frac{\kappa_R\mathcal G^-_{\rm recepta}}{\Omega_P^-}.
\]
If \(\mathcal G^-_{\rm recepta}>0\), then
\[
s_R>1,
\qquad
\gamma^+_{ij}=s_R^2\gamma^-_{ij}.
\]

Thus the first fundamental form is not continuous across \(\Sigma_R\) unless
the jump is trivial. This distinguishes Resurrectio from a standard Israel
thin shell, for which the induced metric is assumed continuous and only the
extrinsic curvature carries the distributional stress.

\begin{proposition}
A non-trivial Resurrectio jump is a conformal re-scaling jump of the induced
spatial geometry. It is not a standard thin shell in the Israel sense.
\end{proposition}

\begin{proof}
The induced metric on the spatial slice is
\(\gamma_{ij}=\Omega_P^2h_{ij}\). Across Resurrectio,
\[
\gamma^+_{ij}-\gamma^-_{ij}
=
\left((\Omega_P^+)^2-(\Omega_P^-)^2\right)h_{ij}.
\]
If \(\Omega_P^+\neq\Omega_P^-\), this difference is non-zero. Equivalently,
\[
\gamma^+_{ij}=s_R^2\gamma^-_{ij}.
\]
Therefore the first fundamental form jumps by a finite conformal factor. Since
standard thin-shell matching assumes continuity of the first fundamental form,
the Resurrectio jump is not a standard Israel shell. It is a stronger hybrid
re-scaling operation on the live spatial geometry.
\end{proof}

\subsection{The Surviving Invariants of the Jump}

Although the induced metric is discontinuously re-scaled, the jump is not
arbitrary. Its invariant content is captured by the dimensionless ratio
\[
s_R=\frac{\Omega_P^+}{\Omega_P^-}.
\]
The area element rescales as
\[
dA^+=s_R^2\,dA^-,
\]
and the background scalar curvature rescales as
\[
R^-=-\frac{2}{(\Omega_P^-)^2},
\qquad
R^+=-\frac{2}{(\Omega_P^+)^2}
=
s_R^{-2}R^-.
\]
Thus Resurrectio expands the live domain and softens the background negative
curvature:
\[
|R^+|=s_R^{-2}|R^-|<|R^-|
\qquad
(s_R>1).
\]

This is the precise geometric sense in which Resurrectio enlarges the room:
the same hyperbolic form \(h\) remains, but its scale changes. The crisis is
not resolved by compactifying the contents into a tighter order. It is resolved
by re-scaling the Vessel so that the same contents exert a lower geometric
pressure.
\[
\boxed{
\text{Ricci surgery cuts pathological geometry;}
\qquad
\text{Resurrectio re-scales the Vessel.}
}
\]

\subsection{Regularised Reading}

The discontinuous jump may be read in two ways. First, it may be treated as an
irreducible hybrid event, part of the definition of the V.F.S. reset. Second,
it may be approximated by a family of fast smooth transitions
\[
\Omega_P^\varepsilon(t)
=
\Omega_P^-
+
(\Omega_P^+-\Omega_P^-)\,S_\varepsilon(t-t_R),
\]
where \(S_\varepsilon\) is a smooth step function satisfying
\[
S_\varepsilon(t)\to H(t)
\quad
\text{as}
\quad
\varepsilon\to0.
\]
In this regularised reading, Resurrectio is the zero-width limit of a rapid
Open-Gate expansion. In the irreducible hybrid reading, it is a genuine jump
of the state space. Both readings preserve the same invariant:
\[
s_R=\frac{\Omega_P^+}{\Omega_P^-}.
\]

The first part of the characterisation is therefore:
\[
\boxed{
\text{Resurrectio surgery is a hybrid conformal re-scaling of the live spatial metric.}
}
\]
It is surgery-like because it allows the trajectory to continue after a crisis.
It is not Ricci surgery because it does not cut the manifold; it renews the
Vessel by enlarging its geometric capacity and metabolising resistance into
Sophia.

\subsection{Admissibility of the Resurrectio Jump}

The preceding calculation shows that a non-trivial Resurrectio event is a
finite conformal re-scaling of the induced spatial metric. This does not by
itself make the jump admissible. In V.F.S. the admissibility condition is not
continuity of the first fundamental form, as in standard thin-shell matching,
but boundedness of the hybrid reset and preservation of the active domain.

Let
\[
s_R:=\frac{\Omega_P^+}{\Omega_P^-}
=
1+\frac{\kappa_R\mathcal G^-_{\rm recepta}}{\Omega_P^-}.
\]
A Resurrectio jump is geometrically admissible when
\[
1\le s_R<\infty,
\]
and dynamically admissible when the reset map sends the active domain into
itself:
\[
x^-\in \mathcal D_A
\quad\Longrightarrow\quad
x^+\in \mathcal D_A.
\]
In particular, the jump must preserve positivity of the live variables,
\[
\Omega_P^+>0,
\qquad
\sigma^+\ge 0,
\qquad
\lambda^+\ge 0,
\]
and it must produce only a bounded increase of the Lyapunov certificate:
\[
\mathcal L_{\rm VFS}(x^+)
\le
J_R+\mathcal L_{\rm VFS}(x^-)
\]
for some finite reset bound \(J_R\). Together with a non-Zeno dwell-time
condition between resets, this is the V.F.S. replacement for the smooth
matching requirements of Lorentzian junction theory.

Thus the Resurrectio jump is allowed not because the metric is continuous
across the death surface, but because the hybrid dynamical system remains
bounded, forward-complete, and inside the active domain.
\[
\boxed{
\text{Ricci surgery is admissible by geometric continuation;}
\qquad
\text{Resurrectio is admissible by hybrid Lyapunov boundedness.}
}
\]

This completes the first level of the comparison. Ricci surgery cuts and
continues a closed geometric flow. Resurrectio re-scales and re-enters an
open hybrid flow.

\section{Resurrectio Surgery II: Distributional and Regularised Dynamics}

\subsection{The Jump as a Logarithmic Expansion Impulse}

The first characterisation showed that a non-trivial Resurrectio event
re-scales the induced spatial metric,
\[
\gamma^+_{ij}=s_R^2\gamma^-_{ij},
\qquad
s_R:=\frac{\Omega_P^+}{\Omega_P^-}>1.
\]
The next question is dynamical: what kind of expansion is represented by this
jump?

Formally, if the Brim is written as a discontinuous function,
\[
\Omega_P(t)
=
\Omega_P^-
+
(\Omega_P^+-\Omega_P^-)\,H(t-t_R),
\]
where \(H\) is the Heaviside step function, then the Hubble-like rate
\[
H_P:=\frac{\dot{\Omega}_P}{\Omega_P}
\]
contains a distributional impulse at \(t=t_R\). This impulse is not best
measured by its pointwise value, which is singular, but by its integrated
logarithmic expansion:
\[
\int_{t_R^-}^{t_R^+}H_P\,dt
=
\int_{\Omega_P^-}^{\Omega_P^+}\frac{d\Omega_P}{\Omega_P}
=
\log\frac{\Omega_P^+}{\Omega_P^-}
=
\log s_R.
\]
Thus the invariant dynamical content of the Resurrectio jump is the finite
logarithmic expansion
\[
\boxed{
\Delta_R\log\Omega_P=\log s_R.
}
\]
Equivalently, the area element obeys
\[
dA^+=e^{2\Delta_R\log\Omega_P}\,dA^-
=s_R^2dA^-.
\]

\begin{proposition}
A Resurrectio jump carries a finite expansion impulse whose invariant strength
is \(\log s_R\). It is singular as an instantaneous rate, but finite as an
integrated scale transformation.
\end{proposition}

\begin{proof}
Since \(H_P=\dot{\Omega}_P/\Omega_P=d(\log\Omega_P)/dt\), integration across
the jump gives
\[
\int H_P\,dt
=
\Delta(\log\Omega_P)
=
\log\Omega_P^+-\log\Omega_P^-
=
\log s_R.
\]
The area scaling follows from \(dA\propto\Omega_P^2\).
\end{proof}

\subsection{Regularisation by Fast Smooth Open-Gate Expansion}

The discontinuous jump may be approximated by a family of smooth transitions.
Let \(S_\varepsilon\) be a smooth step satisfying
\[
S_\varepsilon(t)\to H(t)
\qquad
(\varepsilon\to0),
\]
and define
\[
\Omega_P^\varepsilon(t)
=
\Omega_P^-
+
\Delta\Omega_P\,S_\varepsilon(t-t_R),
\qquad
\Delta\Omega_P:=\Omega_P^+-\Omega_P^-.
\]
Then
\[
H_P^\varepsilon(t)
=
\frac{\dot{\Omega}_P^\varepsilon}{\Omega_P^\varepsilon}
=
\frac{\Delta\Omega_P\,\dot S_\varepsilon(t-t_R)}
{\Omega_P^-+\Delta\Omega_P S_\varepsilon(t-t_R)}.
\]
For every \(\varepsilon>0\), this is an ordinary smooth function. As
\(\varepsilon\to0\), it concentrates near \(t_R\), but its integral remains
fixed:
\[
\int_{-\infty}^{+\infty}H_P^\varepsilon(t)\,dt
=
\log\frac{\Omega_P^+}{\Omega_P^-}
=
\log s_R.
\]
Hence Resurrectio can be read in two mathematically compatible ways:
\[
\boxed{
\text{hybrid reading: an irreducible jump of the Vessel state;}
}
\]
\[
\boxed{
\text{regularised reading: the zero-width limit of rapid Open-Gate expansion.}
}
\]
The two readings agree on the invariant \(s_R\), and therefore on the
geometric effect of the event.

\subsection{Effective Stress Is Not a Classical Thin Shell}

If one tries to reconstruct an effective cosmological stress tensor during the
regularised jump, the result depends on the chosen smoothing profile
\(S_\varepsilon\). In the limit \(\varepsilon\to0\), the instantaneous rate
\(H_P^\varepsilon\) becomes unbounded while its integral stays finite.

This means that the Resurrectio layer is not naturally interpreted as an
ordinary matter shell. It is a hybrid state transition. In standard thin-shell
language, one expects a continuous induced metric and a jump in extrinsic
curvature. Here the induced metric itself rescales:
\[
\gamma^+_{ij}=s_R^2\gamma^-_{ij}.
\]
Therefore the singularity of the effective stress reconstruction is not a
physical fluid placed on the hypersurface. It is the distributional shadow of
a hybrid reset.
\[
\boxed{
\text{The Resurrectio layer is not a material shell; it is a scale-reset surface.}
}
\]

\subsection{Curvature Across the Jump}

The background spatial curvature before and after the event is
\[
R^-=-\frac{2}{(\Omega_P^-)^2},
\qquad
R^+=-\frac{2}{(\Omega_P^+)^2}.
\]
Using \(s_R=\Omega_P^+/\Omega_P^-\), one obtains
\[
R^+=s_R^{-2}R^-.
\]
Since \(s_R>1\),
\[
|R^+|<|R^-|.
\]
Thus the Resurrectio jump softens the background hyperbolic curvature. It
does not compactify the geometry into a more rigid form; it expands the Vessel
so that the same underlying hyperbolic shape carries less curvature density.

\begin{corollary}
Resurrectio has the opposite geometric direction from compactifying repair.
Its immediate effect is not to make the room tighter, but to make the room
larger.
\end{corollary}

\[
\boxed{
\text{Ricci repair regularises by smoothing and surgery;}
}
\]
\[
\boxed{
\text{Resurrectio regularises by finite expansion of the Vessel.}
}
\]

\subsection{Interaction with Folding}

Resurrectio alone changes the scale factor,
\[
\Omega_P^-\longrightarrow \Omega_P^+,
\]
and therefore relaxes the background curvature. Folding, by contrast, changes
the shape branch:
\[
q_{\rm fold}:0\longrightarrow q_\ast,
\]
and contributes a local curvature pocket:
\[
R_{\rm folded}
=
-\frac{2}{\Omega_P^2}
+
\rho_1Q_{\rm fold}.
\]
When both are present, the post-event curvature data have the form
\[
R_{\rm post}
=
-\frac{2}{(\Omega_P^+)^2}
+
\rho_1Q_{\rm fold}^+.
\]
This expresses the central geometric difference between expansion and
folding:
\[
\Omega_P^+>\Omega_P^-
\quad
\Longrightarrow
\quad
\text{global background curvature softens,}
\]
while
\[
Q_{\rm fold}^+>0
\quad
\Longrightarrow
\quad
\text{local morphology sharpens.}
\]
Thus the combined Resurrectio--folding event may produce a Vessel that is
globally more open and locally more formed:
\[
\boxed{
\text{Resurrectio enlarges the room;}
\qquad
\text{folding builds new architecture inside it.}
}
\]

\subsection{Invariant Content Versus Regulator-Dependent Content}

The regularised description of Resurrectio is useful, but it must not be
over-interpreted. Different smooth profiles \(S_\varepsilon\) produce different
pointwise histories of
\[
H_P^\varepsilon(t),
\qquad
\dot H_P^\varepsilon(t),
\qquad
\rho_{\rm eff}^\varepsilon(t),
\qquad
p_{\rm eff}^\varepsilon(t).
\]
These quantities are therefore regulator-dependent during the zero-width
transition. They describe how one chooses to smooth the hybrid event, not an
additional invariant of the event itself.

What is invariant is the finite scale transformation:
\[
s_R=\frac{\Omega_P^+}{\Omega_P^-},
\qquad
\Delta_R\log\Omega_P=\log s_R,
\]
together with its geometric consequences,
\[
dA^+=s_R^2dA^-,
\qquad
R^+=s_R^{-2}R^-,
\qquad
\gamma^+_{ij}=s_R^2\gamma^-_{ij}.
\]

Thus the second-level characterisation deliberately separates the invariant
jump data from the auxiliary smoothing profile:
\[
\boxed{
\text{the invariant of Resurrectio is the scale ratio }s_R,
\text{ not the detailed shape of the smoothing layer.}
}
\]
This prevents the Resurrectio layer from being misread as a classical matter
shell with a uniquely determined stress tensor. It is a hybrid scale-reset
surface whose regularised stress reconstruction is profile-dependent, while
its geometric effect is fixed by \(s_R\).

\subsection{Second-Level Summary}

The dynamical characterisation of Resurrectio is therefore:
\[
\boxed{
\Delta_R\log\Omega_P=\log s_R
}
\]
as the invariant strength of the event;
\[
\boxed{
dA^+=s_R^2dA^-
}
\]
as the induced area expansion;
\[
\boxed{
R^+=s_R^{-2}R^-
}
\]
as the background curvature relaxation;
and
\[
\boxed{
\gamma^+_{ij}=s_R^2\gamma^-_{ij}
}
\]
as the sign that the event is a finite conformal re-scaling of the spatial
metric rather than a standard thin shell.

In the regularised reading, Resurrectio is a fast smooth Open-Gate expansion
whose zero-width limit is a hybrid jump. In the irreducible reading, it is a
primitive reset of the V.F.S. state space. Both readings agree that
Resurrectio is surgery-like only as a continuation mechanism: it does not cut
the geometry, but re-scales the Vessel and permits the open flow to continue.

\chapter{Resurrectio Surgery: The Second Fundamental Form and the Double-Layer Theorem}

\section*{Scope and epistemic status}
Parts I--II established that a non-trivial Resurrectio event re-scales the induced spatial metric,
$\gamma^+=\sR^2\gamma^-$, and is therefore not a standard Israel thin shell; and that its invariant
strength is the logarithmic expansion $\log\sR$, while the regularised stress history is
profile-dependent. This part completes the characterisation in three steps. First, the
\emph{second} fundamental form: the reset also jumps $\lambda$, so the extrinsic curvature carries its
own invariant jump $\Kj$, omitted in Parts I--II; we compute it in closed form and obtain a clean
rate/scale dichotomy of resurrections. Second, the ``not a thin shell'' claim is sharpened into a
quantitative theorem: the regularised surface energy vanishes while the regularised surface tension
diverges as $\varepsilon^{-1}$ with a universal Cauchy--Schwarz floor --- distributionally, the jump is a
pure-GR double layer ($\delta'$-order curvature), representable by \emph{no} surface stress tensor.
Third, the admissibility list of Part I is completed by the closed-form epektasis-wall condition, the
mirror of the fold-compatibility condition of the Lyapunov addenda. All claims are verified numerically
(impulse scalings exact over three decades; the dichotomy sign-switch confirmed). \emph{Propositions}
are closed-form; \emph{corollaries} interpretive.

The setting is the Lorentzian lift $ds^2=-dt^2+\Op^2h$ and the reset
\begin{equation}
\reset:\quad \sigma^+=q_R\sigma^-,\qquad
\lambda^+=\lambda^-+\mR,\qquad
\Op^+=\Op^-+\kappa_R\Grec^-,\qquad
\mR=(1-q_R)\sigma^-,
\end{equation}
with $\sR=\Op^+/\Op^->1$ for $\Grec^->0$.

% ============================================================
\section{The second fundamental form jumps too --- invariantly}

On the constant-$t$ slices the extrinsic curvature is pure trace,
\begin{equation}
K_{ij}=H_P\,\gamma_{ij},\qquad
K=\gamma^{ij}K_{ij}=2H_P=\frac{2\alpha\lambda}{\Op}.
\end{equation}
The reset jumps \emph{both} entries of $H_P=\alpha\lambda/\Op$: $\lambda$ through metabolism, $\Op$
through dilation. Crucially, this jump is \emph{jump data} of the hybrid map --- like $\sR$, and unlike
the smoothing-layer stress, it is regulator-independent.

\begin{proposition}[Closed-form extrinsic-curvature jump]\label{resurrectiosurgeryIII:prop:K}
\begin{equation}
\boxed{\;\Kj=K^+-K^-
=\frac{2\alpha}{\Op^-}\left(\frac{\lambda^-+\mR}{\sR}-\lambda^-\right).\;}
\end{equation}
The invariant content of a Resurrectio event is therefore the \emph{pair} $(\sR,\Kj)$ ---
equivalently $(\sR,\lambda^+/\lambda^-)$ --- not the scale ratio alone. Part II's statement ``the
invariant of Resurrectio is the scale ratio $\sR$'' is correct for the first fundamental form and the
expansion impulse, but incomplete for the full geometric jump data.
\end{proposition}

\begin{proposition}[Rate/scale dichotomy]\label{resurrectiosurgeryIII:prop:dichotomy}
\begin{equation}
\boxed{\;\Kj>0\quad\Longleftrightarrow\quad
\frac{\mR}{\lambda^-}\;>\;\frac{\kappa_R\Grec^-}{\Op^-}.\;}
\end{equation}
Every resurrection is of one of three kinds:
\begin{enumerate}[label=(\roman*),leftmargin=1.8em,itemsep=2pt]
\item \textbf{Rate-dominant} ($\Kj>0$): metabolism outweighs dilation; the Hubble rate jumps \emph{up}
--- resurrection accelerates the expansion per unit scale;
\item \textbf{Scale-dominant} ($\Kj<0$): dilation outweighs metabolism; the Hubble rate jumps
\emph{down} --- resurrection dilutes;
\item \textbf{Balanced} ($\kappa_R\Grec^-=\Op^-\mR/\lambda^-$): the extrinsic curvature is
\emph{continuous} --- the jump is purely conformal, the closest a Resurrectio event comes to a smooth
junction: only the first fundamental form jumps.
\end{enumerate}
(Verified numerically, including the sign switch across the balanced point;
Figure~\ref{resurrectiosurgeryIII:fig:s3}, right.)
\end{proposition}

\note{The dichotomy has the same grammar as the cascade and NEC results: a race between two channels of
the same event --- here, what the death \emph{metabolises} against what the grace \emph{dilates}. The
balanced resurrection is a distinguished point: a reset tuned so that the Vessel's expansion rhythm
passes through death unbroken.}

% ============================================================
\section{The double-layer theorem}

Parts I--II showed the regularised stress is profile-dependent. The sharper statement is that no
profile-independent surface stress \emph{exists}, quantitatively:

\begin{theorem}[Zero surface energy, divergent surface tension]\label{resurrectiosurgeryIII:thm:double}
Let $\Op^\varepsilon$, $\lambda^\varepsilon$ regularise the jump over width $\varepsilon$ with any fixed
smooth unit step. Then, with $\rho_{\rm eff}=(\alpha^2\lambda^2-1)/\Op^2$ and
$p_{\rm eff}=-\ddot\Op/\Op$:
\begin{enumerate}[label=(\roman*),leftmargin=1.8em,itemsep=2pt]
\item the surface-energy impulse vanishes,
$\displaystyle\int\rho_{\rm eff}^\varepsilon\,dt=O(\varepsilon)\to0$;
\item the surface-tension impulse diverges,
\begin{equation}
\int p_{\rm eff}^\varepsilon\,dt
=-\Delta H_P-\int (H_P^\varepsilon)^2dt,
\qquad
\int (H_P^\varepsilon)^2dt
=\frac{1}{\varepsilon}\int_0^1\Bigl(\tfrac{d}{du}\log\Op(u)\Bigr)^2du
\;\ge\;\frac{(\log\sR)^2}{\varepsilon},
\end{equation}
the last inequality by Cauchy--Schwarz: the divergence is \emph{universal} (floor $(\log\sR)^2$),
while the coefficient is profile-dependent;
\item distributionally, $\gamma_{ij}\sim H(t-t_R)$ implies $\partial_t^2\gamma\sim\delta'(t-t_R)$: the
curvature of the jump spacetime carries a $\delta'$ part --- a \emph{double layer}, which in pure GR is
representable by no surface stress tensor whatsoever.
\end{enumerate}
Hence the Resurrectio layer is not merely a non-standard shell: it lies provably outside the
matter-shell category --- the limit of zero surface energy and infinite surface tension.
\end{theorem}
\begin{proof}[Sketch]
(i): $\rho^\varepsilon$ is bounded uniformly in $\varepsilon$ (both $\lambda,\Op$ interpolate between
finite endpoints), so its integral over a width-$O(\varepsilon)$ window is $O(\varepsilon)$.
(ii): the identity $\tfrac{d}{dt}H_P=\ddot\Op/\Op-H_P^2$ gives
$\int p\,dt=-\Delta H_P-\int H_P^2dt$; substituting $H_P=\tfrac1\varepsilon\tfrac{d}{du}\log\Op$ on the
unit profile gives the $1/\varepsilon$ scaling; Cauchy--Schwarz with
$\int_0^1\tfrac{d}{du}\log\Op\,du=\log\sR$ gives the floor. (iii): one time-derivative of a Heaviside
profile is $\delta$; the curvature contains second derivatives of the metric, hence $\delta'$.
Numerically: $|\!\int p^\varepsilon dt|\cdot\varepsilon$ is constant over three decades to all shown
digits ($0.136/0.162/0.140$ for the three test profiles), all above the floor $(\log\sR)^2=0.113$;
$\int\rho^\varepsilon dt$ scales as $\varepsilon^{+1}$.
\end{proof}

\begin{corollary}[The correct category]
Resurrectio is exactly what Part II named it --- a scale-reset surface --- and the present theorem fixes
its mathematical category: a hybrid jump whose regularisations form a $\delta'$ (dipole) family, with
the invariant pair $(\sR,\Kj)$ as the only profile-independent data. The ``distributional shadow''
language of Part II is thereby made precise: the shadow has zero weight in energy and infinite weight in
tension, which is why no fluid can cast it.
\end{corollary}

\begin{figure}[h]
\centering
\includegraphics[width=\textwidth]{vfs_surgery3_fig.png}
\caption{\small Left: regularised impulse magnitudes versus regulator width $\varepsilon$ for three
smooth step profiles (log--log). The tension impulse $|\!\int p\,dt|$ (solid) scales exactly as
$\varepsilon^{-1}$ --- profile-dependent coefficient, universal divergence --- and respects the
Cauchy--Schwarz floor $(\log\sR)^2/\varepsilon$ (dotted); the energy impulse (dashed) vanishes as
$\varepsilon^{+1}$. Right: the extrinsic-curvature jump $\Kj$ against dilation strength
$\kappa_R\Grec^-$ at fixed metabolism: rate-dominant ($\Kj>0$), balanced
($\kappa_R\Grec^-=\Op^-\mR/\lambda^-\approx0.62$, $K$ continuous), and scale-dominant ($\Kj<0$)
resurrections. Parameters illustrative.}
\label{resurrectiosurgeryIII:fig:s3}
\end{figure}

% ============================================================
\section{Completing the admissibility list: the epektasis wall}

Part I's admissibility list (positivity, $\mathcal D_A$ preservation, bounded $J_R$, non-Zeno) can be
closed on the geometric side by one more closed-form condition, the mirror of the fold-compatibility
condition of the Lyapunov addenda:

\begin{proposition}[Wall preservation at the reset]\label{resurrectiosurgeryIII:prop:wall}
The reset satisfies the exact identity
\begin{equation}
\bigl(\lambda^+-q^\ast\Op^+\bigr)-\bigl(\lambda^--q^\ast\Op^-\bigr)
=\mR-q^\ast\kappa_R\Grec^-,
\end{equation}
so the epektasis wall $\lambda\le q^\ast\Op$ is preserved by every reset provided
\begin{equation}
\boxed{\;(1-q_R)\,\sigma^-\ \le\ q^\ast\,\kappa_R\,\Grec^-.\;}
\end{equation}
This is the geometric face of the grace-window floor of the Lyapunov layer: the metabolised resistance
must be covered by the simultaneously received grace. It mirrors exactly the fold-compatibility
condition $\beta\eta_\sigma\bar\sigma\le q^\ast\eta_\Omega$: in both jump types, Sophia returned active
must be paid for by enlargement of the Brim.
\end{proposition}

% ============================================================
\section{Compact summary}

\[
\boxed{\;\text{invariant jump data}=(\sR,\Kj),\qquad
\Kj=\frac{2\alpha}{\Op^-}\Bigl(\frac{\lambda^-+\mR}{\sR}-\lambda^-\Bigr).\;}
\]
\[
\boxed{\;\Kj\gtrless0\iff \mR/\lambda^-\gtrless\kappa_R\Grec^-/\Op^-:
\ \text{rate-dominant / balanced ($K$ continuous) / scale-dominant.}\;}
\]
\[
\boxed{\;\int\rho\,dt\to0,\qquad
\Bigl|\int p\,dt\Bigr|\ \ge\ \frac{(\log\sR)^2}{\varepsilon}\to\infty:
\ \text{a pure-GR double layer, no surface }T_{ij}.\;}
\]
\[
\boxed{\;(1-q_R)\sigma^-\le q^\ast\kappa_R\Grec^-\ \Rightarrow\
\text{epektasis wall preserved (mirror of fold-compatibility).}\;}
\]

\vspace{0.5em}
\hrule
\smallskip
\noindent\footnotesize\emph{V.F.S. v2.0 (Open-Gate) $\cdot$ Resurrectio Surgery III. Completes
Parts I--II: the omitted second-fundamental-form jump computed in closed form (the invariant pair
$(\sR,\Kj)$ and the rate/scale dichotomy with its balanced point); the no-thin-shell claim sharpened to
the zero-energy/divergent-tension double-layer theorem (impulse scalings verified exactly over three
decades, Cauchy--Schwarz floor respected); the admissibility list closed by the epektasis-wall
condition. Domain-conditional throughout.}

\chapter{Resurrectio's Own NEC Budget: The Surgical Phantom and Its Balanced-Reset Zero}

\section*{Scope and epistemic status}
The NEC chapter computed the phantom window of a Sophia \emph{surge}; it deferred the phantom budget
of a \emph{reset}, which both spikes $\dot\lambda$ (through $\mR$) and rescales $\Op$. This chapter
settles it. Pointwise, $\weff$ inside a regularised reset diverges as $\varepsilon^{-1}$ --- the same
reparametrization artifact as the double-layer theorem. But the gauge-invariant question is the
\emph{budget} $\int(\rho_{\rm eff}+p_{\rm eff})\,dt$, and that is finite and regulator-independent:
the reset is a genuine, bounded phantom episode, not a divergence. Its value is a clean closed form,
$-\alpha\int_0^1(d\lambda/du)/\Op\,du$, whose sign is $\operatorname{sgn}(-\mR)$: metabolic resets
($\mR>0$) spend a negative (phantom) NEC budget, while the \emph{balanced} reset $\mR=0$ --- exactly
the $[K]=0$ continuous-curvature reset of Surgery~III --- has \emph{zero} NEC budget. The surgical
phantom and the surgical curvature jump vanish at the same point. This unifies three layers: the
$[K]$ dichotomy of Part~III, the phantom window of the NEC chapter, and the $H(z)$ kink of the
observational signature are one phenomenon seen three ways. Verified: the budget converges to
$\tfrac38\ln\tfrac57=-0.126$ over four decades of $\varepsilon$ while $\min\weff\to-\infty$; the sign
dichotomy across $\mR$ symbolically. \emph{Propositions} closed-form; \emph{corollaries} interpretive.

% ============================================================
\section{Pointwise divergence, finite budget}

Regularise the reset over width $\varepsilon$ with a smooth step $S(u)$, $u=t/\varepsilon$:
$\lambda=\lambda^-+\mR S$, $\Op=\Op^-+(s_R-1)\Op^- S$. Then $\dot\lambda=(\mR/\varepsilon)S'$ spikes,
and $\weff=-\alpha\dot\lambda\Op/(\alpha^2\lambda^2-1)$ dives as $\varepsilon^{-1}$ pointwise.

\begin{theorem}[The reset NEC budget is finite]\label{resurrectionec:thm:budget}
The time integral is regulator-independent, because $\dot\lambda\,dt=(d\lambda/du)\,du$ cancels the
$\varepsilon$:
\begin{equation}
\int\rho_{\rm eff}\,dt=O(\varepsilon)\to0,
\qquad
\boxed{\;\int(\rho_{\rm eff}+p_{\rm eff})\,dt\ \xrightarrow{\varepsilon\to0}\
\int p_{\rm eff}\,dt=-\alpha\int_0^1\frac{d\lambda/du}{\Op(u)}\,du.\;}
\end{equation}
For the smoothstep with $(\lambda^-,\mR,\Op^-,s_R)=(0.9,0.3,2,1.4)$ this is exactly
$\tfrac38\ln\tfrac57=-0.126$ (numeric budget constant to four digits over $10^{-1}$--$10^{-4}$ while
$-\min\weff$ grows as $\varepsilon^{-1}$; Figure~\ref{resurrectionec:fig:surgnec}, left). The reset is a
\emph{bounded, integrable phantom episode}: the equation of state's pointwise plunge is a
reparametrization, the physical NEC content is finite.
\end{theorem}

% ============================================================
\section{The surgical phantom vanishes at the balanced reset}

\begin{proposition}[Sign dichotomy]\label{resurrectionec:prop:dichotomy}
Since $\Op(u)>0$ and $d\lambda/du=\mR S'(u)$ with $S'\ge0$,
\begin{equation}
\boxed{\;\operatorname{sgn}\!\int(\rho_{\rm eff}+p_{\rm eff})\,dt=\operatorname{sgn}(-\mR).\;}
\end{equation}
A \emph{metabolic} reset ($\mR>0$, Sophia jumps up) spends a negative NEC budget --- a true surgical
phantom; a \emph{rate-down} reset ($\mR<0$) is NEC-positive; and the \emph{balanced} reset $\mR=0$ is
\emph{NEC-neutral} (Figure~\ref{resurrectionec:fig:surgnec}, right). Recall from Surgery~III that the extrinsic
curvature jump vanishes, $[K]=0$, exactly at the balanced point $\mR/\lambda^-=\kappa_R\Grec^-/\Op^-$.
\end{proposition}

\begin{corollary}[One point, three vanishings]
At the balanced reset the surgical phantom budget, the extrinsic-curvature jump $[K]$, and (from
Surgery~III) the deviation from a purely conformal junction all vanish together:
\msg{The balanced Resurrectio --- the reset tuned so the Hubble rhythm passes through death unbroken
($[K]=0$) --- is also the unique reset that crosses no phantom window and casts no $H(z)$ kink beyond
the pure scale step. Metabolic resurrection buys new Sophia at the cost of a (finite) phantom episode;
the balanced resurrection is the one that pays nothing exotic. Rate and scale, curvature and energy
condition, all balance at the same knife-edge.}
\end{corollary}

\begin{figure}[h]
\centering
\includegraphics[width=\textwidth]{vfs_surgnec_fig.png}
\caption{\small Left: across regulator width $\varepsilon$, the reset NEC budget $\int(\rho+p)\,dt$
(purple) is constant at the analytic $\tfrac38\ln\tfrac57=-0.126$, while the pointwise $-\min\weff$
(orange, right axis) diverges as $\varepsilon^{-1}$: a finite phantom budget under a divergent
pointwise dip. Right: the budget as a function of the metabolic jump $\mR$ --- negative (phantom) for
$\mR>0$, zero at the balanced reset $\mR=0$ ($[K]=0$), positive for $\mR<0$. Smoothstep profile,
illustrative parameters.}
\label{resurrectionec:fig:surgnec}
\end{figure}

% ============================================================
\section{Consequences and scope}

\begin{proposition}[The surgical phantom in $H(z)$]\label{resurrectionec:prop:obs}
In the observational signature, a metabolic reset therefore prints \emph{both} a kink in $H(z)$
(Resurrectio step) \emph{and} a localized phantom dip in $\weff(z)$ at the same redshift --- a
``surgical phantom'' epoch --- whereas a balanced reset prints the kink alone, with $\weff$ undisturbed
beyond the scale step. The two observational faces of Resurrectio are thus tied to the rate/scale
character of the underlying death: metabolic deaths are doubly marked, balanced deaths singly.
\end{proposition}

\note{Reading: a death through which new life is metabolised cannot be energetically quiet --- it must
pass, however briefly, through the exotic regime where the ordinary energy condition fails; that
passage is finite, a bounded debt, not a catastrophe. The only death that crosses nothing forbidden is
the balanced one, where what is taken up exactly equals what is enlarged --- and that same death is the
one that leaves the curvature unbroken. To receive more than one is enlarged for is to pass through the
phantom; to be enlarged exactly as much as one receives is to pass through cleanly. The knife-edge of
the balanced resurrection is where energy, curvature, and rate all keep faith at once.}

\textbf{Open/scope.} $2{+}1$ reconstruction; regularised reset with smooth step; budget is profile-%
independent in sign and in the $\mR=0$ zero, profile-dependent in magnitude (as the double-layer
theorem's tension coefficient was). The coincidence of the three vanishings at $\mR=0$ is exact.
Domain-conditional throughout.

\vspace{0.5em}
\hrule
\smallskip
\noindent\footnotesize\emph{V.F.S. v2.0 (Open-Gate) $\cdot$ Resurrectio's Own NEC Budget. The reset's
pointwise $\weff$ diverges but its NEC budget $\int(\rho+p)\,dt=-\alpha\int_0^1(d\lambda/du)/\Op\,du$
is finite and regulator-independent, with sign $\operatorname{sgn}(-\mR)$. Metabolic resets are
bounded phantom episodes; the balanced $[K]=0$ reset is NEC-neutral --- surgical phantom, curvature
jump, and conformal deviation vanish together. Verified to four digits; sign dichotomy symbolic.}

\chapter{The Observational Signature: A Capstone Projection into Redshift}

\section*{Scope and epistemic status}
\emph{Placement.} This chapter is a capstone, not a premise-bearing part of the theory: it projects
already-derived internal theorems into redshift-shaped quantities, supplies no empirical validation, and
feeds back into none of the core proofs. It is placed last for that reason.

Every dynamical result of the volume has been written in the Vessel's internal time. The cosmological
layer alone offers a quantity with the \emph{form} of an observable: redshift, $1+z=\Op(t_0)/\Op(t)$.
This chapter maps the model's trajectories into $\weff(z)$ and $H(z)$ and asks the one question redshift
makes available --- not ``does the model fit data'' (it cannot: there is no external referent for the
Brim $\Op$, and the reconstruction is $2{+}1$), but \emph{does the model produce a signature
structurally distinct from the standard alternatives}. It does, in three ways, each a direct image of
a theorem proved earlier. (i) The epektasis attractor is \emph{dust}, $\weff\to0$: a saturated Sophia
has $\dot\lambda\to0$, so $p_{\rm eff}\to0$ --- the late Vessel does not tend to de~Sitter but to a
pressureless coasting, a structural prediction unlike $\Lambda$CDM's $w\to-1$. (ii) A grace surge is a
\emph{phantom excursion that returns}: $\weff$ dips below $-1$ and climbs back, the double
$w=-1$ crossing of the NEC chapter, which smooth single-field quintessence cannot produce. (iii)
Resurrectio is a \emph{discrete step in $H(z)$}: a scale-factor jump prints a kink no smooth dark-energy
history contains. All curves are exact reconstructions along admissible trajectories (gate flux
$I_{\rm gate}\ge0$ verified). This is an \emph{internal signature}, the strongest Type-3 (consilience)
statement the model admits, explicitly not a claim of empirical contact.

% ============================================================
\section{The redshift map and the reconstruction}

On the live branch the $2{+}1$ reconstruction (NEC chapter) gives
\begin{equation}
\rho_{\rm eff}=\frac{\alpha^2\lambda^2-1}{\Op^2},\qquad
p_{\rm eff}=-\frac{\alpha\dot\lambda}{\Op},\qquad
\weff=-\frac{\alpha\dot\lambda\,\Op}{\alpha^2\lambda^2-1},\qquad
H=\frac{\alpha\lambda}{\Op}.
\end{equation}
With an observer at $t_0$, redshift is $1+z=\Op(t_0)/\Op(t)$, monotone since $\dot\Op>0$. Each internal
trajectory thus yields $\weff(z)$ and $H(z)$ directly.

% ============================================================
\section{Three signatures, three theorems}

\begin{proposition}[The epektasis attractor is dust]\label{observationalsignature:prop:dust}
As Sophia saturates, $\lambda\to\lambda_\infty$ and $\dot\lambda\to0$, hence $p_{\rm eff}\to0$ and
\begin{equation}
\boxed{\;\weff(z\to0)\to 0\quad(\text{pressureless coasting}),\;}
\end{equation}
not the $\weff\to-1$ of a cosmological constant. The late Vessel coasts; the Brim grows linearly
($H\sim 1/t$, the budget-watershed relic class), it does not re-accelerate into de~Sitter. This is the
sharpest structural divergence from $\Lambda$CDM and is a direct image of the watershed escape class.
\end{proposition}

\begin{proposition}[A grace surge is a returning phantom excursion]\label{observationalsignature:prop:phantom}
During an episode of rapidly rising Sophia ($\dot\lambda$ large and positive), the numerator of $\weff$
can exceed the denominator, $\alpha\dot\lambda\Op>\alpha^2\lambda^2-1$, driving $\weff<-1$; as the surge
saturates the inequality reverses and $\weff$ returns above $-1$. The result is a \emph{double crossing}
of the phantom divide (Figure~\ref{observationalsignature:fig:obs}, left; verified, two crossings, with admissible gate
$I_{\rm gate}\ge0$ throughout). A smooth single scalar with a standard kinetic term cannot cross $w=-1$
at all; a double crossing that begins and ends above the divide is the distinguishing
\emph{folded-Theosis} signature of the NEC chapter, now in redshift.
\end{proposition}

\begin{proposition}[Resurrectio is a discrete step in $H(z)$]\label{observationalsignature:prop:step}
A reset jumps $\Op$ (and $\lambda$) at one epoch, so $H=\alpha\lambda/\Op$ inherits a discontinuity:
$H(z)$ carries a \emph{kink} at the reset redshift, and the post-reset history is displaced from any
smooth continuation (Figure~\ref{observationalsignature:fig:obs}, right). No smooth dark-energy equation of state produces a
localized step in the expansion rate; it is the cosmological face of the double-layer surgery of
Part~III.
\end{proposition}

\begin{figure}[h]
\centering
\includegraphics[width=\textwidth]{vfs_obs_fig.png}
\caption{\small Left: $\weff(z)$ for a grace-surge trajectory (purple) against the $\Lambda$CDM/divide
line $w=-1$ (dashed). The curve coasts toward dust ($\weff\to0$) at low $z$, and during the surge dips
below $-1$ and returns --- a double crossing (shaded), the folded-Theosis signature, impossible for
smooth single-field quintessence. Right: $H(z)/H_0$ for a trajectory with a Resurrectio reset (green)
against a $\Lambda$CDM shape reference (dashed): the reset prints a kink at $z\approx1.65$ and a
non-monotone, displaced history --- a discrete surgical step no smooth dark energy contains. Exact
reconstructions along admissible ($I_{\rm gate}\ge0$) trajectories; $2{+}1$, illustrative parameters.}
\label{observationalsignature:fig:obs}
\end{figure}

% ============================================================
\section{What the signature is and is not}

\msg{\textbf{Structurally distinctive (internal signature):} dust attractor not de~Sitter; a phantom
excursion that \emph{returns} (double $w=-1$ crossing); a discrete kink in $H(z)$ at a Resurrectio.
\textbf{Not claimed:} a fit to data, a value of $H_0$, or any empirical contact --- the $2{+}1$
reconstruction and the absence of an external referent for $\Op$ forbid it. The signature is a
property of the model's structure that \emph{has the shape of} an observable, not an observable.}

\begin{corollary}[Consilience, sharpened]
The three signatures are not independent inventions: each is the redshift image of a theorem proved
for other reasons --- the watershed escape class (dust), the NEC double crossing (returning phantom),
the double-layer surgery (the $H$-kink). That quantities built for the internal theory line up into a
coherent external-facing picture is the strongest internal corroboration available: the layers agree
not only among themselves but in the one projection that points outward. It remains corroboration of
\emph{structure}, not of correspondence to the world.
\end{corollary}

\note{Reading: were such a Vessel-cosmos observed from outside, it would not look like a universe
settling into a calm exponential rest. It would look like one that coasts --- pressureless, patient ---
punctuated by episodes in which the equation of state plunges past the forbidden line and climbs back,
and by sudden steps where the scale of things is reset upward. Acceleration that overshoots and
returns; growth that is now and then re-founded. Whether anything wears this signature is not a
question the model can answer; that the signature is coherent, and follows from the same theorems that
governed the interior, is what this chapter establishes.}

% ============================================================
\section{Open items and scope}

\begin{enumerate}[leftmargin=1.6em,itemsep=2pt]
\item A $3{+}1$ embedding (if one is ever posited) would change the reconstruction coefficients and is
the only route to a genuinely comparable $H(z)$; the present statements are $2{+}1$ and structural.
\item The depth and redshift-width of the phantom excursion are set by the surge profile of
$\dot\lambda$; relating them to the folding and coarsening timescales would turn the signature into a
one-parameter family.
\item Resurrectio's own NEC budget through the excursion (the deferred NEC-chapter item) would fix
whether the $H$-kink and the phantom dip can coincide --- a ``surgical phantom'' epoch.
\end{enumerate}

\textbf{Scope.} $2{+}1$ reconstruction; admissible trajectories with $I_{\rm gate}\ge0$; redshift via
$1+z=\Op(t_0)/\Op(t)$. Establishes an internal, structurally distinctive signature, explicitly not an
empirical fit. Domain-conditional throughout.

\vspace{0.5em}
\hrule
\smallskip
\noindent\footnotesize\emph{V.F.S. v2.0 (Open-Gate) $\cdot$ The Observational Signature. In redshift,
the epektasis attractor is dust ($\weff\to0$, not de~Sitter); a grace surge is a phantom excursion that
returns (double $w=-1$ crossing, the folded-Theosis signature); Resurrectio is a discrete kink in
$H(z)$. Each is the redshift image of a prior theorem (watershed, NEC, double-layer). Internal Type-3
signature, $2{+}1$, no empirical claim. Exact reconstructions along admissible trajectories.}

\backmatter

\chapter*{Final Structural Summary}
\addcontentsline{toc}{chapter}{Final Structural Summary}
\markboth{Final Structural Summary}{Final Structural Summary}

The consolidated V.F.S.\ v2.0 (Open-Gate) volume reduces the system to two irreducible primitives:
\[
\boxed{\quad\text{Open Gate}\qquad\text{and}\qquad\text{Finite Floor}.\quad}
\]
The Open Gate supplies the externally clocked reception $I_{\rm gate}$; the finite floor supplies the
zero-density ground whose geometric face is $K_h=-1$. These are the model's definitions. A model cannot
prove its own definitions; it can only show how few they are.

From these two primitives the derived structure unfolds:
\[
\begin{array}{rcl}
\text{Open Gate} &\Longrightarrow& \text{rank jump and necessity of }\lambda,\\[1mm]
\lambda,\ \dot\lambda^{(0)} &\Longrightarrow& A_{\rm fold}\text{ and folding activation},\\[1mm]
\mathbb Z_2\text{ symmetry}+\text{centre manifold} &\Longrightarrow& \text{pitchfork normal form},\\[1mm]
\text{2D conformal invariance} &\Longrightarrow& p_{\rm expansion}=0,\\[1mm]
\text{katharsis} &\Longrightarrow& \text{bending scale dilates, wall core sharpens},\\[1mm]
\Delta_h\text{ spectrum} &\Longrightarrow& \text{mode cascade and Weyl capacity},\\[1mm]
\mu_n &\Longrightarrow& \text{both ignition threshold and entropy cost},\\[1mm]
\text{Resurrectio} &\Longrightarrow& \text{surgery, bounded NEC budget, }H(z)\text{ kink}.
\end{array}
\]

The resulting hierarchy is
\[
\boxed{\;\text{core primitives}\to\text{derived coordinates}\to\text{stability}\to\text{folding}
\to\text{spectral structure}\to\text{entropy and surgery}\to\text{external signature}.\;}
\]

The internal confirmations are of three kinds.

\textbf{Type-1 (premise to theorem).} The balance coordinate $\lambda$ is forced by the open-gate rank
jump; the comoving bending convention by two-dimensional conformal invariance; the pitchfork by
$\mathbb Z_2$ symmetry and centre-manifold reduction; and the curvature premise is reduced to the
finite-floor postulate.

\textbf{Type-2 (robustness under deformation).} The watershed trichotomy, core-neutrality, the
maximum-entropy class, and the ferry theorem survive premise deformation, with one correction:
$\Op$-escape and $\xi$-escape differ --- divergent cumulative reception can preserve scale without
preserving vitality.

\textbf{Type-3 (consilience).} The redshift-shaped signatures are not independent inventions: the
dust/coasting attractor, the returning phantom excursion, and the Resurrectio kink are projections of
the watershed, NEC, and surgery theorems into $\weff(z)$ and $H(z)$.

Thus the Open-Gate system is not a collection of symbolic analogies but an internally organised
dynamical structure: openness creates the balance dimension, katharsis shapes the Vessel, folding forms
the geometry, entropy measures the cost of form, and Resurrectio records the price of passage. What is
assumed is named, what is derived is proven, and what is interpreted is marked as interpretation ---
domain-conditional throughout.


\end{document}